% Options for packages loaded elsewhere
% Options for packages loaded elsewhere
\PassOptionsToPackage{unicode}{hyperref}
\PassOptionsToPackage{hyphens}{url}
\PassOptionsToPackage{dvipsnames,svgnames,x11names}{xcolor}
%
\documentclass[
  english,
  11pt,
  a4paper,
]{article}
\usepackage{xcolor}
\usepackage[margin=25mm]{geometry}
\usepackage{amsmath,amssymb}
\setcounter{secnumdepth}{-\maxdimen} % remove section numbering
\usepackage{iftex}
\ifPDFTeX
  \usepackage[T1]{fontenc}
  \usepackage[utf8]{inputenc}
  \usepackage{textcomp} % provide euro and other symbols
\else % if luatex or xetex
  \usepackage{unicode-math} % this also loads fontspec
  \defaultfontfeatures{Scale=MatchLowercase}
  \defaultfontfeatures[\rmfamily]{Ligatures=TeX,Scale=1}
\fi
\usepackage{lmodern}
\ifPDFTeX\else
  % xetex/luatex font selection
\fi
% Use upquote if available, for straight quotes in verbatim environments
\IfFileExists{upquote.sty}{\usepackage{upquote}}{}
\IfFileExists{microtype.sty}{% use microtype if available
  \usepackage[]{microtype}
  \UseMicrotypeSet[protrusion]{basicmath} % disable protrusion for tt fonts
}{}
\makeatletter
\@ifundefined{KOMAClassName}{% if non-KOMA class
  \IfFileExists{parskip.sty}{%
    \usepackage{parskip}
  }{% else
    \setlength{\parindent}{0pt}
    \setlength{\parskip}{6pt plus 2pt minus 1pt}}
}{% if KOMA class
  \KOMAoptions{parskip=half}}
\makeatother
% Make \paragraph and \subparagraph free-standing
\makeatletter
\ifx\paragraph\undefined\else
  \let\oldparagraph\paragraph
  \renewcommand{\paragraph}{
    \@ifstar
      \xxxParagraphStar
      \xxxParagraphNoStar
  }
  \newcommand{\xxxParagraphStar}[1]{\oldparagraph*{#1}\mbox{}}
  \newcommand{\xxxParagraphNoStar}[1]{\oldparagraph{#1}\mbox{}}
\fi
\ifx\subparagraph\undefined\else
  \let\oldsubparagraph\subparagraph
  \renewcommand{\subparagraph}{
    \@ifstar
      \xxxSubParagraphStar
      \xxxSubParagraphNoStar
  }
  \newcommand{\xxxSubParagraphStar}[1]{\oldsubparagraph*{#1}\mbox{}}
  \newcommand{\xxxSubParagraphNoStar}[1]{\oldsubparagraph{#1}\mbox{}}
\fi
\makeatother

\usepackage{color}
\usepackage{fancyvrb}
\newcommand{\VerbBar}{|}
\newcommand{\VERB}{\Verb[commandchars=\\\{\}]}
\DefineVerbatimEnvironment{Highlighting}{Verbatim}{commandchars=\\\{\}}
% Add ',fontsize=\small' for more characters per line
\usepackage{framed}
\definecolor{shadecolor}{RGB}{241,243,245}
\newenvironment{Shaded}{\begin{snugshade}}{\end{snugshade}}
\newcommand{\AlertTok}[1]{\textcolor[rgb]{0.68,0.00,0.00}{#1}}
\newcommand{\AnnotationTok}[1]{\textcolor[rgb]{0.37,0.37,0.37}{#1}}
\newcommand{\AttributeTok}[1]{\textcolor[rgb]{0.40,0.45,0.13}{#1}}
\newcommand{\BaseNTok}[1]{\textcolor[rgb]{0.68,0.00,0.00}{#1}}
\newcommand{\BuiltInTok}[1]{\textcolor[rgb]{0.00,0.23,0.31}{#1}}
\newcommand{\CharTok}[1]{\textcolor[rgb]{0.13,0.47,0.30}{#1}}
\newcommand{\CommentTok}[1]{\textcolor[rgb]{0.37,0.37,0.37}{#1}}
\newcommand{\CommentVarTok}[1]{\textcolor[rgb]{0.37,0.37,0.37}{\textit{#1}}}
\newcommand{\ConstantTok}[1]{\textcolor[rgb]{0.56,0.35,0.01}{#1}}
\newcommand{\ControlFlowTok}[1]{\textcolor[rgb]{0.00,0.23,0.31}{\textbf{#1}}}
\newcommand{\DataTypeTok}[1]{\textcolor[rgb]{0.68,0.00,0.00}{#1}}
\newcommand{\DecValTok}[1]{\textcolor[rgb]{0.68,0.00,0.00}{#1}}
\newcommand{\DocumentationTok}[1]{\textcolor[rgb]{0.37,0.37,0.37}{\textit{#1}}}
\newcommand{\ErrorTok}[1]{\textcolor[rgb]{0.68,0.00,0.00}{#1}}
\newcommand{\ExtensionTok}[1]{\textcolor[rgb]{0.00,0.23,0.31}{#1}}
\newcommand{\FloatTok}[1]{\textcolor[rgb]{0.68,0.00,0.00}{#1}}
\newcommand{\FunctionTok}[1]{\textcolor[rgb]{0.28,0.35,0.67}{#1}}
\newcommand{\ImportTok}[1]{\textcolor[rgb]{0.00,0.46,0.62}{#1}}
\newcommand{\InformationTok}[1]{\textcolor[rgb]{0.37,0.37,0.37}{#1}}
\newcommand{\KeywordTok}[1]{\textcolor[rgb]{0.00,0.23,0.31}{\textbf{#1}}}
\newcommand{\NormalTok}[1]{\textcolor[rgb]{0.00,0.23,0.31}{#1}}
\newcommand{\OperatorTok}[1]{\textcolor[rgb]{0.37,0.37,0.37}{#1}}
\newcommand{\OtherTok}[1]{\textcolor[rgb]{0.00,0.23,0.31}{#1}}
\newcommand{\PreprocessorTok}[1]{\textcolor[rgb]{0.68,0.00,0.00}{#1}}
\newcommand{\RegionMarkerTok}[1]{\textcolor[rgb]{0.00,0.23,0.31}{#1}}
\newcommand{\SpecialCharTok}[1]{\textcolor[rgb]{0.37,0.37,0.37}{#1}}
\newcommand{\SpecialStringTok}[1]{\textcolor[rgb]{0.13,0.47,0.30}{#1}}
\newcommand{\StringTok}[1]{\textcolor[rgb]{0.13,0.47,0.30}{#1}}
\newcommand{\VariableTok}[1]{\textcolor[rgb]{0.07,0.07,0.07}{#1}}
\newcommand{\VerbatimStringTok}[1]{\textcolor[rgb]{0.13,0.47,0.30}{#1}}
\newcommand{\WarningTok}[1]{\textcolor[rgb]{0.37,0.37,0.37}{\textit{#1}}}

\usepackage{longtable,booktabs,array}
\newcounter{none} % for unnumbered tables
\usepackage{calc} % for calculating minipage widths
% Correct order of tables after \paragraph or \subparagraph
\usepackage{etoolbox}
\makeatletter
\patchcmd\longtable{\par}{\if@noskipsec\mbox{}\fi\par}{}{}
\makeatother
% Allow footnotes in longtable head/foot
\IfFileExists{footnotehyper.sty}{\usepackage{footnotehyper}}{\usepackage{footnote}}
\makesavenoteenv{longtable}
\usepackage{graphicx}
\makeatletter
\newsavebox\pandoc@box
\newcommand*\pandocbounded[1]{% scales image to fit in text height/width
  \sbox\pandoc@box{#1}%
  \Gscale@div\@tempa{\textheight}{\dimexpr\ht\pandoc@box+\dp\pandoc@box\relax}%
  \Gscale@div\@tempb{\linewidth}{\wd\pandoc@box}%
  \ifdim\@tempb\p@<\@tempa\p@\let\@tempa\@tempb\fi% select the smaller of both
  \ifdim\@tempa\p@<\p@\scalebox{\@tempa}{\usebox\pandoc@box}%
  \else\usebox{\pandoc@box}%
  \fi%
}
% Set default figure placement to htbp
\def\fps@figure{htbp}
\makeatother


% definitions for citeproc citations
\NewDocumentCommand\citeproctext{}{}
\NewDocumentCommand\citeproc{mm}{%
  \begingroup\def\citeproctext{#2}\cite{#1}\endgroup}
\makeatletter
 % allow citations to break across lines
 \let\@cite@ofmt\@firstofone
 % avoid brackets around text for \cite:
 \def\@biblabel#1{}
 \def\@cite#1#2{{#1\if@tempswa , #2\fi}}
\makeatother
\newlength{\cslhangindent}
\setlength{\cslhangindent}{1.5em}
\newlength{\csllabelwidth}
\setlength{\csllabelwidth}{3em}
\newenvironment{CSLReferences}[2] % #1 hanging-indent, #2 entry-spacing
 {\begin{list}{}{%
  \setlength{\itemindent}{0pt}
  \setlength{\leftmargin}{0pt}
  \setlength{\parsep}{0pt}
  % turn on hanging indent if param 1 is 1
  \ifodd #1
   \setlength{\leftmargin}{\cslhangindent}
   \setlength{\itemindent}{-1\cslhangindent}
  \fi
  % set entry spacing
  \setlength{\itemsep}{#2\baselineskip}}}
 {\end{list}}
\usepackage{calc}
\newcommand{\CSLBlock}[1]{\hfill\break\parbox[t]{\linewidth}{\strut\ignorespaces#1\strut}}
\newcommand{\CSLLeftMargin}[1]{\parbox[t]{\csllabelwidth}{\strut#1\strut}}
\newcommand{\CSLRightInline}[1]{\parbox[t]{\linewidth - \csllabelwidth}{\strut#1\strut}}
\newcommand{\CSLIndent}[1]{\hspace{\cslhangindent}#1}

\ifLuaTeX
\usepackage[bidi=basic,shorthands=off]{babel}
\else
\usepackage[bidi=default,shorthands=off]{babel}
\fi
\ifLuaTeX
  \usepackage{selnolig} % disable illegal ligatures
\fi


\setlength{\emergencystretch}{3em} % prevent overfull lines

\providecommand{\tightlist}{%
  \setlength{\itemsep}{0pt}\setlength{\parskip}{0pt}}


% Novo pacote para usar input{appendix_a.tex}
%\usepackage{ragged2e}


\usepackage{pdfpages}
\usepackage{float}
\usepackage{xurl}
\usepackage{caption}
\captionsetup[table]{labelformat=empty}
\captionsetup[figure]{labelformat=empty}
\makeatletter
\@ifpackageloaded{caption}{}{\usepackage{caption}}
\AtBeginDocument{%
\ifdefined\contentsname
  \renewcommand*\contentsname{Table of contents}
\else
  \newcommand\contentsname{Table of contents}
\fi
\ifdefined\listfigurename
  \renewcommand*\listfigurename{List of Figures}
\else
  \newcommand\listfigurename{List of Figures}
\fi
\ifdefined\listtablename
  \renewcommand*\listtablename{List of Tables}
\else
  \newcommand\listtablename{List of Tables}
\fi
\ifdefined\figurename
  \renewcommand*\figurename{Figure}
\else
  \newcommand\figurename{Figure}
\fi
\ifdefined\tablename
  \renewcommand*\tablename{Table}
\else
  \newcommand\tablename{Table}
\fi
}
\@ifpackageloaded{float}{}{\usepackage{float}}
\floatstyle{ruled}
\@ifundefined{c@chapter}{\newfloat{codelisting}{h}{lop}}{\newfloat{codelisting}{h}{lop}[chapter]}
\floatname{codelisting}{Listing}
\newcommand*\listoflistings{\listof{codelisting}{List of Listings}}
\makeatother
\makeatletter
\makeatother
\makeatletter
\@ifpackageloaded{caption}{}{\usepackage{caption}}
\@ifpackageloaded{subcaption}{}{\usepackage{subcaption}}
\makeatother
\usepackage{bookmark}
\IfFileExists{xurl.sty}{\usepackage{xurl}}{} % add URL line breaks if available
\urlstyle{same}
\hypersetup{
  pdftitle={Supplementary material},
  pdfauthor={Victor Rafael Rezende Celestino; Artur Guerra Rosa; Andréia Elizabeth Silva Barros; Gabriela Corsano; Luis Javier Zeballos; Rodolfo Dondo; Silvia Araujo dos Reis},
  pdflang={en},
  colorlinks=true,
  linkcolor={blue},
  filecolor={Maroon},
  citecolor={Blue},
  urlcolor={Blue},
  pdfcreator={LaTeX via pandoc}}


\title{Supplementary material}
\usepackage{etoolbox}
\makeatletter
\providecommand{\subtitle}[1]{% add subtitle to \maketitle
  \apptocmd{\@title}{\par {\large #1 \par}}{}{}
}
\makeatother
\subtitle{Computational Scalability of Deterministic and Two-Stage
Stochastic MILP Models for Grain-Storage Network Design and Hub
Location}
\author{Victor Rafael Rezende Celestino \and Artur Guerra
Rosa \and Andréia Elizabeth Silva Barros \and Gabriela Corsano \and Luis
Javier Zeballos \and Rodolfo Dondo \and Silvia Araujo dos Reis}
\date{2026-11-18}
\begin{document}
\maketitle

\renewcommand*\contentsname{Table of contents}
{
\setcounter{tocdepth}{1}
\tableofcontents
}

\newpage{}

These appendices accompany the article as separate documents. Appendix A is a systematic literature review of grain facility location models; Appendix B describes Human-supervised AI-assisted research; Appendix C presents deterministic and two-stage stochastic MILP implementation in Python and GAMS; Appendix D provides data and preparation of the experimental instances; Appendix E illustrates geographic decisions and origin-destination flows.

\newpage{}

\section{Appendix A Systematic Literature Review of Grain Facility-Location Models}\label{appendix-a-systematic-literature-review-of-grain-facility-location-models}

\includepdf[pages=-,pagecommand={\thispagestyle{empty}},fitpaper=true]{appendix_a_original.pdf}

%\input{appendix_a}

\newpage{}

\section{Appendix B Human-supervised AI-assisted
research}\label{appendix-b-human-supervised-ai-assisted-research}

\subsection{B1 Purpose and scope}\label{b1-purpose-and-scope}

AI assistants supported the development of this work through
mathematical discussion, implementation, technical criticism and
scientific writing. Human researchers defined the planning problem,
selected assumptions and experimental configurations, reviewed proposed
changes and interpreted the resulting investment and operating
decisions. This appendix explains how those activities were connected
through mathematical review and computational verification.

The procedure follows six phases: Human Scientific Intent, AI-Assisted
Implementation, Technical Verification, Scientific Validation, Human
Scientific Decision, and Provenance/Reporting. Section B3 explains these
phases through examples from the shared-investment stochastic model;
Section B4 distinguishes the four levels of verification used to assess
the work.

\subsection{B2 Tools and human
responsibilities}\label{b2-tools-and-human-responsibilities}

Gurobot was used for formulation and implementation discussions. Gemini
and GPT systems supported technical criticism and alternative
interpretations. ChatGPT Classic supported coordination and writing,
while Codex supported Python implementation, tests, diagnostics and
documentation. Researchers reviewed the proposed changes and decided
which ones to adopt.

The numerical tools had distinct roles. Python organized the data and
calculations. The native Gurobi implementation and the common Pyomo
formulation provided the mathematical-programming pipelines evaluated
here. Pyomo's APPSI interfaces connected that formulation to Gurobi and
HiGHS. GitHub retained versioned source and experiment records, and
Slurm scheduled the HPC calculations. Researchers executed the
experimental commands and assessed their outputs.

\subsection{B3 From the scientific question to a computational
result}\label{b3-from-the-scientific-question-to-a-computational-result}

The six phases connect a planning question to an implemented and
interpreted experiment. Table B.1 illustrates their application to the
two-stage model. A discrepancy found during verification returns the
work to the relevant assumption or implementation.

\begin{longtable}[]{@{}
  >{\raggedright\arraybackslash}p{(\linewidth - 4\tabcolsep) * \real{0.3333}}
  >{\raggedright\arraybackslash}p{(\linewidth - 4\tabcolsep) * \real{0.3333}}
  >{\raggedright\arraybackslash}p{(\linewidth - 4\tabcolsep) * \real{0.3333}}@{}}
\caption{Table B.1. Research phases and their
application.}\tabularnewline
\toprule\noalign{}
\begin{minipage}[b]{\linewidth}\raggedright
Phase
\end{minipage} & \begin{minipage}[b]{\linewidth}\raggedright
Research activity
\end{minipage} & \begin{minipage}[b]{\linewidth}\raggedright
Example from this work
\end{minipage} \\
\midrule\noalign{}
\endfirsthead
\toprule\noalign{}
\begin{minipage}[b]{\linewidth}\raggedright
Phase
\end{minipage} & \begin{minipage}[b]{\linewidth}\raggedright
Research activity
\end{minipage} & \begin{minipage}[b]{\linewidth}\raggedright
Example from this work
\end{minipage} \\
\midrule\noalign{}
\endhead
\bottomrule\noalign{}
\endlastfoot
P1 Human Scientific Intent & Define the planning question, decisions and
assumptions & Locate and size storage facilities before uncertain supply
and demand are realized \\
P2 AI-Assisted Implementation & Translate reviewed equations into code
and propose checks & Represent common investment decisions and
scenario-dependent grain flows \\
P3 Technical Verification & Check executable behavior and mathematical
properties & Verify that investment is charged once and operating costs
are probability weighted \\
P4 Scientific Validation & Examine whether experiments address the
question and how their outputs should be interpreted & Assess
feasibility, objective bounds, resources and agricultural service
jointly \\
P5 Human Scientific Decision & Accept, revise or reject interpretations
and methodological changes & Interpret penalized plans using unmet
demand and storage/handling indicators \\
P6 Provenance/Reporting & Associate reported results with their
implemented experiment & Retain the source, input, scenario, solver and
validation records for each result \\
\end{longtable}

An ilustrative example documents a substantive redirection concerning the research objective. Early
experiments used a hierarchical formulation in which capacity-related and economic criteria had different priorities. Reviewing the interpretation of those outputs led the researchers to enforce the single objective reported here: investment is counted once, operating costs are probability weighted, and the three artificial penalties remain explicit. The earlier experiments were retained as a separate historical cohort rather than relabeled as scalar results. This change concerned the scientific question expressed by the model, not merely whether the program executed successfully. The
\href{https://github.com/unb-lamfo-or-ai-hpc/model-agrologistic/blob/2dbd780b1dfdc77d8e31f164e5d3bcc3b2708410/docs/mvp2_completion_handoff_20261022.md}{retained
formulation history} links the archived hierarchy and scalar contract.
The episode illustrates human responsibility for accepting model interpretation.

\subsection{B4 Four levels of
verification}\label{b4-four-levels-of-verification}

The workflow in Figure B.1 links the research phases to repeated review
and verification. A proposed change is examined at the levels
appropriate to its effect on the scientific calculation.

\begin{figure}[H]

{\centering \pandocbounded{\includegraphics[keepaspectratio]{human_ai_workflow.png}}

}

\caption{Figure B.1. Human-supervised research workflow.}

\end{figure}%

The return arrows represent revision after a discrepancy is identified.
For example, a successful program is followed by checks of investment
accounting and probability weights, and then by examination of
original-unit balances and the agricultural meaning of the returned
plan. Table B.2 distinguishes these complementary levels.

\begin{longtable}[]{@{}
  >{\raggedright\arraybackslash}p{(\linewidth - 4\tabcolsep) * \real{0.3333}}
  >{\raggedright\arraybackslash}p{(\linewidth - 4\tabcolsep) * \real{0.3333}}
  >{\raggedright\arraybackslash}p{(\linewidth - 4\tabcolsep) * \real{0.3333}}@{}}
\caption{Table B.2. Complementary levels of
verification.}\tabularnewline
\toprule\noalign{}
\begin{minipage}[b]{\linewidth}\raggedright
Level
\end{minipage} & \begin{minipage}[b]{\linewidth}\raggedright
Question examined
\end{minipage} & \begin{minipage}[b]{\linewidth}\raggedright
Evidence used in this work
\end{minipage} \\
\midrule\noalign{}
\endfirsthead
\toprule\noalign{}
\begin{minipage}[b]{\linewidth}\raggedright
Level
\end{minipage} & \begin{minipage}[b]{\linewidth}\raggedright
Question examined
\end{minipage} & \begin{minipage}[b]{\linewidth}\raggedright
Evidence used in this work
\end{minipage} \\
\midrule\noalign{}
\endhead
\bottomrule\noalign{}
\endlastfoot
V1 Software functionality & Does the intended operation execute and
return the required products? & Data-loading, model-construction and
output tests \\
V2 Mathematical consistency & Does the implementation express the
reviewed equations and units? & Analytical examples, deterministic
specialization, investment accounting and objective reconstruction \\
V3 Protocol conformity & Was the specified instance and experimental
configuration used? & Source/input identities, scenario membership,
solver settings and resource records \\
V4 Empirical scientific support & What conclusions do the observed
decisions and resources support? & Accepted original residual checks,
objective bounds, service/capacity indicators and qualified
interpretation \\
\end{longtable}

Small examples examined investment charged once, probability weighting,
capacity units and activation conditions. Large-instance outputs were
checked against material balances, eligibility and capacity
relationships in their original units, with all fourteen objective
components reconstructed separately. These checks support different
conclusions: the relative gap quantifies uncertainty about the scalar
objective, while service and capacity indicators explain the physical
consequences of the plan.

\subsection{B5 Traceability of a computational
result}\label{b5-traceability-of-a-computational-result}

For each experiment, the retained record identifies the source version, input workbook, facility population, scenario design, solver settings, resource limits, decisions and validation outcome. This association
permits another researcher to connect a numerical table to the calculation that produced it. Partial and failed attempts are retained alongside successful ones to avoid selective reporting.
An integrity replay checks consistency of retained records and their hashes. A fresh numerical reproduction must generate new decisions and apply the independent validation procedure to them.

\subsection{B6 Limitations}\label{b6-limitations}

It was not possible to make available the complete conversational record. Different assistants may repeat
a common error, so agreement between them is not independent scientific review. Tests verify specified properties and examples, but they do not establish general correctness or substitute for substantive modeling judgment. Restricted inputs can also limit independent reconstruction of a large experiment.

The literature reconciliation separately checks bibliographic identity, ordered family names where registry metadata were available, and claims
supported by accessible primary passages. Abstract- or preview-level access does not establish that an unreported feature is absent. The present
article's comparison is bounded by the source-access record and its acknowledged mathematical provenance.
The study has no controlled comparison for measuring time saved, productivity or error reductions caused by AI assistance. Its potential methodological contribution is the inspectable connection between
proposals, human decisions, implementations and numerical evidence. The authors retain responsibility for all assumptions, reported results and
conclusions.

\newpage{}

\section{Appendix C Deterministic and Two-Stage Stochastic MILP
Implementations in Python and
GAMS}\label{appendix-c-deterministic-and-two-stage-stochastic-milp-implementations-in-python-and-gams}

\subsection{C1 From the planning problem to an executable
model}\label{c1-from-the-planning-problem-to-an-executable-model}

The implementation follows the grain-storage formulation attributed to Rosa (2026) and presented in the main article. Facility opening and
capacity investment are chosen before supply and demand are realized.
Grain movements, inventory and the permitted shortages or capacity exceedances are then determined for each operating trajectory. The purpose of this appendix is to explain how those decisions are represented in software and how the resulting calculations are checked.
A reader can follow two successive questions: how the equations become variables and constraints (Sections C2--C4), and what makes a returned answer acceptable (Section C5). Section C7 provides one implementation index and a practical route through public examples and the separate input-export procedure. The question related to which input arrays define the case is deffered to Appendix D.

The deterministic (DET) and stochastic (STO) models use the same economic objective and physical relationships. DET evaluates the original deterministic supply and demand arrays. STO represents nine complete supply-demand trajectories with a common investment vector and
scenario-dependent operation. The model considers 60 operating periods, but the information structure has two
decision stages.

Table C.1 relates the transport notation to the identifiers retained in the implementation. In the legacy Python labels, D denotes the
Portuguese \emph{depósito} (warehouse). This convention must be distinguished from topology D, which permits direct origin-to-destination movements.

\begin{longtable}[]{@{}
  >{\raggedright\arraybackslash}p{(\linewidth - 4\tabcolsep) * \real{0.3333}}
  >{\raggedright\arraybackslash}p{(\linewidth - 4\tabcolsep) * \real{0.3333}}
  >{\raggedright\arraybackslash}p{(\linewidth - 4\tabcolsep) * \real{0.3333}}@{}}
\caption{Table C.1. Correspondence between physical transport families
and implementation labels.}\tabularnewline
\toprule\noalign{}
\begin{minipage}[b]{\linewidth}\raggedright
Physical movement
\end{minipage} & \begin{minipage}[b]{\linewidth}\raggedright
Notation used in this appendix
\end{minipage} & \begin{minipage}[b]{\linewidth}\raggedright
Python label / GAMS variable
\end{minipage} \\
\midrule\noalign{}
\endfirsthead
\toprule\noalign{}
\begin{minipage}[b]{\linewidth}\raggedright
Physical movement
\end{minipage} & \begin{minipage}[b]{\linewidth}\raggedright
Notation used in this appendix
\end{minipage} & \begin{minipage}[b]{\linewidth}\raggedright
Python label / GAMS variable
\end{minipage} \\
\midrule\noalign{}
\endhead
\bottomrule\noalign{}
\endlastfoot
Origin to warehouse & OW & OD / xOW \\
Warehouse to destination & WD & DC / xWD \\
Warehouse to warehouse & WW & DD / xWW \\
Direct origin to destination & Direct OD & OC / xOD \\
\end{longtable}

Topology W excludes direct OD movements; topology D permits them alongside warehouse-mediated paths. Permission does not require positive
direct flow. The native Gurobi cohort comprises nine W observations
across h215--h500 and two additional D observations at h215. The latter
form a separate topology comparison, rather than additional points in
the warehouse-only scaling sequence.

\subsection{C2 Shared investment and scenario-dependent
operation}\label{c2-shared-investment-and-scenario-dependent-operation}

The common implementation first prepares effective inputs and selects
the admissible product-specific network. It then creates opening,
expansion and bulk-handling decisions only for eligible facilities.
Investment quantities and their activation binaries have no scenario
index. Operational flows, stock and slacks are defined over the selected
arcs, products and periods, with a scenario index in STO. An absent arc
has no decision variable.

This organization makes the distinction between stages explicit. Each
scenario uses the same installed capacity, while its operating decisions
respond to its own supply and demand. The objective charges investment
once and takes the probability-weighted sum of operating and penalty
costs. Scenario probabilities do not multiply material balances or
capacity limits. A scenario's grain must still satisfy the physical
equations in its original units.

The following expression makes the shared-investment accounting
explicit. Let \(W^N\), \(W^X\) and \(W^B\) denote candidate,
expansion-eligible and bulk-handling-eligible facilities; \(O\), \(W\)
and \(C^D\) denote origins, warehouses and domestic destinations.
Products and periods are indexed by \(P\) and \(T\), and \(S\) is the
scenario set. The first-stage vector is \(X=(y,z,g,k,e,b)\): the first
three entries are activation binaries and the remaining entries are the
corresponding investment quantities. For each scenario, \(Y^s\) contains
arc flows \(x\), end stocks \(I\), domestic shortages \(u\) and
static/reception exceedances \(q^K,q^R\).

The scalar objective and its fourteen components are

\[
\begin{aligned}
\min_{X,\{Y^s\}}\ Z={}&C^{INV}(X)
+\sum_{s\in S}\pi_s
 \left[C^{TR}(Y^s)+C^{HOLD}(Y^s)+C^{PEN}(Y^s)\right],
\\
C^{INV}(X)={}&\sum_{w\in W^N}(F^N_wy_w+c^N_wk_w)
 +\sum_{w\in W^X}(F^X_wz_w+c^X_we_w)
 +\sum_{w\in W^B}(F^B_wg_w+c^B_wb_w),\\
C^{TR}(Y)={}&\sum_{t\in T}\sum_{(o,w,p)\in A^{OW}}
 (f^O_o\ell^{OW}_{ow}+\tau_w)x^{OW}_{owpt}\\
&+\sum_{t\in T}\sum_{(w,c,p)\in A^{WD}}f^W_w\ell^{WD}_{wc}x^{WD}_{wcpt}\\
&+\sum_{t\in T}\sum_{(w,v,p)\in A^{WW}}
 (\alpha f^W_w\ell^{WW}_{wv}+\tau_v)x^{WW}_{wvpt}\\
&+\sum_{t\in T}\sum_{(o,c,p)\in A^{OD}}f^O_o\ell^{OD}_{oc}x^{OD}_{ocpt},\\
C^{HOLD}(Y)={}&\sum_{w\in W,p\in P,t\in T}h_{wp}I_{wpt},\\
C^{PEN}(Y)={}&\sum_{c\in C^D,p\in P,t\in T}\lambda^D_{cp}u_{cpt}
 +\sum_{w\in W,t\in T}(\lambda^K_wq^K_{wt}+\lambda^R_wq^R_{wt}).
\end{aligned}
\tag{C.1}
\]

In the component expressions, \(Y\) denotes one scenario's operating
vector. Fixed investment charges \(F\) are in BRL and quantity-dependent
charges \(c\) are in BRL per investment-tonne. Each transport
coefficient combines a distance \(\ell\) in km and a freight rate \(f\)
in BRL/(t km), with a receiving charge \(\tau\) in BRL/t where
applicable; \(\alpha\) multiplies interwarehouse freight. A holding
tariff is quoted in BRL/(t period); the coefficient \(h\) charges one
modeled period of stock, so \(hI\) is that period's BRL cost. The three
\(\lambda\) coefficients price their respective shortage,
stock-exceedance and reception-overflow quantities at each period. The
scenario weights satisfy \(\sum_s\pi_s=1\). Stock and period throughput
consequently enter different physical checks even when both are measured
in tonnes.

For DET, the original arrays are embedded as one scenario with
probability one. The resulting solution is projected back to the
deterministic output format. Thus, DET is neither the solution for a
conveniently selected STO realization nor a solution for
scenario-averaged data. This specialization is also useful for
verification: a one-scenario STO case and its DET counterpart should
agree on objective components and feasibility, allowing for alternative
optimal decision vectors.

The physical equations are implemented in a corresponding sequence.
Supply leaves each origin; warehouse inventory links successive periods;
domestic arrivals plus shortage satisfy demand; and exports remain
optional within their specified limits. Static capacity bounds stock,
whereas reception and shipping capacities bound period throughput after
converting daily rates using operating days. Static and reception
exceedances have separate slacks. Shipping, conservation and route
eligibility remain hard constraints. For a closed candidate, the
emergency slacks must also be zero.

\subsection{C3 Native gurobipy
implementation}\label{c3-native-gurobipy-implementation}

This implementation builds and solves the extensive form directly
through the Gurobi Python API. It supplies the historical W STO scale
series and the h215 W/D comparisons. The native runs used a
14,400-second solver limit, a 1\% relative-gap target and seed 42. The
h215 W/D comparisons used four threads. The h300 W STO block contains
one observation at each of 1, 2, 4, 8 and 16 threads, and the h400 and
h500 W STO observations used one thread. The original h300 runs shared
one reservation, with fixed execution order and uncontrolled cache state
and shared-system load. This block describes the tested configurations.

The h215 W/D comparison shares the original input workbook, the
warehouse population, costs, scenario data and common route families. D
adds 592 direct origin-to-customer arcs, equivalent to 35,520 DET or
319,680 STO flow variables. The validated D incumbents are below the
corresponding W lower bounds. The retained bound intervals bracket the
reduction in the optimal penalized objective at 35.819--36.175\% for DET
and 43.313--44.083\% for STO. The intervals are mathematical bounds on
the penalized objective. Expected shortages increase from 55.16 to 68.48
million tonnes in DET and from 52.40 to 68.74 million tonnes in STO.
Artificial penalties account for approximately 95.39\% and 95.38\% of
the corresponding D totals. The original h215 runs record native times
of 103.93 and 2,582.59 seconds for D and 122.51 and 2,445.52 seconds for
W in DET and STO, respectively. The separate runs have different
attained gaps and describe their own effort. The W records provide the
historical control for the topology comparison.

Gurobi uses a common extensive-form builder for DET and STO. Its cost
construction separates six investment components, four transport
components, inventory holding and three artificial penalties: fourteen
components in total. The scalar model uses one optimization call under
the stated time allowance.

The native builder expresses closed-candidate implications using
indicators. An active candidate's redundant slack has no arbitrary
imposed upper bound. Section C4 explains the bounded representative used
for the common Pyomo formulation and why the representation must be
declared when interpreting native and Pyomo results.

Table C.2 links the main formulation to the principal software
components. The software reading index in Section C7 identifies the
implementations. Original experiments retain their own executed
versions.

\begin{longtable}[]{@{}
  >{\raggedright\arraybackslash}p{(\linewidth - 4\tabcolsep) * \real{0.3333}}
  >{\raggedright\arraybackslash}p{(\linewidth - 4\tabcolsep) * \real{0.3333}}
  >{\raggedright\arraybackslash}p{(\linewidth - 4\tabcolsep) * \real{0.3333}}@{}}
\caption{Table C.2. Mathematical relationships and their implementation
roles.}\tabularnewline
\toprule\noalign{}
\begin{minipage}[b]{\linewidth}\raggedright
Main equations
\end{minipage} & \begin{minipage}[b]{\linewidth}\raggedright
Computational role
\end{minipage} & \begin{minipage}[b]{\linewidth}\raggedright
Principal implementation
\end{minipage} \\
\midrule\noalign{}
\endfirsthead
\toprule\noalign{}
\begin{minipage}[b]{\linewidth}\raggedright
Main equations
\end{minipage} & \begin{minipage}[b]{\linewidth}\raggedright
Computational role
\end{minipage} & \begin{minipage}[b]{\linewidth}\raggedright
Principal implementation
\end{minipage} \\
\midrule\noalign{}
\endhead
\bottomrule\noalign{}
\endlastfoot
D1/S1 & Investment, operation and penalties & Investment/scenario cost
builders and the scalar objective \\
D2--D3 & Eligible investment and effective capacities & First-stage
constraints and static/reception/shipping expressions \\
D4--D6 / S2--S4 & Supply and inventory conservation & Scenario
physical-constraint construction \\
D7--D8 / S5--S6 & Domestic service and export limits & Destination
balance and upper-bound rows \\
D9--D10 / S7--S8 & Capacity and closed-candidate conditions & Capacity
rows and native indicators \\
D11/S9 & Domains and shared decisions & Eligible first-stage keys and
nonnegative recourse variables \\
\end{longtable}

The extensive-form builders implement these relationships; the
mathematical input contract prepares the effective arrays, and the
independent validator reconstructs balances and costs from exported
decisions. The worked example below illustrates this separation without
requiring a source-code lookup.

\subsection{C4 Pyomo and the bounded representation of candidate
slacks}\label{c4-pyomo-and-the-bounded-representation-of-candidate-slacks}

The Pyomo implementation and its bound derivation provide a common
scalar formulation for Gurobi and HiGHS. It uses the same effective
data, investment expressions and physical balances. The principal
representational difference concerns candidate emergency slacks. The
native model allows an open candidate's slack to be unbounded above. A
finite bound must therefore be justified mathematically, rather than
selected because it appears numerically large.

For scenario \(s\) and period \(t\), let \(G_{st}\) denote total
exogenous supply, and let \(I^0\) denote total initial inventory. Let
\(\overline H_{wt}\) be a finite upper bound on warehouse \(w\)'s
\emph{period} shipping capacity, obtained from the original daily rates,
operating days and maximum eligible investments. Ignoring mutual
exclusion when calculating this upper bound can only enlarge it. Define

\[
M^K_{st}=I^0+\sum_{\tau\leq t}G_{s\tau},
\qquad
M^R_{st}=G_{st}+\sum_{w\in W}\overline H_{wt}.
\tag{C.2}
\]

The first bound follows by summing warehouse balances. Interwarehouse
flows cancel, while deliveries to customers and other withdrawals are
nonnegative. Consequently, the stock at any one warehouse cannot exceed
initial stock plus cumulative supply. Since installed static capacity is
nonnegative, the minimum required static exceedance is also bounded by
\(M^K_{st}\).

The reception bound must account for circulation. Arrivals from origins
cannot exceed current exogenous supply. Arrivals from other warehouses
are bounded by their hard shipping capacities, even if grain traverses
several warehouses. Their sum therefore provides a valid upper bound on
incoming interwarehouse flow. Current supply alone would not provide
such a bound.

For each candidate warehouse with opening variable \(y_w\), Pyomo
imposes

\[
0\leq q^K_{swt}\leq M^K_{st}y_w,
\qquad
0\leq q^R_{swt}\leq M^R_{st}y_w.
\tag{C.3}
\]

When the candidate is closed, these constraints reproduce the native
zero-slack implication. When it is open, take any feasible native
physical and investment decisions and reduce its slack to the minimum
nonnegative excess required by the corresponding capacity inequality.
The bounds above contain that representative. Conversely, any solution
satisfying the bounded constraints also satisfies the original native
implication and the unchanged physical rows.

With nonnegative candidate slack coefficients and scenario
probabilities, reducing redundant slack cannot increase the objective.
The two representations therefore have the same physical/investment
projection and optimal scalar value. They do \textbf{not} have identical
full feasible sets in the redundant slack coordinates. Nor does the
proof establish identical relaxations, numerical conditioning or search
behavior.

This equivalence is established under the explicitly selected bounded\_scalar\_slacks\_v1 policy, implemented in the repository's \href{https://github.com/unb-lamfo-or-ai-hpc/model-agrologistic/blob/main/src/logic/pyomo_bounds.py}{Pyomo bounds module}. The formulation preserves the physical and investment projection and the optimal scalar objective value under the stated assumptions, without requiring identical feasible regions in redundant slack coordinates.
Intuitively, the transformation removes unnecessarily large slack values while retaining a representative of every feasible physical plan and investment vector.
It does not guarantee that two solver interfaces will explore the same search tree.

The implementation checks the proof's finite, nonnegative data conditions and records the derived bounds' identity and scope. Gurobi and HiGHS use Pyomo's APPSI interfaces. Each receives the reviewed linear expressions, variable domains and bounds from this common formulation.

To compare inputs to these adapters, the ordered algebra record includes
variable domains and bounds, objective coefficients and constant, and
constraint coefficients and bounds. Matching records describe the common
Pyomo algebra before solver transformations; they do not assert
identical presolved models. Construction, transfer, optimization and
validation are therefore measured separately. Public analytical and
intermediate-size controls test correctness and operational behavior,
but do not predict h215 runtime or memory. Large-instance interpretation
must retain these formulation, interface and resource scopes.

\subsubsection{C4.1 Parallel search and
repeatability}\label{c4.1-parallel-search-and-repeatability}

Gurobi's
\href{https://docs.gurobi.com/projects/optimizer/en/current/reference/parameters.html\#concurrentmip}{Threads
and ConcurrentMIP settings} control different aspects of execution. A
setting of \texttt{ConcurrentMIP=1} uses one MIP search, whose parallel
algorithms can use the permitted threads. Higher concurrent-MIP settings
divide the thread budget among independent MIP searches. The
common-Pyomo campaign used one search.

The solver seed initializes its pseudo-random choices. Branch selection,
node order, cuts and heuristics also respond to the model and the
evolving search state. Thread count can change the trajectory at the
same seed. Gurobi describes repeatability for fixed model and settings
on the same system, subject to time-dependent stopping and algorithm
choices
(\href{https://support.gurobi.com/hc/en-us/articles/360031636051-Is-Gurobi-deterministic}{technical
note}). The time-limited repetitions at seed 42 characterize observed
execution variation for the declared settings.

\paragraph{C4.2 Parallel tree search in HiGHS
1.15.1}\label{parallel-tree-search-in-highs-1.15.1}

A MILP solution involves several kinds of work. Presolve simplifies the
model; an LP relaxation supplies a bound; branch-and-bound explores
remaining integer decisions. Parallelism in one part does not imply
parallel execution of the whole calculation.

The HiGHS 1.15 series provides an initial parallel MIP implementation,
included in the 1.15.1 release used here. Setting \texttt{parallel="on"}
allows several workers to explore the tree, including cuts and
heuristics away from the root. Symmetry detection, clique-table queries
and analytic-centre calculations also offer parallel work. However,
parallel simplex is disabled for LP relaxations inside MIP in this
release. The parallel capability of a standalone LP solver must
therefore not be assumed for a MIP's root relaxation. These distinctions
follow the
\href{https://github.com/ERGO-Code/HiGHS/blob/v1.15.1/docs/src/parallel.md}{HiGHS
1.15.1 documentation} and
\href{https://github.com/ERGO-Code/HiGHS/releases/tag/v1.15.1}{release
description}.

The common-Pyomo campaign profile sets \texttt{threads} to 1, 2, 4, 8 or
16, \texttt{parallel="on"} and \texttt{random\_seed=42}, with an
eight-hour native limit and a 1\% relative-gap target. Each calculation
starts in a fresh process; requested options and values returned by the
native solver are retained. A thread ceiling is not a measurement of
active workers. Additional threads may provide little benefit when
root-relaxation work dominates or memory bandwidth limits progress.
Actual logs, CPU observations and elapsed times must support any
parallel-performance claim.

The completed h215 HiGHS observation belongs to a different execution
profile: its log records four threads, one maximum search worker and
parallel search off. Its accepted objective and bound remain valid, but
its runtime cannot establish the performance of the prospective parallel
profile. Fixing seed 42 specifies solver-specific pseudorandom choices;
it neither constructs a shared initial incumbent for Gurobi and HiGHS
nor guarantees identical parallel trajectories. Repeated runs under that
fixed seed quantify observed execution variation and do not represent
independent random problem instances.

\subsubsection{C4.3 A calculation from workbook to validated
result}\label{c4.2-a-calculation-from-workbook-to-validated-result}

Both Pyomo pipelines follow the same scientific sequence.

\begin{enumerate}
\def\labelenumi{\arabic{enumi}.}
\tightlist
\item
  Read the identified workbook, warehouse order and scenario definition.
  Convert the source quantities to the effective arrays described in
  Appendix D and select the product-specific admissible directed arcs.
\item
  Build the shared investment block and scenario-dependent balances,
  capacities and costs. Derive the bounded candidate-slack
  representation.
\item
  Transfer the Pyomo expressions to the native engine through APPSI:
  gurobipy for Gurobi and highspy for HiGHS. Record requested options
  and, where retained, native values and the ordered common-algebra
  fingerprint.
\item
  Optimize in a fresh process. Preserve the incumbent, bound, native
  status and solver runtime even when the target is not attained.
\item
  Export decisions, reconstruct costs and check original-unit residuals.
  Interpret financial cost, artificial penalties, service and capacity
  exceedances separately.
\end{enumerate}

Pyomo supplies algebra and options; Gurobi still selects its own
presolve, relaxation methods, cuts and heuristics. In the completed
Pyomo observation, Method, NodeMethod, Presolve and Cuts were -1
(automatic), Heuristics was 0.05 and ConcurrentMIP was 1. TimeLimit,
MIPGap, Threads, Seed, FeasibilityTol and IntFeasTol were explicitly
controlled. This is an automatic-strategy profile with declared
overrides, rather than an untouched default configuration. ConcurrentMIP
counts independent MIP searches; its value of one does not restrict the
engine to one worker thread. The
\href{https://docs.gurobi.com/projects/optimizer/en/current/reference/parameters.html}{native
parameter reference} defines these controls.

For HiGHS, the thread ceiling and parallel-search option must both be
reported, as explained in C4.1. Parallelism within a solve differs from
running two independent Slurm jobs at once. The common-Pyomo campaign
protocol allocates one node per calculation; its worker affinity, solver
thread ceiling and allocation CPU count are recorded separately. Neither
use of Pyomo nor seed 42 makes parallel execution or returned investment
vectors identical.

Data preparation, common-model construction, matrix identification,
solver transfer, native optimization, extraction and independent
validation have different timing boundaries. Some recorded regions
enclose other regions; their values must not be added to reconstruct an
end-to-end time. The microsecond data-read field at an already-loaded
model boundary is not the time required to read and prepare a large
workbook.

\subsubsection{C4.4 Finite campaign and progressive population
qualification}\label{c4.3-finite-campaign-and-progressive-population-qualification}

The campaign design crosses two solvers, three warehouse populations
(h300, h400 and h500), five thread ceilings (1, 2, 4, 8 and 16) and
three temporal repetitions: 90 planned runs across 30
solver-population-thread settings. Every run retains seed 42, a
28,800-second native limit and a 1\% target for the scalar objective.
The first h300 observations remain a separate matched cohort from the
historical native Gurobi results; a new campaign cell does not
retroactively replicate or replace a historical observation.

The
\href{https://github.com/unb-lamfo-or-ai-hpc/model-agrologistic/blob/82c2c19ce619f1251043829b10a4702ac88a7cdd/experiments/pyomo_scaling_serial_schedule_v1.json}{finite
schedule} defines one ordered sequence per solver. Parallel independent
jobs and parallel work inside a solver are different levels of
concurrency. Each scientific job requests an exclusive node and 192 GiB
for a 12-hour workflow envelope. Its worker affinity contains 16 logical
CPU IDs; the solver thread ceiling varies by configuration and the
scheduler's allocation count is recorded separately. A thread ceiling of
one therefore does not mean that only one logical CPU was allocated to
the complete workflow.

At each population, the one- and sixteen-thread pilots for both solvers
precede the remaining configurations. Advancement checks the recorded
inputs, ordered algebra, derived bounds, hardware class and operational
closure. The sampled worker-tree peak must retain 20\% headroom relative
to the requested memory. Live output-growth and personal-quota checks
separately protect the evidence reserve. Public analytical qualification
checks software and dependency behavior; it does not predict the memory
requirements of an unexecuted population.

A regular time limit without an incumbent remains a censored outcome,
with no incumbent objective, service plan or incumbent-based gap.
Interrupted or inconsistent execution closes progression and preserves
the evidence. Neither kind of outcome is silently retried or replaced by
a different solver. The final consolidation retained all 90 run states.
The campaign returned 34 observations: 20 validated at the 1\% target
and 14 without a usable incumbent. The scheduler cancelled the other 56
solve jobs before they ran. The h400 pilot control was the first failed
control. Section C8 reconciles those states and documents the memory
criterion.

\subsection{C5 Checking a returned
solution}\label{c5-checking-a-returned-solution}

Independent validation reconstructs the fourteen cost components and
evaluates the original physical and investment constraints from the
exported decisions. It uses absolute tolerance \(10^{-5}\) and relative
tolerance \(10^{-8}\), together with the explicit sparse-export rounding
allowance. Local constraints are examined individually. For a check with
left and right sides \(L\) and \(R\), the tolerance is
\(10^{-5}+10^{-8}\max(|L|,|R|)\) plus the applicable export-rounding
allowance. The residual is \(|L-R|\) for an equality and \(\max(0,L-R)\)
for an upper-bound inequality.

Three outcomes must remain distinct. A usable incumbent supplies
decisions that can be checked. An accepted incumbent also passes the
independent residual and cost checks. A quality-qualified result
additionally meets the stated gap target using a valid bound. A solver
or scheduler status alone does not establish these properties. When no
usable incumbent is available, financial costs, service levels and an
incumbent-based gap are unavailable.

The common gap uses the incumbent objective in its denominator, whereas
a solver may use a different native definition. Both are retained with
their definitions. A small gap for the penalized total remains a
statement about that total; it is not a relative-accuracy guarantee for
the financial subtotal or a proof that nominal capacity and full service
can be achieved without slack.

\subsubsection{C5.1 An illustrative balance and objective
check}\label{an-illustrative-balance-and-objective-check}

Consider a teaching example with one product, one operating day, one
origin, one candidate warehouse and one domestic destination. Two
scenarios have probability \(1/2\). A specified plan opens the warehouse
and installs 5 t of static capacity at BRL4/t. For this teaching example
only, a handling coefficient of 3 day\(^{-1}\) gives reception and
shipping rates of 15 t/day. The experimental coefficient is specified
separately in Appendix D. Initial stock and fixed opening cost are zero.
Table C.3 gives the plan's flows and the corresponding end stocks.

\begin{longtable}[]{@{}
  >{\raggedright\arraybackslash}p{(\linewidth - 4\tabcolsep) * \real{0.2727}}
  >{\raggedleft\arraybackslash}p{(\linewidth - 4\tabcolsep) * \real{0.3636}}
  >{\raggedleft\arraybackslash}p{(\linewidth - 4\tabcolsep) * \real{0.3636}}@{}}
\caption{Table C.3. Illustrative scenario balances.}\tabularnewline
\toprule\noalign{}
\begin{minipage}[b]{\linewidth}\raggedright
Quantity (t within the period, except end stock in t)
\end{minipage} & \begin{minipage}[b]{\linewidth}\raggedleft
Scenario A
\end{minipage} & \begin{minipage}[b]{\linewidth}\raggedleft
Scenario B
\end{minipage} \\
\midrule\noalign{}
\endfirsthead
\toprule\noalign{}
\begin{minipage}[b]{\linewidth}\raggedright
Quantity (t within the period, except end stock in t)
\end{minipage} & \begin{minipage}[b]{\linewidth}\raggedleft
Scenario A
\end{minipage} & \begin{minipage}[b]{\linewidth}\raggedleft
Scenario B
\end{minipage} \\
\midrule\noalign{}
\endhead
\bottomrule\noalign{}
\endlastfoot
Supply and origin-to-warehouse flow & 10 & 14 \\
Domestic demand and warehouse-to-destination flow & 8 & 12 \\
End stock & 2 & 2 \\
Shortage / static exceedance / reception overflow & 0 / 0 / 0 & 0 / 0 /
0 \\
\end{longtable}

The warehouse balances are \(0+10-8=2\) and \(0+14-12=2\). End stock is
below 5 t, and each scenario's receipts and shipments are below the
daily handling limits. A free terminal-stock policy permits the
remaining 2 t. The example demonstrates checks on a specified feasible
plan.

If each transport leg costs BRL1/t, receiving has no additional charge
and holding costs BRL1/(t period), the operating totals are
\(10+8+2=20\) and \(14+12+2=28\) BRL. Shared investment costs
\(5(4)=20\) BRL, so Equation C.1 gives \(20+(20+28)/2=44\) BRL. The
validator checks the two balances separately and reconstructs the same
investment, transport and holding totals. These illustrative values
explain the accounting procedure; the experimental input values are
reported in Appendix D.

\subsection{C6 GAMS translation and verification
limits}\label{c6-gams-translation-and-verification-limits}

This supplement provides \href{https://github.com/unb-lamfo-or-ai-hpc/model-agrologistic/blob/main/manuscript/supplement/appendix_c/model_det.gms}{model\_det.gms} and \href{https://github.com/unb-lamfo-or-ai-hpc/model-agrologistic/blob/main/manuscript/supplement/appendix_c/model_sto.gms}{model\_sto.gms} as entry
files, with \href{https://github.com/unb-lamfo-or-ai-hpc/model-agrologistic/blob/main/manuscript/supplement/appendix_c/model_common.inc}{model\_common.inc} containing the shared objective,
constraints and domains.
The \href{https://github.com/unb-lamfo-or-ai-hpc/model-agrologistic/blob/main/scripts/export_scalar_gams_data.py}{data-only exporter} writes already prepared
effective Python inputs to \href{https://github.com/unb-lamfo-or-ai-hpc/model-agrologistic/blob/main/manuscript/supplement/appendix_c/README.md}{data.inc} and produces an identifier and checksum record. This avoids manual transcription of the large arrays
while retaining the mapping to original facilities, products, periods
and scenarios.

The translation supports the declared scalable-candidate,
daily-handling-factor, free-terminal and three-enabled-slack
configuration. Unsupported settings are rejected. Closed-candidate
indicators require the solver-specific mapping described by
\href{https://www.gams.com/latest/docs/UG_LanguageFeatures.html\#UG_LanguageFeatures_IndicatorConstraints}{GAMS};
omitting that mapping would change the model. The provided source does
not replace it with an arbitrary slack bound.

Structural inspection and exporter tests examine indexing, units, costs,
scenarios and input preservation. The GAMS translation is provided for reference only and it has not
been compiled or numerically validated in this study. Numerical verification requires comparing objective components and original-unit residuals on the same
small input; alternative optimal decision vectors can represent the same scalar value.

\subsection{C7 Implementation index and reproduction
sequence}\label{c7-implementation-index-and-reproduction-sequence}

Table C.4 groups the principal reading links by scientific task. Paths
are relative to the repository root: mathematical software is under
src/logic, executable public examples under tests, and the GAMS
translation under manuscript/supplement/appendix\_c.~The links identify
a fixed reading snapshot. Each numerical result remains associated with
its original executed source.

\begin{longtable}[]{@{}
  >{\raggedright\arraybackslash}p{(\linewidth - 2\tabcolsep) * \real{0.5000}}
  >{\raggedright\arraybackslash}p{(\linewidth - 2\tabcolsep) * \real{0.5000}}@{}}
\caption{Table C.4. Implementation and verification
index.}\tabularnewline
\toprule\noalign{}
\begin{minipage}[b]{\linewidth}\raggedright
Task
\end{minipage} & \begin{minipage}[b]{\linewidth}\raggedright
Repository files
\end{minipage} \\
\midrule\noalign{}
\endfirsthead
\toprule\noalign{}
\begin{minipage}[b]{\linewidth}\raggedright
Task
\end{minipage} & \begin{minipage}[b]{\linewidth}\raggedright
Repository files
\end{minipage} \\
\midrule\noalign{}
\endhead
\bottomrule\noalign{}
\endlastfoot
Displayed equations and effective inputs &
\href{https://github.com/unb-lamfo-or-ai-hpc/model-agrologistic/blob/c80dca69c247753d49f06925ec712fcd43bb5a56/manuscript/formulation.qmd}{formulation.qmd};
\href{https://github.com/unb-lamfo-or-ai-hpc/model-agrologistic/blob/c80dca69c247753d49f06925ec712fcd43bb5a56/src/logic/mathematical_contract.py}{mathematical\_contract.py} \\
Native Gurobi &
\href{https://github.com/unb-lamfo-or-ai-hpc/model-agrologistic/blob/c80dca69c247753d49f06925ec712fcd43bb5a56/src/logic/optimization_gurobipy_stochastic.py}{optimization\_gurobipy\_stochastic.py} \\
Common Pyomo algebra and bounds &
\href{https://github.com/unb-lamfo-or-ai-hpc/model-agrologistic/blob/c80dca69c247753d49f06925ec712fcd43bb5a56/src/logic/optimization_pyomo.py}{optimization\_pyomo.py};
\href{https://github.com/unb-lamfo-or-ai-hpc/model-agrologistic/blob/c80dca69c247753d49f06925ec712fcd43bb5a56/src/logic/pyomo_bounds.py}{pyomo\_bounds.py} \\
Ordered common linear algebra &
\href{https://github.com/unb-lamfo-or-ai-hpc/model-agrologistic/blob/c80dca69c247753d49f06925ec712fcd43bb5a56/src/logic/pyomo_linear_matrix.py}{pyomo\_linear\_matrix.py} \\
Independent residual and cost validation &
\href{https://github.com/unb-lamfo-or-ai-hpc/model-agrologistic/blob/c80dca69c247753d49f06925ec712fcd43bb5a56/src/logic/solution_validation.py}{solution\_validation.py} \\
GAMS entry points and common equations &
\href{https://github.com/unb-lamfo-or-ai-hpc/model-agrologistic/blob/c80dca69c247753d49f06925ec712fcd43bb5a56/manuscript/supplement/appendix_c/model_det.gms}{model\_det.gms};
\href{https://github.com/unb-lamfo-or-ai-hpc/model-agrologistic/blob/c80dca69c247753d49f06925ec712fcd43bb5a56/manuscript/supplement/appendix_c/model_sto.gms}{model\_sto.gms};
\href{https://github.com/unb-lamfo-or-ai-hpc/model-agrologistic/blob/c80dca69c247753d49f06925ec712fcd43bb5a56/manuscript/supplement/appendix_c/model_common.inc}{model\_common.inc} \\
Data-only GAMS export and instructions &
\href{https://github.com/unb-lamfo-or-ai-hpc/model-agrologistic/blob/c80dca69c247753d49f06925ec712fcd43bb5a56/scripts/export_scalar_gams_data.py}{export\_scalar\_gams\_data.py};
\href{https://github.com/unb-lamfo-or-ai-hpc/model-agrologistic/blob/c80dca69c247753d49f06925ec712fcd43bb5a56/manuscript/supplement/appendix_c/README.md}{README.md} \\
Public mathematical / Excel examples &
\href{https://github.com/unb-lamfo-or-ai-hpc/model-agrologistic/blob/c80dca69c247753d49f06925ec712fcd43bb5a56/tests/test_scalar_contract.py}{test\_scalar\_contract.py};
\href{https://github.com/unb-lamfo-or-ai-hpc/model-agrologistic/blob/c80dca69c247753d49f06925ec712fcd43bb5a56/tests/test_excel_loader.py}{test\_excel\_loader.py} \\
Environment and adapter instructions &
\href{https://github.com/unb-lamfo-or-ai-hpc/model-agrologistic/blob/c80dca69c247753d49f06925ec712fcd43bb5a56/README.md}{README.md};
\href{https://github.com/unb-lamfo-or-ai-hpc/model-agrologistic/blob/c80dca69c247753d49f06925ec712fcd43bb5a56/docs/pyomo_native_adapters.md}{pyomo\_native\_adapters.md} \\
\end{longtable}

The index separates data preparation, model construction, solver
transfer and validation. A reader can therefore follow one physical
relationship from its effective input through the builder to the
residual check, rather than treating a solver's status as the complete
numerical assessment.

Begin with the equations and the small accounting example above. In an
existing compatible Python environment, public input examples can be
examined with:

\begin{Shaded}
\begin{Highlighting}[]
\ExtensionTok{python} \AttributeTok{{-}m}\NormalTok{ pytest tests/test\_excel\_loader.py }\AttributeTok{{-}q}
\end{Highlighting}
\end{Shaded}

Then obtain the exact processed workbook, population and network
definition described in Appendix D. Record the source, scenario design,
solver/interface settings and stopping rules for a fresh calculation.
Validate the returned decisions in original units, reconstruct all
fourteen cost components and compare the incumbent and bound using the
declared gap definition.

For the GAMS translation, the exporter writes already prepared effective
arrays to data.inc and records identifier mappings and checksums in
data\_manifest.json. Its inputs are a version-2 case manifest, a case
index, a workbook and expected hash, and a new output directory. The
resulting data remain bound to that workbook and case. The compilation
and numerical-verification boundary is stated in Section C6.

\subsection{C8 Consolidated common-Pyomo
results}\label{c8-consolidated-common-pyomo-results}

The NPAD
\href{https://github.com/unb-lamfo-or-ai-hpc/model-agrologistic/blob/docs/editorial-reconciliation-20261105/docs/evidence/pyomo_serial_final_export_20261109/configurations.csv}{configuration
ledger} and
\href{https://github.com/unb-lamfo-or-ai-hpc/model-agrologistic/blob/docs/editorial-reconciliation-20261105/docs/evidence/pyomo_serial_final_export_20261109/results.csv}{result
ledger} join one-to-one by the 90 planned cell identifiers. The
\href{https://github.com/unb-lamfo-or-ai-hpc/model-agrologistic/blob/docs/editorial-reconciliation-20261105/docs/evidence/pyomo_serial_final_export_20261109/final_result.json}{final
receipt} binds both tables by SHA-256. The
\href{https://github.com/unb-lamfo-or-ai-hpc/model-agrologistic/blob/docs/editorial-reconciliation-20261105/docs/evidence/pyomo_serial_final_export_20261109/reconciliation.json}{local
reconciliation} records the verified transfer and pilot
checks. The original NPAD collector verified the archives used for the
final export. The local check compared the transferred tables and
receipts with those identities.

\subsubsection{C8.1 Repeated h300 W STO
comparison}\label{c8.1-repeated-h300-w-sto-comparison}

All 30 h300 cells returned observations, including the four original
pilot cells. Gurobi reached the 1\% target in 15 of 15 runs. HiGHS
reached it in all three one-thread runs. The 12 HiGHS runs at 2, 4, 8
and 16 threads ended without a usable incumbent. Table C.5 reports the
three repetitions per setting, with median
(minimum--maximum) over three repetitions. Native and end-to-end times describe termination. The dash denotes an unavailable incumbent-based gap. The gap is reported for validated incumbents and the resource observations cover every executed run.

\begin{longtable}[]{@{}
  >{\raggedright\arraybackslash}p{(\linewidth - 14\tabcolsep) * \real{0.1250}}
  >{\raggedleft\arraybackslash}p{(\linewidth - 14\tabcolsep) * \real{0.1250}}
  >{\raggedleft\arraybackslash}p{(\linewidth - 14\tabcolsep) * \real{0.1250}}
  >{\raggedleft\arraybackslash}p{(\linewidth - 14\tabcolsep) * \real{0.1250}}
  >{\raggedleft\arraybackslash}p{(\linewidth - 14\tabcolsep) * \real{0.1250}}
  >{\raggedleft\arraybackslash}p{(\linewidth - 14\tabcolsep) * \real{0.1250}}
  >{\raggedleft\arraybackslash}p{(\linewidth - 14\tabcolsep) * \real{0.1250}}
  >{\raggedleft\arraybackslash}p{(\linewidth - 14\tabcolsep) * \real{0.1250}}@{}}
\caption{Table C.5. Final h300 W STO outcomes}\tabularnewline
\toprule\noalign{}
\begin{minipage}[b]{\linewidth}\raggedright
Solver
\end{minipage} & \begin{minipage}[b]{\linewidth}\raggedleft
Threads
\end{minipage} & \begin{minipage}[b]{\linewidth}\raggedleft
At target
\end{minipage} & \begin{minipage}[b]{\linewidth}\raggedleft
Without incumbent
\end{minipage} & \begin{minipage}[b]{\linewidth}\raggedleft
Common gap, \%
\end{minipage} & \begin{minipage}[b]{\linewidth}\raggedleft
Native terminal time, min
\end{minipage} & \begin{minipage}[b]{\linewidth}\raggedleft
End-to-end time, min
\end{minipage} & \begin{minipage}[b]{\linewidth}\raggedleft
Sampled worker-tree peak, GiB
\end{minipage} \\
\midrule\noalign{}
\endfirsthead
\toprule\noalign{}
\begin{minipage}[b]{\linewidth}\raggedright
Solver
\end{minipage} & \begin{minipage}[b]{\linewidth}\raggedleft
Threads
\end{minipage} & \begin{minipage}[b]{\linewidth}\raggedleft
At target
\end{minipage} & \begin{minipage}[b]{\linewidth}\raggedleft
Without incumbent
\end{minipage} & \begin{minipage}[b]{\linewidth}\raggedleft
Common gap, \%
\end{minipage} & \begin{minipage}[b]{\linewidth}\raggedleft
Native terminal time, min
\end{minipage} & \begin{minipage}[b]{\linewidth}\raggedleft
End-to-end time, min
\end{minipage} & \begin{minipage}[b]{\linewidth}\raggedleft
Sampled worker-tree peak, GiB
\end{minipage} \\
\midrule\noalign{}
\endhead
\bottomrule\noalign{}
\endlastfoot
Gurobi & 1 & 3 & 0 & 0.004105 (0.004105--0.004105) & 48.23
(48.15--61.50) & 91.97 (91.79--105.90) & 74.42 (74.21--74.54) \\
Gurobi & 2 & 3 & 0 & 0.004122 (0.004122--0.004122) & 66.18
(51.93--68.60) & 108.94 (96.10--111.42) & 82.51 (82.51--82.66) \\
Gurobi & 4 & 3 & 0 & 0.004120 (0.004120--0.004120) & 65.14
(56.40--76.99) & 113.14 (100.24--120.39) & 100.65 (99.35--100.86) \\
Gurobi & 8 & 3 & 0 & 0.004120 (0.004120--0.004120) & 93.29
(93.11--93.85) & 136.00 (135.82--137.85) & 99.52 (99.50--100.78) \\
Gurobi & 16 & 3 & 0 & 0.004120 (0.004120--0.004120) & 96.39
(71.90--101.34) & 141.15 (115.08--144.87) & 100.88 (100.26--103.36) \\
HiGHS & 1 & 3 & 0 & 0.006879 (0.006879--0.006879) & 401.43
(392.37--404.08) & 443.35 (433.55--444.57) & 84.82 (84.26--85.15) \\
HiGHS & 2 & 0 & 3 & --- & 480.79 (480.62--480.83) & 513.14
(512.26--514.54) & 89.50 (89.48--90.29) \\
HiGHS & 4 & 0 & 3 & --- & 480.61 (480.60--480.80) & 513.04
(512.52--518.63) & 90.68 (90.32--90.89) \\
HiGHS & 8 & 0 & 3 & --- & 480.62 (480.62--480.63) & 513.16
(512.71--513.35) & 90.52 (88.97--91.24) \\
HiGHS & 16 & 0 & 3 & --- & 480.60 (480.60--480.63) & 512.78
(512.67--512.85) & 90.39 (90.26--90.50) \\
\end{longtable}

For each solver and thread ceiling, select the three h300 W STO rows
from the result ledger by cell identifier and keep the original
validation and classification. Divide each timing value in seconds by 60
and each sampled peak in bytes by \(2^{30}\). Compute medians and
extrema on the unrounded values. The attainment denominator is three for
every h300 setting. The gaps in Table C.5 are expressed in percent and
use only rows with an accepted incumbent. A run without an incumbent
contributes its observed time and memory but no gap.

For h300 Gurobi, Table C.6 compares the median time at each thread
ceiling with the one-thread baseline. Each
median uses three originally validated runs at the common 1\% target. The attained gaps are given in Table C.5.
It gives terminal-time ratios
\(S(p)=T(1)/T(p)\) under the common stopping policy and nominal
efficiency \(S(p)/p\). The final export contains terminal clocks.
First-attainment timing and measured CPU utilization would require
separate evidence. The HiGHS configurations above one thread are
reported through their observed time and outcome counts in Table C.5.

\begin{longtable}[]{@{}
  >{\raggedleft\arraybackslash}p{(\linewidth - 8\tabcolsep) * \real{0.2000}}
  >{\raggedleft\arraybackslash}p{(\linewidth - 8\tabcolsep) * \real{0.2000}}
  >{\raggedleft\arraybackslash}p{(\linewidth - 8\tabcolsep) * \real{0.2000}}
  >{\raggedleft\arraybackslash}p{(\linewidth - 8\tabcolsep) * \real{0.2000}}
  >{\raggedleft\arraybackslash}p{(\linewidth - 8\tabcolsep) * \real{0.2000}}@{}}
\caption{Table C.6. Descriptive h300 Gurobi terminal-time ratios}\tabularnewline
\toprule\noalign{}
\begin{minipage}[b]{\linewidth}\raggedleft
Threads
\end{minipage} & \begin{minipage}[b]{\linewidth}\raggedleft
Native \(S(p)\)
\end{minipage} & \begin{minipage}[b]{\linewidth}\raggedleft
Native \(S(p)/p\)
\end{minipage} & \begin{minipage}[b]{\linewidth}\raggedleft
End-to-end \(S(p)\)
\end{minipage} & \begin{minipage}[b]{\linewidth}\raggedleft
End-to-end \(S(p)/p\)
\end{minipage} \\
\midrule\noalign{}
\endfirsthead
\toprule\noalign{}
\begin{minipage}[b]{\linewidth}\raggedleft
Threads
\end{minipage} & \begin{minipage}[b]{\linewidth}\raggedleft
Native \(S(p)\)
\end{minipage} & \begin{minipage}[b]{\linewidth}\raggedleft
Native \(S(p)/p\)
\end{minipage} & \begin{minipage}[b]{\linewidth}\raggedleft
End-to-end \(S(p)\)
\end{minipage} & \begin{minipage}[b]{\linewidth}\raggedleft
End-to-end \(S(p)/p\)
\end{minipage} \\
\midrule\noalign{}
\endhead
\bottomrule\noalign{}
\endlastfoot
1 & 1.000 & 1.000 & 1.000 & 1.000 \\
2 & 0.729 & 0.364 & 0.844 & 0.422 \\
4 & 0.740 & 0.185 & 0.813 & 0.203 \\
8 & 0.517 & 0.065 & 0.676 & 0.085 \\
16 & 0.500 & 0.031 & 0.652 & 0.041 \\
\end{longtable}

The
\href{https://github.com/unb-lamfo-or-ai-hpc/model-agrologistic/blob/dc59fca89d2b437c5066bedca1e6d45eb70cc141/docs/evidence/pyomo_partial_readback_review_20261105.json}{5
November readback} remains available as a dated preview. The final
ledger supplies the repetition-level values used in these tables.

\subsubsection{C8.2 h400 pilots and the memory
gate}\label{c8.2-h400-pilots-and-the-memory-gate}

All four pilots used the same input bindings, ordered linear matrix and
derived bounds, hardware class and worker affinity. The common ordered
matrix has 42,106,944 variables, including 492 binary variables,
1,435,062 linear constraints and 161,245,766 nonzeros. Both Gurobi
incumbents were independently accepted and share the same penalized
objective, BRL 7.593395 trillion, and economic subtotal, BRL 381.573
billion. HiGHS reached the native time limit in both pilots, with a
retained lower bound and absent incumbent and gap.

\begin{longtable}[]{@{}
  >{\raggedright\arraybackslash}p{(\linewidth - 14\tabcolsep) * \real{0.1250}}
  >{\raggedleft\arraybackslash}p{(\linewidth - 14\tabcolsep) * \real{0.1250}}
  >{\raggedleft\arraybackslash}p{(\linewidth - 14\tabcolsep) * \real{0.1250}}
  >{\raggedright\arraybackslash}p{(\linewidth - 14\tabcolsep) * \real{0.1250}}
  >{\raggedleft\arraybackslash}p{(\linewidth - 14\tabcolsep) * \real{0.1250}}
  >{\raggedleft\arraybackslash}p{(\linewidth - 14\tabcolsep) * \real{0.1250}}
  >{\raggedleft\arraybackslash}p{(\linewidth - 14\tabcolsep) * \real{0.1250}}
  >{\raggedleft\arraybackslash}p{(\linewidth - 14\tabcolsep) * \real{0.1250}}@{}}
\caption{Table C.7. Individual h400 W STO pilot outcomes and resource
measurements.}\tabularnewline
\toprule\noalign{}
\begin{minipage}[b]{\linewidth}\raggedright
Solver
\end{minipage} & \begin{minipage}[b]{\linewidth}\raggedleft
Threads
\end{minipage} & \begin{minipage}[b]{\linewidth}\raggedleft
Job
\end{minipage} & \begin{minipage}[b]{\linewidth}\raggedright
Outcome
\end{minipage} & \begin{minipage}[b]{\linewidth}\raggedleft
Common gap, \%
\end{minipage} & \begin{minipage}[b]{\linewidth}\raggedleft
Native terminal time, min
\end{minipage} & \begin{minipage}[b]{\linewidth}\raggedleft
End-to-end time, min
\end{minipage} & \begin{minipage}[b]{\linewidth}\raggedleft
Sampled worker-tree peak, GiB
\end{minipage} \\
\midrule\noalign{}
\endfirsthead
\toprule\noalign{}
\begin{minipage}[b]{\linewidth}\raggedright
Solver
\end{minipage} & \begin{minipage}[b]{\linewidth}\raggedleft
Threads
\end{minipage} & \begin{minipage}[b]{\linewidth}\raggedleft
Job
\end{minipage} & \begin{minipage}[b]{\linewidth}\raggedright
Outcome
\end{minipage} & \begin{minipage}[b]{\linewidth}\raggedleft
Common gap, \%
\end{minipage} & \begin{minipage}[b]{\linewidth}\raggedleft
Native terminal time, min
\end{minipage} & \begin{minipage}[b]{\linewidth}\raggedleft
End-to-end time, min
\end{minipage} & \begin{minipage}[b]{\linewidth}\raggedleft
Sampled worker-tree peak, GiB
\end{minipage} \\
\midrule\noalign{}
\endhead
\bottomrule\noalign{}
\endlastfoot
Gurobi & 1 & 2232379 & Validated at target & 0.836784 & 64.98 & 138.48 &
113.41 \\
Gurobi & 16 & 2232382 & Validated at target & 0.836784 & 130.57 & 205.33
& 160.24 \\
HiGHS & 1 & 2232380 & Time limit, no incumbent & --- & 481.04 & 535.14 &
126.38 \\
HiGHS & 16 & 2232381 & Time limit, no incumbent & --- & 481.78 & 542.97
& 150.91 \\
\end{longtable}

The h400 pilot control (job 2232383) rejected progression. Replaying the
pinned gate on the four matching receipts reproduced the memory-headroom
error. The gate allows a sampled peak of at most 80\% of the requested
192 GiB, or 153.60 GiB. The Gurobi 16-thread pilot (job 2232382) reached
160.24 GiB, leaving 16.54\% headroom. The other three sampled peaks were
below the gate ceiling. All four pilot records passed the
operational-closure checks. The Gurobi pilot retains its accepted
numerical result. The memory measure is sampled worker-tree peak RSS.

The scheduler cancelled 26 remaining h400 cells and all 30 h500 cells.
Their numerical fields remain empty in the final ledger. The h400 pilot
times are individual observations. The comparison across populations
therefore uses repeated Pyomo evidence at h300, individual Pyomo
evidence at h400, and a separate historical native Gurobi series through
h500.

\subsection{C9 h400 follow-up evidence and population
extension}\label{c9-h400-follow-up-evidence-and-population-extension}

The original gate in C8.2 used a sampled worker-tree peak ceiling of
153.60 GiB, equal to 80\% of the 192 GiB allocation. The E1 and E2
follow-up kept the 192 GiB request and checked both sampled worker-tree
RSS and the Slurm job cgroup v1 peak. The applicable ceilings were
163.20 and 172.80 GiB, respectively. The monitored scopes have different
meanings from the Slurm step MaxRSS.

The
\href{https://github.com/unb-lamfo-or-ai-hpc/model-agrologistic/blob/docs/editorial-reconciliation-20261105/docs/evidence/h400_e1_e2_final_20261112/verification.json}{verified
follow-up readback} binds the original E1 closing gate, the six accepted
E2 gates, the collection and memory receipts, and the
\href{https://github.com/unb-lamfo-or-ai-hpc/model-agrologistic/blob/docs/editorial-reconciliation-20261105/docs/evidence/h400_e1_e2_final_20261112/h400-e1-e2.csv}{ten-cell
comparison CSV}. The E1 and E2 run identifiers remain distinct. Their
shared ordered matrix, bounds, hardware class and input fingerprints
permit a descriptive comparison under the follow-up profile. The
original full archives remain in the NPAD claims. Table C.8 presents h400 follow-up terminal times and memory peaks.
Gurobi returned a validated incumbent at the 1\% target for every ceiling. HiGHS reached the time limit without a usable incumbent in every run.

\begin{longtable}[]{@{}
  >{\raggedright\arraybackslash}p{(\linewidth - 12\tabcolsep) * \real{0.1429}}
  >{\raggedleft\arraybackslash}p{(\linewidth - 12\tabcolsep) * \real{0.1429}}
  >{\raggedright\arraybackslash}p{(\linewidth - 12\tabcolsep) * \real{0.1429}}
  >{\raggedleft\arraybackslash}p{(\linewidth - 12\tabcolsep) * \real{0.1429}}
  >{\raggedleft\arraybackslash}p{(\linewidth - 12\tabcolsep) * \real{0.1429}}
  >{\raggedleft\arraybackslash}p{(\linewidth - 12\tabcolsep) * \real{0.1429}}
  >{\raggedleft\arraybackslash}p{(\linewidth - 12\tabcolsep) * \real{0.1429}}@{}}
\caption{Table C.8. h400 follow-up terminal times and memory peaks}\tabularnewline
\toprule\noalign{}
\begin{minipage}[b]{\linewidth}\raggedright
Solver
\end{minipage} & \begin{minipage}[b]{\linewidth}\raggedleft
Threads
\end{minipage} & \begin{minipage}[b]{\linewidth}\raggedright
Stage / job
\end{minipage} & \begin{minipage}[b]{\linewidth}\raggedleft
Native, min
\end{minipage} & \begin{minipage}[b]{\linewidth}\raggedleft
End-to-end, min
\end{minipage} & \begin{minipage}[b]{\linewidth}\raggedleft
Worker peak, GiB
\end{minipage} & \begin{minipage}[b]{\linewidth}\raggedleft
Cgroup peak, GiB
\end{minipage} \\
\midrule\noalign{}
\endfirsthead
\toprule\noalign{}
\begin{minipage}[b]{\linewidth}\raggedright
Solver
\end{minipage} & \begin{minipage}[b]{\linewidth}\raggedleft
Threads
\end{minipage} & \begin{minipage}[b]{\linewidth}\raggedright
Stage / job
\end{minipage} & \begin{minipage}[b]{\linewidth}\raggedleft
Native, min
\end{minipage} & \begin{minipage}[b]{\linewidth}\raggedleft
End-to-end, min
\end{minipage} & \begin{minipage}[b]{\linewidth}\raggedleft
Worker peak, GiB
\end{minipage} & \begin{minipage}[b]{\linewidth}\raggedleft
Cgroup peak, GiB
\end{minipage} \\
\midrule\noalign{}
\endhead
\bottomrule\noalign{}
\endlastfoot
Gurobi & 1 & E1 / 2247718 & 87.36 & 162.60 & 113.31 & 113.63 \\
Gurobi & 2 & E2 / 2250119 & 80.03 & 154.83 & 133.13 & 133.46 \\
Gurobi & 4 & E2 / 2250123 & 139.63 & 214.98 & 162.10 & 162.89 \\
Gurobi & 8 & E2 / 2250127 & 150.36 & 226.35 & 161.47 & 162.94 \\
Gurobi & 16 & E1 / 2247727 & 134.60 & 207.78 & 162.99 & 163.45 \\
HiGHS & 1 & E1 / 2247723 & 481.13 & 534.88 & 124.82 & 126.62 \\
HiGHS & 2 & E2 / 2250121 & 481.39 & 535.59 & 151.33 & 151.76 \\
HiGHS & 4 & E2 / 2250125 & 481.35 & 536.08 & 150.27 & 151.65 \\
HiGHS & 8 & E2 / 2250129 & 481.37 & 536.96 & 151.28 & 151.68 \\
HiGHS & 16 & E1 / 2247725 & 481.28 & 534.81 & 149.67 & 151.52 \\
\end{longtable}

The five Gurobi gaps round to 0.836784\% and the reported incumbent
objective and economic subtotal coincide across the grid. HiGHS has a
finite reported lower bound of BRL 3,917,886.51 in each return, with
empty incumbent objective and gap fields. The ten cells passed their
operational criteria, including the applicable memory ceilings and the
cgroup counter checks. The closest sampled worker peak was within 0.21
GiB of its ceiling. This small margin concerns the observed run and
matters when planning repetitions.

The Gurobi ratios use the single E1 one-thread observation separately
for each clock. Dividing each ratio by the configured thread ceiling
gives native efficiencies of 100.0\%, 54.6\%, 15.6\%, 7.3\% and 4.1\%,
and end-to-end efficiencies of 100.0\%, 52.5\%, 18.9\%, 9.0\% and 4.9\%,
in ascending thread order. These are descriptive ratios of terminal
clocks. The HiGHS returns supply attainment, timing and resource
observations for all five settings, with no validated baseline for a
speedup calculation.

A future repeated h400 comparison would quantify variability at each
thread ceiling. The h500 extension would qualify its own resource
profile and retain the effective inputs, ordered matrix identity, solver
options, time budget and validation criteria for each cell.

\subsection{C10 Root relaxation and branch-and-bound diagnostics}\label{c10-root-relaxation-and-branch-and-bound-diagnostics}

The earlier Gurobi h300 W STO records all report one processed branch-and-bound node under the historical automatic method and heuristic settings. This count describes the realized search under its stopping policy; it does not establish a structural one-node bound on the MILP. In the historical common-Pyomo h300 runs, the matrix contains 25,107,544 variables (392 binaries) and 94,967,766 nonzero coefficients. The initial linear-programming (LP) relaxation can therefore be computationally substantial even when no extended search tree is observed. These observations motivate separating root-relaxation effort, incumbent attainment and gap certification. A small final MIP gap is a certificate for the penalized scalar objective, not for the financial subtotal or physical service. The controlled root-method experiment is a distinct exploratory cohort and is not combined with the historical thread-ceiling data in Sections C8--C9.

\subsection{C11 Closed h300 root-method block}\label{c11-partial-h300-root-method-experiment}

This appendix reports all 60 configurations in the closed first temporal block. The design crossed \texttt{Method=1,2,4,6}, thread ceilings of 1, 2, 4, 8 and 16 on distinct physical cores, and three profiles: standard (\texttt{Heuristics=0.05}, \texttt{NoRelHeurWork=0}); H0 (\texttt{Heuristics=0}, \texttt{NoRelHeurWork=0}); and NR (\texttt{Heuristics=0.05}, \texttt{NoRelHeurWork=600}). The common settings included Gurobi 13.0.3, seed 42, a 28,800-second native limit, a 1\% MIP-gap target, \texttt{ConcurrentMIP=1}, CPU-only execution and \texttt{SoftMemLimit=160} in Gurobi's memory units. Slurm requested 192 GiB per job. The retained affinity records identify the configured number of distinct physical cores. NoRel was controlled by deterministic work, rather than a wall-time heuristic budget.

The 49 Slurm-completed tasks comprised 33 accepted solutions, nine time-limited returns and seven native memory-limited returns. Seven other tasks stopped at the external memory guard. One PDHG/H0 execution at eight cores was cancelled during optimization, and three configurations were cancelled before starting. The remaining two temporal blocks, comprising 120 planned configurations, were not submitted. Tables~\ref{tab:c11-method} and~\ref{tab:c11-jobs} retain this coverage explicitly.

\begin{table}[htbp]\centering\small
\caption{Table~\thetable. Outcomes by method and heuristic profile in the closed first h300 block. Each row contains five planned core settings. GMem denotes Gurobi's native soft memory limit, Guard the external memory monitor, and Cancel researcher-requested cancellation.}
\label{tab:c11-method}
\begin{tabular}{llrrrrr}\toprule
Method & Profile & Accepted & Time & GMem & Guard & Cancel\\\midrule
1 & Std & 4 & 0 & 0 & 0 & 1\\
1 & H0 & 3 & 0 & 2 & 0 & 0\\
1 & NR & 3 & 0 & 0 & 1 & 1\\
2 & Std & 4 & 0 & 0 & 0 & 1\\
2 & H0 & 3 & 0 & 2 & 0 & 0\\
2 & NR & 3 & 0 & 0 & 2 & 0\\
4 & Std & 5 & 0 & 0 & 0 & 0\\
4 & H0 & 3 & 0 & 2 & 0 & 0\\
4 & NR & 3 & 0 & 0 & 2 & 0\\
6 & Std & 2 & 3 & 0 & 0 & 0\\
6 & H0 & 0 & 3 & 1 & 0 & 1\\
6 & NR & 0 & 3 & 0 & 2 & 0\\
\midrule Total & & 33 & 9 & 7 & 7 & 4\\\bottomrule
\end{tabular}\end{table}

The seven native logs end with \emph{Memory limit reached} and zero solutions, despite the application's original generic error label. They occur for H0 at eight or sixteen cores. Their original summary and rejected-validation records remain unchanged. All seven external guard stops occur for NR at eight or sixteen cores; their reported OOM-kill and fail-count deltas are zero. The maximum retained job-cgroup peak among guard stops was 184.10 GiB. These are resource-censored observations, and no feasible plan or service value is assigned to them. The cancelled running task lacks a final memory receipt and solver terminal clock; its Slurm step RSS remains a separate accounting observation.

\begingroup\small\setlength{\tabcolsep}{3pt}
\begin{longtable}{@{}rrlrlrrrr@{}}
\caption{Table~\thetable. All configurations in the closed h300 block. Root and native terminal times are minutes from native log events; memory is the job-cgroup peak in GiB. Time denotes no usable incumbent at the time limit. NA preserves unavailable events. Task is the original array index.}\label{tab:c11-jobs}\\
\toprule Task & Method & Profile & Cores & Outcome & Root & Native & Nodes & Peak\\\midrule\endfirsthead
\multicolumn{9}{l}{\small Closed h300 block (continued)}\\
\toprule Task & Method & Profile & Cores & Outcome & Root & Native & Nodes & Peak\\\midrule\endhead
0 & 2 & H0 & 8 & GRB mem & 69.5 & 93.9 & 17 & 153.31\\
1 & 6 & Std & 1 & Time & 473.1 & 480.3 & 1 & 70.93\\
2 & 4 & Std & 8 & Accepted & 47.5 & 65.3 & 1 & 101.40\\
3 & 2 & Std & 1 & Accepted & 122.3 & 137.8 & 1 & 74.48\\
4 & 2 & Std & 8 & Accepted & 68.4 & 85.3 & 1 & 85.19\\
5 & 6 & H0 & 2 & Time & 474.3 & 480.2 & 1 & 69.79\\
6 & 4 & NR & 8 & Guard & NA & NA & NA & 173.40\\
7 & 4 & Std & 4 & Accepted & 37.7 & 55.1 & 1 & 89.30\\
8 & 4 & H0 & 4 & Accepted & 39.9 & 63.1 & 39 & 112.19\\
9 & 4 & H0 & 8 & GRB mem & 48.1 & 72.6 & 28 & 154.40\\
10 & 1 & NR & 4 & Accepted & 26.3 & 87.9 & 43 & 120.32\\
11 & 1 & Std & 16 & Accepted & 29.4 & 50.1 & 1 & 82.29\\
12 & 6 & Std & 2 & Time & 474.0 & 480.3 & 1 & 70.93\\
13 & 4 & Std & 16 & Accepted & 48.8 & 66.5 & 1 & 103.94\\
14 & 6 & H0 & 4 & Time & 474.8 & 480.2 & 1 & 69.75\\
15 & 1 & NR & 2 & Accepted & 24.4 & 74.8 & 1 & 94.12\\
16 & 6 & NR & 2 & Time & 455.3 & 480.4 & 1 & 94.12\\
17 & 2 & NR & 8 & Guard & NA & NA & NA & 174.55\\
18 & 2 & Std & 2 & Accepted & 94.9 & 114.6 & 1 & 74.79\\
19 & 2 & NR & 16 & Guard & NA & NA & NA & 177.57\\
20 & 1 & H0 & 4 & Accepted & 27.3 & 51.5 & 47 & 113.96\\
21 & 6 & NR & 8 & Guard & NA & NA & NA & 173.99\\
22 & 2 & H0 & 2 & Accepted & 94.9 & 119.3 & 25 & 87.08\\
23 & 2 & H0 & 1 & Accepted & 115.9 & 138.5 & 19 & 74.97\\
24 & 4 & NR & 16 & Guard & NA & NA & NA & 174.82\\
25 & 4 & H0 & 1 & Accepted & 25.9 & 45.4 & 17 & 74.25\\
26 & 4 & NR & 4 & Accepted & 39.0 & 93.3 & 1 & 120.27\\
27 & 6 & Std & 16 & Accepted & 296.8 & 312.9 & 1 & 80.05\\
28 & 6 & NR & 1 & Time & 451.5 & 480.7 & 1 & 80.20\\
29 & 1 & Std & 2 & Accepted & 25.1 & 41.7 & 1 & 75.68\\
30 & 1 & Std & 8 & Accepted & 30.6 & 53.5 & 1 & 82.39\\
31 & 2 & H0 & 4 & Accepted & 84.4 & 110.0 & 39 & 113.79\\
32 & 6 & H0 & 1 & Time & 474.1 & 480.2 & 1 & 69.75\\
33 & 4 & H0 & 2 & Accepted & 35.9 & 57.5 & 27 & 88.67\\
34 & 2 & Std & 16 & Accepted & 61.3 & 78.1 & 1 & 85.38\\
35 & 6 & Std & 8 & Accepted & 343.9 & 360.4 & 1 & 82.41\\
36 & 4 & NR & 1 & Accepted & 26.9 & 81.6 & 1 & 80.19\\
37 & 6 & H0 & 16 & GRB mem & 296.8 & 318.6 & 10 & 162.29\\
38 & 4 & Std & 2 & Accepted & 35.1 & 50.8 & 1 & 83.02\\
39 & 1 & Std & 1 & Accepted & 26.3 & 41.1 & 1 & 75.73\\
40 & 1 & H0 & 2 & Accepted & 25.5 & 47.2 & 25 & 87.87\\
41 & 6 & Std & 4 & Time & 473.0 & 480.1 & 1 & 79.17\\
42 & 2 & NR & 2 & Accepted & 94.7 & 135.1 & 1 & 94.14\\
43 & 4 & H0 & 16 & GRB mem & 49.6 & 72.0 & 10 & 162.30\\
44 & 1 & NR & 16 & Guard & NA & NA & NA & 184.10\\
45 & 4 & Std & 1 & Accepted & 26.9 & 41.7 & 1 & 75.79\\
46 & 4 & NR & 2 & Accepted & 34.9 & 91.3 & 1 & 94.06\\
47 & 1 & H0 & 8 & GRB mem & 28.8 & 52.1 & 27 & 153.82\\
48 & 2 & NR & 1 & Accepted & 121.4 & 158.8 & 1 & 80.21\\
49 & 1 & NR & 1 & Accepted & 26.4 & 80.4 & 1 & 80.18\\
50 & 1 & H0 & 16 & GRB mem & 29.8 & 52.2 & 10 & 162.06\\
51 & 2 & NR & 4 & Accepted & 83.7 & 139.4 & 1 & 121.32\\
52 & 2 & H0 & 16 & GRB mem & 59.9 & 85.0 & 10 & 162.51\\
53 & 1 & H0 & 1 & Accepted & 26.6 & 47.5 & 17 & 74.39\\
54 & 6 & NR & 16 & Guard & NA & NA & NA & 180.59\\
55 & 6 & NR & 4 & Time & 449.8 & 480.1 & 1 & 121.30\\
56 & 6 & H0 & 8 & Cancel & NA & NA & NA & NA\\
57 & 1 & NR & 8 & Not run & NA & NA & NA & NA\\
58 & 1 & Std & 4 & Not run & NA & NA & NA & NA\\
59 & 2 & Std & 4 & Not run & NA & NA & NA & NA\\
\bottomrule\end{longtable}\endgroup

Root clocks for time-limited cases measure effort before interruption, not a completed LP solve. The original source, plan and matrix guard are preserved. Accepted certificates use absolute tolerance $10^{-5}$ and relative tolerance $10^{-8}$, with zero failed checks and objective reconstruction against the original inputs. Table~\ref{tab:c11-quality} reports the accepted scalar and agricultural outcomes. The scalar dual bound certifies the penalized objective; it supplies no separate economic-only MIP gap.

\begingroup\small\setlength{\tabcolsep}{3pt}
\begin{longtable}{@{}rrlr rrrr@{}}
\caption{Table~\thetable. Accepted h300 outcomes: total objective in trillion BRL, economic cost in billion BRL, scalar relative gap in percent and domestic service in percent. Std, H0 and NR are defined above. Economic-only gaps are unavailable.}\label{tab:c11-quality}\\
\toprule Task & Method & Profile & Cores & Objective & Economic & Gap & Service\\\midrule\endfirsthead
\multicolumn{8}{l}{\small Accepted h300 outcomes (continued)}\\
\toprule Task & Method & Profile & Cores & Objective & Economic & Gap & Service\\\midrule\endhead
2 & 4 & Std & 8 & 9.289109 & 385.502 & 0.004120 & 91.3256\\
3 & 2 & Std & 1 & 9.289109 & 385.502 & 0.004122 & 91.3256\\
4 & 2 & Std & 8 & 9.289109 & 385.502 & 0.004123 & 91.3256\\
7 & 4 & Std & 4 & 9.289109 & 385.502 & 0.004120 & 91.3256\\
8 & 4 & H0 & 4 & 9.288940 & 385.450 & 0.001350 & 91.3258\\
10 & 1 & NR & 4 & 9.289205 & 385.409 & 0.004199 & 91.3237\\
11 & 1 & Std & 16 & 9.289109 & 385.502 & 0.004105 & 91.3256\\
13 & 4 & Std & 16 & 9.289109 & 385.502 & 0.004120 & 91.3256\\
15 & 1 & NR & 2 & 9.288940 & 385.438 & 0.002283 & 91.3248\\
18 & 2 & Std & 2 & 9.289109 & 385.502 & 0.004107 & 91.3256\\
20 & 1 & H0 & 4 & 9.288940 & 385.450 & 0.001359 & 91.3258\\
22 & 2 & H0 & 2 & 9.289214 & 385.421 & 0.004904 & 91.3247\\
23 & 2 & H0 & 1 & 9.289359 & 385.486 & 0.006014 & 91.3252\\
25 & 4 & H0 & 1 & 9.289372 & 385.408 & 0.006175 & 91.3254\\
26 & 4 & NR & 4 & 9.288940 & 385.438 & 0.002284 & 91.3248\\
27 & 6 & Std & 16 & 9.289109 & 385.502 & 0.004122 & 91.3256\\
29 & 1 & Std & 2 & 9.289109 & 385.502 & 0.004120 & 91.3256\\
30 & 1 & Std & 8 & 9.289109 & 385.502 & 0.004105 & 91.3256\\
31 & 2 & H0 & 4 & 9.289205 & 385.409 & 0.004223 & 91.3237\\
33 & 4 & H0 & 2 & 9.289228 & 385.343 & 0.005053 & 91.3249\\
34 & 2 & Std & 16 & 9.289109 & 385.502 & 0.004123 & 91.3256\\
35 & 6 & Std & 8 & 9.289109 & 385.502 & 0.004123 & 91.3256\\
36 & 4 & NR & 1 & 9.289122 & 385.425 & 0.004248 & 91.3258\\
38 & 4 & Std & 2 & 9.289109 & 385.502 & 0.004122 & 91.3256\\
39 & 1 & Std & 1 & 9.289109 & 385.502 & 0.004120 & 91.3256\\
40 & 1 & H0 & 2 & 9.289214 & 385.421 & 0.004902 & 91.3247\\
42 & 2 & NR & 2 & 9.289109 & 385.502 & 0.004107 & 91.3256\\
45 & 4 & Std & 1 & 9.289109 & 385.502 & 0.004120 & 91.3256\\
46 & 4 & NR & 2 & 9.288940 & 385.438 & 0.002284 & 91.3248\\
48 & 2 & NR & 1 & 9.289109 & 385.502 & 0.004122 & 91.3256\\
49 & 1 & NR & 1 & 9.289122 & 385.425 & 0.004248 & 91.3258\\
51 & 2 & NR & 4 & 9.289205 & 385.409 & 0.005145 & 91.3237\\
53 & 1 & H0 & 1 & 9.289372 & 385.408 & 0.006175 & 91.3254\\
\bottomrule\end{longtable}\endgroup

Under standard heuristics, barrier's observed one-to-sixteen-core root-time ratio is 1.99 and its terminal speedup is 1.77; nominal terminal efficiency is 11.0\%. Dual simplex and deterministic concurrent optimization have terminal ratios below one at sixteen cores, while accepted PDHG observations lack an accepted one-core baseline. These are single-configuration observations in one temporal block. Accepted node counts range from one to 47. The results locate the method-dependent response in root relaxation and subsequent incumbent search, while differences in paths and hardware-level contention remain unisolated.

The final archive, all 14 transfer-ledger files and all 1,763 inventory member hashes were checked locally. The 49 expected-but-absent records belong to guard-stopped or cancelled configurations and remain listed in the technical evidence. These absences do not reduce the 33 accepted certificates. Complete vectors and raw receipts remain in custody; the repository contains aggregate results, identities and the audit record. The earlier 28-task checkpoint is retained as provenance and is superseded for final outcome counts.

\subsection{C12 Network utilization, model size and objective
certificates}\label{c12-network-utilization-model-size-and-objective-certificates}

The tables report Gurobi incumbents from the historical native and
Pyomo/APPSI cohorts. DET and STO denote deterministic and stochastic
formulations. W routes flow through warehouses and D also permits direct
delivery. T denotes the solver thread ceiling. Matrix counts refer to
the original MILP, including any formulation-specific rows before
presolve. Network nodes count origin, warehouse and destination roles.
Arcs count directed family and product coordinates before period and
scenario expansion. Used arcs count coordinates with positive exported
flow in at least one period and scenario.

Nonzero continuous entries use the retained export threshold of
\(10^{-7}\) and cover flow, inventory, unmet demand, capacity slack and
installed capacities. Binary counts exclude the availability flag of
existing warehouses. ND indicates that the required vector or
certificate was unavailable in the retrieved evidence. The economic
subtotal is evaluated on the incumbent and has no separate optimality
certificate. The same penalized MILP bound cannot be used as an
economic-objective bound. LP in the notes identifies the observed
initial root-relaxation time.

The h300 Pyomo logs report one explored node for all fifteen
observations, with the root LP objective printed at approximately BRL
9.28872 trillion. The matching incumbent has a penalized objective of
approximately BRL 9.28911 trillion. The archived solution vector
measures incumbent sparsity. The root LP logs record the objective,
iterations and time.

The
\href{https://github.com/unb-lamfo-or-ai-hpc/model-agrologistic/blob/f8ee519b45477c593930692e71601cf1f406896e/manuscript/scalar_evidence.json}{historical
scalar evidence} binds the native observations to their archived
certificates and output hashes. The
\href{https://github.com/unb-lamfo-or-ai-hpc/model-agrologistic/blob/0f63c10eb1abf88318a4d9c6afc75fcac97f6266/docs/evidence/h300_selected_graph_release_20261101.json}{h300
admissibility record} identifies the selected graph separately from the
exported flows.

In the native stochastic W series, used arcs rise from 4,090 to 7,023
from h215 to h500, while the share falls from 16.29\% to 6.09\%. The
available arc set expands faster than the observed set used at least
once.

\subsubsection{h215}\label{h215}

\textbf{Model and network}

{\def\LTcaptype{none} % do not increment counter
\begingroup\footnotesize
\begin{longtable}[]{@{}
  >{\raggedright\arraybackslash}p{(\linewidth - 14\tabcolsep) * \real{0.1250}}
  >{\raggedright\arraybackslash}p{(\linewidth - 14\tabcolsep) * \real{0.1250}}
  >{\raggedright\arraybackslash}p{(\linewidth - 14\tabcolsep) * \real{0.1250}}
  >{\raggedright\arraybackslash}p{(\linewidth - 14\tabcolsep) * \real{0.1250}}
  >{\raggedright\arraybackslash}p{(\linewidth - 14\tabcolsep) * \real{0.1250}}
  >{\raggedright\arraybackslash}p{(\linewidth - 14\tabcolsep) * \real{0.1250}}
  >{\raggedright\arraybackslash}p{(\linewidth - 14\tabcolsep) * \real{0.1250}}
  >{\raggedright\arraybackslash}p{(\linewidth - 14\tabcolsep) * \real{0.1250}}@{}}
\toprule\noalign{}
\begin{minipage}[b]{\linewidth}\raggedright
Setting
\end{minipage} & \begin{minipage}[b]{\linewidth}\raggedright
Rows
\end{minipage} & \begin{minipage}[b]{\linewidth}\raggedright
Columns
\end{minipage} & \begin{minipage}[b]{\linewidth}\raggedright
Nonzeros
\end{minipage} & \begin{minipage}[b]{\linewidth}\raggedright
Network nodes
\end{minipage} & \begin{minipage}[b]{\linewidth}\raggedright
Arcs
\end{minipage} & \begin{minipage}[b]{\linewidth}\raggedright
Used arcs (\%)
\end{minipage} & \begin{minipage}[b]{\linewidth}\raggedright
B\&B nodes
\end{minipage} \\
\midrule\noalign{}
\endhead
\bottomrule\noalign{}
\endlastfoot
DET W / Native / 4T & 74,017 & 1,562,174 & 5,786,346 & 289 & 25,112 &
12.01 & 1 \\
STO W / Native / 4T & 661,057 & 14,054,654 & 52,068,826 & 289 & 25,112 &
16.29 & 1 \\
DET D / Native / 4T & 74,017 & 1,597,694 & 5,857,386 & 289 & 25,704 &
11.34 & 1 \\
STO D / Native / 4T & 661,057 & 14,374,334 & 52,708,186 & 289 & 25,704 &
15.30 & 1 \\
STO W / Pyomo / 4T & 735,577 & 14,054,654 & 52,217,866 & 289 & 25,112 &
16.39 & 1 \\
\end{longtable}
\endgroup
}

\textbf{Incumbent, certificate and variable usage}

{\def\LTcaptype{none} % do not increment counter
\begingroup\footnotesize
\begin{longtable}[]{@{}
  >{\raggedright\arraybackslash}p{(\linewidth - 18\tabcolsep) * \real{0.1000}}
  >{\raggedright\arraybackslash}p{(\linewidth - 18\tabcolsep) * \real{0.1000}}
  >{\raggedright\arraybackslash}p{(\linewidth - 18\tabcolsep) * \real{0.1000}}
  >{\raggedright\arraybackslash}p{(\linewidth - 18\tabcolsep) * \real{0.1000}}
  >{\raggedright\arraybackslash}p{(\linewidth - 18\tabcolsep) * \real{0.1000}}
  >{\raggedright\arraybackslash}p{(\linewidth - 18\tabcolsep) * \real{0.1000}}
  >{\raggedright\arraybackslash}p{(\linewidth - 18\tabcolsep) * \real{0.1000}}
  >{\raggedright\arraybackslash}p{(\linewidth - 18\tabcolsep) * \real{0.1000}}
  >{\raggedright\arraybackslash}p{(\linewidth - 18\tabcolsep) * \real{0.1000}}
  >{\raggedright\arraybackslash}p{(\linewidth - 18\tabcolsep) * \real{0.1000}}@{}}
\toprule\noalign{}
\begin{minipage}[b]{\linewidth}\raggedright
Setting
\end{minipage} & \begin{minipage}[b]{\linewidth}\raggedright
Total objective (BRL tn)
\end{minipage} & \begin{minipage}[b]{\linewidth}\raggedright
Total gap (\%)
\end{minipage} & \begin{minipage}[b]{\linewidth}\raggedright
Economic cost (BRL bn)
\end{minipage} & \begin{minipage}[b]{\linewidth}\raggedright
Economic gap (\%)
\end{minipage} & \begin{minipage}[b]{\linewidth}\raggedright
Integers
\end{minipage} & \begin{minipage}[b]{\linewidth}\raggedright
Nonzero integers (\%)
\end{minipage} & \begin{minipage}[b]{\linewidth}\raggedright
Cont. vars
\end{minipage} & \begin{minipage}[b]{\linewidth}\raggedright
Nonzero continuous (\%)
\end{minipage} & \begin{minipage}[b]{\linewidth}\raggedright
Notes
\end{minipage} \\
\midrule\noalign{}
\endhead
\bottomrule\noalign{}
\endlastfoot
DET W / Native / 4T & 11.437342 & 0.000812 & 389.810 & ND & 307 & 70.03
& 1,561,867 & 1.962 & Full vector \\
STO W / Native / 4T & 13.071397 & 0.740159 & 393.902 & ND & 307 & 70.03
& 14,054,347 & 1.954 & Full vector \\
DET D / Native / 4T & 7.340492 & 0.552427 & 338.715 & ND & 307 & 70.03 &
1,597,387 & 1.449 & Full vector \\
STO D / Native / 4T & 7.355002 & 0.623718 & 340.109 & ND & 307 & 70.03 &
14,374,027 & 1.451 & Full vector \\
STO W / Pyomo / 4T & 13.071397 & 0.740159 & 393.902 & ND & 307 & 70.03 &
14,054,347 & 1.946 & Full vector \\
\end{longtable}
\endgroup
}

\subsubsection{h300}\label{h300}

\textbf{Model and network}

{\def\LTcaptype{none} % do not increment counter
\begingroup\footnotesize
\begin{longtable}[]{@{}
  >{\raggedright\arraybackslash}p{(\linewidth - 14\tabcolsep) * \real{0.1250}}
  >{\raggedright\arraybackslash}p{(\linewidth - 14\tabcolsep) * \real{0.1250}}
  >{\raggedright\arraybackslash}p{(\linewidth - 14\tabcolsep) * \real{0.1250}}
  >{\raggedright\arraybackslash}p{(\linewidth - 14\tabcolsep) * \real{0.1250}}
  >{\raggedright\arraybackslash}p{(\linewidth - 14\tabcolsep) * \real{0.1250}}
  >{\raggedright\arraybackslash}p{(\linewidth - 14\tabcolsep) * \real{0.1250}}
  >{\raggedright\arraybackslash}p{(\linewidth - 14\tabcolsep) * \real{0.1250}}
  >{\raggedright\arraybackslash}p{(\linewidth - 14\tabcolsep) * \real{0.1250}}@{}}
\toprule\noalign{}
\begin{minipage}[b]{\linewidth}\raggedright
Setting
\end{minipage} & \begin{minipage}[b]{\linewidth}\raggedright
Rows
\end{minipage} & \begin{minipage}[b]{\linewidth}\raggedright
Columns
\end{minipage} & \begin{minipage}[b]{\linewidth}\raggedright
Nonzeros
\end{minipage} & \begin{minipage}[b]{\linewidth}\raggedright
Network nodes
\end{minipage} & \begin{minipage}[b]{\linewidth}\raggedright
Arcs
\end{minipage} & \begin{minipage}[b]{\linewidth}\raggedright
Used arcs (\%)
\end{minipage} & \begin{minipage}[b]{\linewidth}\raggedright
B\&B nodes
\end{minipage} \\
\midrule\noalign{}
\endhead
\bottomrule\noalign{}
\endlastfoot
STO W / Native / 1T & 890,642 & 25,107,544 & 94,635,126 & 374 & 45,240 &
11.18 & 1 \\
STO W / Native / 2T & 890,642 & 25,107,544 & 94,635,126 & 374 & 45,240 &
11.18 & 1 \\
STO W / Native / 4T & 890,642 & 25,107,544 & 94,635,126 & 374 & 45,240 &
11.18 & 1 \\
STO W / Native / 8T & 890,642 & 25,107,544 & 94,635,126 & 374 & 45,240 &
11.18 & 1 \\
STO W / Native / 16T & 890,642 & 25,107,544 & 94,635,126 & 374 & 45,240
& 11.18 & 1 \\
STO W / Pyomo / 1T r1 & 1,056,962 & 25,107,544 & 94,967,766 & 374 &
45,240 & 11.23 & 1 \\
STO W / Pyomo / 1T r2 & 1,056,962 & 25,107,544 & 94,967,766 & 374 &
45,240 & ND & 1 \\
STO W / Pyomo / 1T r3 & 1,056,962 & 25,107,544 & 94,967,766 & 374 &
45,240 & ND & 1 \\
STO W / Pyomo / 2T r1 & 1,056,962 & 25,107,544 & 94,967,766 & 374 &
45,240 & ND & 1 \\
STO W / Pyomo / 2T r2 & 1,056,962 & 25,107,544 & 94,967,766 & 374 &
45,240 & ND & 1 \\
STO W / Pyomo / 2T r3 & 1,056,962 & 25,107,544 & 94,967,766 & 374 &
45,240 & ND & 1 \\
STO W / Pyomo / 4T r1 & 1,056,962 & 25,107,544 & 94,967,766 & 374 &
45,240 & ND & 1 \\
STO W / Pyomo / 4T r2 & 1,056,962 & 25,107,544 & 94,967,766 & 374 &
45,240 & ND & 1 \\
STO W / Pyomo / 4T r3 & 1,056,962 & 25,107,544 & 94,967,766 & 374 &
45,240 & ND & 1 \\
STO W / Pyomo / 8T r1 & 1,056,962 & 25,107,544 & 94,967,766 & 374 &
45,240 & ND & 1 \\
STO W / Pyomo / 8T r2 & 1,056,962 & 25,107,544 & 94,967,766 & 374 &
45,240 & ND & 1 \\
STO W / Pyomo / 8T r3 & 1,056,962 & 25,107,544 & 94,967,766 & 374 &
45,240 & ND & 1 \\
STO W / Pyomo / 16T r1 & 1,056,962 & 25,107,544 & 94,967,766 & 374 &
45,240 & ND & 1 \\
STO W / Pyomo / 16T r2 & 1,056,962 & 25,107,544 & 94,967,766 & 374 &
45,240 & ND & 1 \\
STO W / Pyomo / 16T r3 & 1,056,962 & 25,107,544 & 94,967,766 & 374 &
45,240 & ND & 1 \\
\end{longtable}
\endgroup
}

\textbf{Incumbent, certificate and variable usage}

{\def\LTcaptype{none} % do not increment counter
\begingroup\footnotesize
\begin{longtable}[]{@{}
  >{\raggedright\arraybackslash}p{(\linewidth - 18\tabcolsep) * \real{0.1000}}
  >{\raggedright\arraybackslash}p{(\linewidth - 18\tabcolsep) * \real{0.1000}}
  >{\raggedright\arraybackslash}p{(\linewidth - 18\tabcolsep) * \real{0.1000}}
  >{\raggedright\arraybackslash}p{(\linewidth - 18\tabcolsep) * \real{0.1000}}
  >{\raggedright\arraybackslash}p{(\linewidth - 18\tabcolsep) * \real{0.1000}}
  >{\raggedright\arraybackslash}p{(\linewidth - 18\tabcolsep) * \real{0.1000}}
  >{\raggedright\arraybackslash}p{(\linewidth - 18\tabcolsep) * \real{0.1000}}
  >{\raggedright\arraybackslash}p{(\linewidth - 18\tabcolsep) * \real{0.1000}}
  >{\raggedright\arraybackslash}p{(\linewidth - 18\tabcolsep) * \real{0.1000}}
  >{\raggedright\arraybackslash}p{(\linewidth - 18\tabcolsep) * \real{0.1000}}@{}}
\toprule\noalign{}
\begin{minipage}[b]{\linewidth}\raggedright
Setting
\end{minipage} & \begin{minipage}[b]{\linewidth}\raggedright
Total objective (BRL tn)
\end{minipage} & \begin{minipage}[b]{\linewidth}\raggedright
Total gap (\%)
\end{minipage} & \begin{minipage}[b]{\linewidth}\raggedright
Economic cost (BRL bn)
\end{minipage} & \begin{minipage}[b]{\linewidth}\raggedright
Economic gap (\%)
\end{minipage} & \begin{minipage}[b]{\linewidth}\raggedright
Integers
\end{minipage} & \begin{minipage}[b]{\linewidth}\raggedright
Nonzero integers (\%)
\end{minipage} & \begin{minipage}[b]{\linewidth}\raggedright
Cont. vars
\end{minipage} & \begin{minipage}[b]{\linewidth}\raggedright
Nonzero continuous (\%)
\end{minipage} & \begin{minipage}[b]{\linewidth}\raggedright
Notes
\end{minipage} \\
\midrule\noalign{}
\endhead
\bottomrule\noalign{}
\endlastfoot
STO W / Native / 1T & 9.289109 & 0.004119 & 385.502 & ND & 392 & 76.53 &
25,107,152 & 1.256 & Full vector \\
STO W / Native / 2T & 9.289109 & 0.004119 & 385.502 & ND & 392 & 76.53 &
25,107,152 & 1.256 & Full vector \\
STO W / Native / 4T & 9.289109 & 0.004119 & 385.502 & ND & 392 & 76.53 &
25,107,152 & 1.256 & Full vector \\
STO W / Native / 8T & 9.289109 & 0.004119 & 385.502 & ND & 392 & 76.53 &
25,107,152 & 1.256 & Full vector \\
STO W / Native / 16T & 9.289109 & 0.004119 & 385.502 & ND & 392 & 76.53
& 25,107,152 & 1.256 & Full vector \\
STO W / Pyomo / 1T r1 & 9.289109 & 0.004105 & 385.502 & ND & 392 & 76.53
& 25,107,152 & 1.249 & Full vector · 2231228 · LP 29.0 min \\
STO W / Pyomo / 1T r2 & 9.289109 & 0.004105 & 385.502 & ND & 392 & ND &
25,107,152 & ND & Vector ND · 2232366 · LP 37.3 min \\
STO W / Pyomo / 1T r3 & 9.289109 & 0.004105 & 385.502 & ND & 392 & ND &
25,107,152 & ND & Vector ND · 2232373 · LP 28.9 min \\
STO W / Pyomo / 2T r1 & 9.289109 & 0.004122 & 385.502 & ND & 392 & ND &
25,107,152 & ND & Vector ND · 2232355 · LP 46.7 min \\
STO W / Pyomo / 2T r2 & 9.289109 & 0.004122 & 385.502 & ND & 392 & ND &
25,107,152 & ND & Vector ND · 2232359 · LP 48.7 min \\
STO W / Pyomo / 2T r3 & 9.289109 & 0.004122 & 385.502 & ND & 392 & ND &
25,107,152 & ND & Vector ND · 2232375 · LP 36.3 min \\
STO W / Pyomo / 4T r1 & 9.289109 & 0.004120 & 385.502 & ND & 392 & ND &
25,107,152 & ND & Vector ND · 2232353 · LP 45.8 min \\
STO W / Pyomo / 4T r2 & 9.289109 & 0.004120 & 385.502 & ND & 392 & ND &
25,107,152 & ND & Vector ND · 2232360 · LP 55.3 min \\
STO W / Pyomo / 4T r3 & 9.289109 & 0.004120 & 385.502 & ND & 392 & ND &
25,107,152 & ND & Vector ND · 2232376 · LP 39.5 min \\
STO W / Pyomo / 8T r1 & 9.289109 & 0.004120 & 385.502 & ND & 392 & ND &
25,107,152 & ND & Vector ND · 2232356 · LP 71.1 min \\
STO W / Pyomo / 8T r2 & 9.289109 & 0.004120 & 385.502 & ND & 392 & ND &
25,107,152 & ND & Vector ND · 2232365 · LP 70.7 min \\
STO W / Pyomo / 8T r3 & 9.289109 & 0.004120 & 385.502 & ND & 392 & ND &
25,107,152 & ND & Vector ND · 2232369 · LP 71.2 min \\
STO W / Pyomo / 16T r1 & 9.289109 & 0.004120 & 385.502 & ND & 392 & ND &
25,107,152 & ND & Vector ND · 2231593 · LP 54.3 min \\
STO W / Pyomo / 16T r2 & 9.289109 & 0.004120 & 385.502 & ND & 392 & ND &
25,107,152 & ND & Vector ND · 2232363 · LP 78.7 min \\
STO W / Pyomo / 16T r3 & 9.289109 & 0.004120 & 385.502 & ND & 392 & ND &
25,107,152 & ND & Vector ND · 2232370 · LP 77.0 min \\
\end{longtable}
\endgroup
}

\subsubsection{h400}\label{h400}

\textbf{Model and network}

{\def\LTcaptype{none} % do not increment counter
\begingroup\footnotesize
\begin{longtable}[]{@{}
  >{\raggedright\arraybackslash}p{(\linewidth - 14\tabcolsep) * \real{0.1250}}
  >{\raggedright\arraybackslash}p{(\linewidth - 14\tabcolsep) * \real{0.1250}}
  >{\raggedright\arraybackslash}p{(\linewidth - 14\tabcolsep) * \real{0.1250}}
  >{\raggedright\arraybackslash}p{(\linewidth - 14\tabcolsep) * \real{0.1250}}
  >{\raggedright\arraybackslash}p{(\linewidth - 14\tabcolsep) * \real{0.1250}}
  >{\raggedright\arraybackslash}p{(\linewidth - 14\tabcolsep) * \real{0.1250}}
  >{\raggedright\arraybackslash}p{(\linewidth - 14\tabcolsep) * \real{0.1250}}
  >{\raggedright\arraybackslash}p{(\linewidth - 14\tabcolsep) * \real{0.1250}}@{}}
\toprule\noalign{}
\begin{minipage}[b]{\linewidth}\raggedright
Setting
\end{minipage} & \begin{minipage}[b]{\linewidth}\raggedright
Rows
\end{minipage} & \begin{minipage}[b]{\linewidth}\raggedright
Columns
\end{minipage} & \begin{minipage}[b]{\linewidth}\raggedright
Nonzeros
\end{minipage} & \begin{minipage}[b]{\linewidth}\raggedright
Network nodes
\end{minipage} & \begin{minipage}[b]{\linewidth}\raggedright
Arcs
\end{minipage} & \begin{minipage}[b]{\linewidth}\raggedright
Used arcs (\%)
\end{minipage} & \begin{minipage}[b]{\linewidth}\raggedright
B\&B nodes
\end{minipage} \\
\midrule\noalign{}
\endhead
\bottomrule\noalign{}
\endlastfoot
STO W / Native / 1T & 1,160,742 & 42,106,944 & 160,697,126 & 474 &
76,320 & 8.05 & 1 \\
STO W / Pyomo / 1T E1 & 1,435,062 & 42,106,944 & 161,245,766 & 474 &
76,320 & ND & 1 \\
STO W / Pyomo / 1T PR148 & 1,435,062 & 42,106,944 & 161,245,766 & 474 &
76,320 & ND & 1 \\
STO W / Pyomo / 2T E2 & 1,435,062 & 42,106,944 & 161,245,766 & 474 &
76,320 & ND & 1 \\
STO W / Pyomo / 4T E2 & 1,435,062 & 42,106,944 & 161,245,766 & 474 &
76,320 & ND & 1 \\
STO W / Pyomo / 8T E2 & 1,435,062 & 42,106,944 & 161,245,766 & 474 &
76,320 & ND & 1 \\
STO W / Pyomo / 16T E1 & 1,435,062 & 42,106,944 & 161,245,766 & 474 &
76,320 & ND & 1 \\
STO W / Pyomo / 16T PR148 & 1,435,062 & 42,106,944 & 161,245,766 & 474 &
76,320 & ND & 1 \\
\end{longtable}
\endgroup
}

\textbf{Incumbent, certificate and variable usage}

{\def\LTcaptype{none} % do not increment counter
\begingroup\footnotesize
\begin{longtable}[]{@{}
  >{\raggedright\arraybackslash}p{(\linewidth - 18\tabcolsep) * \real{0.1000}}
  >{\raggedright\arraybackslash}p{(\linewidth - 18\tabcolsep) * \real{0.1000}}
  >{\raggedright\arraybackslash}p{(\linewidth - 18\tabcolsep) * \real{0.1000}}
  >{\raggedright\arraybackslash}p{(\linewidth - 18\tabcolsep) * \real{0.1000}}
  >{\raggedright\arraybackslash}p{(\linewidth - 18\tabcolsep) * \real{0.1000}}
  >{\raggedright\arraybackslash}p{(\linewidth - 18\tabcolsep) * \real{0.1000}}
  >{\raggedright\arraybackslash}p{(\linewidth - 18\tabcolsep) * \real{0.1000}}
  >{\raggedright\arraybackslash}p{(\linewidth - 18\tabcolsep) * \real{0.1000}}
  >{\raggedright\arraybackslash}p{(\linewidth - 18\tabcolsep) * \real{0.1000}}
  >{\raggedright\arraybackslash}p{(\linewidth - 18\tabcolsep) * \real{0.1000}}@{}}
\toprule\noalign{}
\begin{minipage}[b]{\linewidth}\raggedright
Setting
\end{minipage} & \begin{minipage}[b]{\linewidth}\raggedright
Total objective (BRL tn)
\end{minipage} & \begin{minipage}[b]{\linewidth}\raggedright
Total gap (\%)
\end{minipage} & \begin{minipage}[b]{\linewidth}\raggedright
Economic cost (BRL bn)
\end{minipage} & \begin{minipage}[b]{\linewidth}\raggedright
Economic gap (\%)
\end{minipage} & \begin{minipage}[b]{\linewidth}\raggedright
Integers
\end{minipage} & \begin{minipage}[b]{\linewidth}\raggedright
Nonzero integers (\%)
\end{minipage} & \begin{minipage}[b]{\linewidth}\raggedright
Cont. vars
\end{minipage} & \begin{minipage}[b]{\linewidth}\raggedright
Nonzero continuous (\%)
\end{minipage} & \begin{minipage}[b]{\linewidth}\raggedright
Notes
\end{minipage} \\
\midrule\noalign{}
\endhead
\bottomrule\noalign{}
\endlastfoot
STO W / Native / 1T & 7.593395 & 0.836784 & 381.573 & ND & 492 & 81.30 &
42,106,452 & 0.875 & Full vector \\
STO W / Pyomo / 1T E1 & 7.593395 & 0.836784 & 381.573 & ND & 492 & ND &
42,106,452 & ND & Vector ND · 2247718 \\
STO W / Pyomo / 1T PR148 & 7.593395 & 0.836784 & 381.573 & ND & 492 & ND
& 42,106,452 & ND & Vector ND · 2232379 \\
STO W / Pyomo / 2T E2 & 7.593395 & 0.836784 & 381.573 & ND & 492 & ND &
42,106,452 & ND & Vector ND · 2250119 \\
STO W / Pyomo / 4T E2 & 7.593395 & 0.836784 & 381.573 & ND & 492 & ND &
42,106,452 & ND & Vector ND · 2250123 \\
STO W / Pyomo / 8T E2 & 7.593395 & 0.836784 & 381.573 & ND & 492 & ND &
42,106,452 & ND & Vector ND · 2250127 \\
STO W / Pyomo / 16T E1 & 7.593395 & 0.836784 & 381.573 & ND & 492 & ND &
42,106,452 & ND & Vector ND · 2247727 \\
STO W / Pyomo / 16T PR148 & 7.593395 & 0.836784 & 381.573 & ND & 492 &
ND & 42,106,452 & ND & Vector ND · 2232382 \\
\end{longtable}
\endgroup
}

\subsubsection{h500}\label{h500}

\textbf{Model and network}

{\def\LTcaptype{none} % do not increment counter
\begingroup\footnotesize
\begin{longtable}[]{@{}
  >{\raggedright\arraybackslash}p{(\linewidth - 14\tabcolsep) * \real{0.1250}}
  >{\raggedright\arraybackslash}p{(\linewidth - 14\tabcolsep) * \real{0.1250}}
  >{\raggedright\arraybackslash}p{(\linewidth - 14\tabcolsep) * \real{0.1250}}
  >{\raggedright\arraybackslash}p{(\linewidth - 14\tabcolsep) * \real{0.1250}}
  >{\raggedright\arraybackslash}p{(\linewidth - 14\tabcolsep) * \real{0.1250}}
  >{\raggedright\arraybackslash}p{(\linewidth - 14\tabcolsep) * \real{0.1250}}
  >{\raggedright\arraybackslash}p{(\linewidth - 14\tabcolsep) * \real{0.1250}}
  >{\raggedright\arraybackslash}p{(\linewidth - 14\tabcolsep) * \real{0.1250}}@{}}
\toprule\noalign{}
\begin{minipage}[b]{\linewidth}\raggedright
Setting
\end{minipage} & \begin{minipage}[b]{\linewidth}\raggedright
Rows
\end{minipage} & \begin{minipage}[b]{\linewidth}\raggedright
Columns
\end{minipage} & \begin{minipage}[b]{\linewidth}\raggedright
Nonzeros
\end{minipage} & \begin{minipage}[b]{\linewidth}\raggedright
Network nodes
\end{minipage} & \begin{minipage}[b]{\linewidth}\raggedright
Arcs
\end{minipage} & \begin{minipage}[b]{\linewidth}\raggedright
Used arcs (\%)
\end{minipage} & \begin{minipage}[b]{\linewidth}\raggedright
B\&B nodes
\end{minipage} \\
\midrule\noalign{}
\endhead
\bottomrule\noalign{}
\endlastfoot
STO W / Native / 1T & 1,430,842 & 63,426,344 & 244,039,126 & 574 &
115,400 & 6.09 & 1 \\
\end{longtable}
\endgroup
}

\textbf{Incumbent, certificate and variable usage}

{\def\LTcaptype{none} % do not increment counter
\begingroup\footnotesize
\begin{longtable}[]{@{}
  >{\raggedright\arraybackslash}p{(\linewidth - 18\tabcolsep) * \real{0.1000}}
  >{\raggedright\arraybackslash}p{(\linewidth - 18\tabcolsep) * \real{0.1000}}
  >{\raggedright\arraybackslash}p{(\linewidth - 18\tabcolsep) * \real{0.1000}}
  >{\raggedright\arraybackslash}p{(\linewidth - 18\tabcolsep) * \real{0.1000}}
  >{\raggedright\arraybackslash}p{(\linewidth - 18\tabcolsep) * \real{0.1000}}
  >{\raggedright\arraybackslash}p{(\linewidth - 18\tabcolsep) * \real{0.1000}}
  >{\raggedright\arraybackslash}p{(\linewidth - 18\tabcolsep) * \real{0.1000}}
  >{\raggedright\arraybackslash}p{(\linewidth - 18\tabcolsep) * \real{0.1000}}
  >{\raggedright\arraybackslash}p{(\linewidth - 18\tabcolsep) * \real{0.1000}}
  >{\raggedright\arraybackslash}p{(\linewidth - 18\tabcolsep) * \real{0.1000}}@{}}
\toprule\noalign{}
\begin{minipage}[b]{\linewidth}\raggedright
Setting
\end{minipage} & \begin{minipage}[b]{\linewidth}\raggedright
Total objective (BRL tn)
\end{minipage} & \begin{minipage}[b]{\linewidth}\raggedright
Total gap (\%)
\end{minipage} & \begin{minipage}[b]{\linewidth}\raggedright
Economic cost (BRL bn)
\end{minipage} & \begin{minipage}[b]{\linewidth}\raggedright
Economic gap (\%)
\end{minipage} & \begin{minipage}[b]{\linewidth}\raggedright
Integers
\end{minipage} & \begin{minipage}[b]{\linewidth}\raggedright
Nonzero integers (\%)
\end{minipage} & \begin{minipage}[b]{\linewidth}\raggedright
Cont. vars
\end{minipage} & \begin{minipage}[b]{\linewidth}\raggedright
Nonzero continuous (\%)
\end{minipage} & \begin{minipage}[b]{\linewidth}\raggedright
Notes
\end{minipage} \\
\midrule\noalign{}
\endhead
\bottomrule\noalign{}
\endlastfoot
STO W / Native / 1T & 6.297240 & 0.207432 & 373.173 & ND & 592 & 84.46 &
63,425,752 & 0.630 & Full vector \\
\end{longtable}
\endgroup
}

\newpage{}

\section{Appendix D Data and Preparation of the Experimental
Instances}\label{appendix-d-data-and-preparation-of-the-experimental-instances}

\subsection{D1 Data sources and scope}\label{d1-data-sources-and-scope}

The four tested instances (h215, h300, h400, h500), detailed in Table D.5, share the same supply, demand, model parameters, scenario inputs and freight,
storage and investment cost tables. The 146 existing warehouses and the
69 h215 candidates are retained across all four workbooks; h300, h400
and h500 contain 154, 254 and 354 candidates, respectively. The
identifier \texttt{CDA} distinguishes a warehouse in the registry and
its processed workbook. The selection record provides the ordered
membership for that population. The historical sources below document
how the inputs were constructed.

An experimental instance combines a facility population, supply and
demand over time, investment and operating costs, capacity rules, road
distances and an uncertainty design. This appendix describes these
inputs and the transformations needed to reproduce the evaluated model.

Data is imported from a spreadsheet. It supplies the model
schema, supply and demand records, costs and a 215-facility population.
Additional candidate locations come from
another spreadsheet. Spreadsheet versions and structures
are identified in the
\href{https://github.com/unb-lamfo-or-ai-hpc/model-agrologistic/blob/1b54023f4e0c0c27aba64883b1e9cf25817ad0f8/data/manifests/warehouse_population_source_v020.json}{registry
manifest}. The two sources have different roles: the former workbook defines the reference instance; the latter registry provides potential additions to its candidate set.

Appendix F details \href{https://doi.org/10.5281/zenodo.22751909}{Zenodo version 0.2.0
dataset}, which provides the deposited source material, subject to the stated file-level access and reuse conditions. Full file identities, worksheet inventories and experimental input paths are listed in the
\href{https://github.com/unb-lamfo-or-ai-hpc/model-agrologistic/blob/1b54023f4e0c0c27aba64883b1e9cf25817ad0f8/docs/evidence/appendix_input_reproduction_20261028.json}{input
reproduction guide}. Readers should compare those identities before using a downloaded workbook; a matching filename or facility count alone does not identify an input.

The data spreadsheet contains 14 worksheets and 215 facility records: 146 existing warehouses and 69 candidates. The Warehouses spreadsheet has 18,086 facility records: 146
existing sites and 17,940 candidates. Candidate with nonpositive static capacity are excluded from the eligible selection pool.

The effective model contains 37 origin identifiers, 27 domestic destination roles, ten export roles, two products and 60 planning periods. An identifier denotes a model entity; coincident coordinates do not merge entities. A destination may have both domestic and export roles, so the ten export records should not be interpreted as ten
distinct physical ports.

PAREI A ANÁLISE AQUI

\subsection{D2 Quantities and units}\label{d2-quantities-and-units}

Supply and demand are read from Oferta and Demanda. Each record
identifies a product, location, period and quantity. Facility records
use CDA as the identifier; names and coordinates are descriptive
attributes. The input reader aggregates additive supply or demand
entries that share the same effective keys. Missing mandatory
identifiers are treated as input errors. Initial inventory is zero where
no initial-stock record is supplied.

\begin{longtable}[]{@{}
  >{\raggedright\arraybackslash}p{(\linewidth - 4\tabcolsep) * \real{0.3333}}
  >{\raggedright\arraybackslash}p{(\linewidth - 4\tabcolsep) * \real{0.3333}}
  >{\raggedright\arraybackslash}p{(\linewidth - 4\tabcolsep) * \real{0.3333}}@{}}
\caption{Table D.1. Units used in the experimental
model.}\tabularnewline
\toprule\noalign{}
\begin{minipage}[b]{\linewidth}\raggedright
Quantity
\end{minipage} & \begin{minipage}[b]{\linewidth}\raggedright
Meaning
\end{minipage} & \begin{minipage}[b]{\linewidth}\raggedright
Effective unit
\end{minipage} \\
\midrule\noalign{}
\endfirsthead
\toprule\noalign{}
\begin{minipage}[b]{\linewidth}\raggedright
Quantity
\end{minipage} & \begin{minipage}[b]{\linewidth}\raggedright
Meaning
\end{minipage} & \begin{minipage}[b]{\linewidth}\raggedright
Effective unit
\end{minipage} \\
\midrule\noalign{}
\endhead
\bottomrule\noalign{}
\endlastfoot
Supply and domestic demand & Grain available or required within a
product-period record & t within the period \\
Export limit & Upper bound on optional arrivals within a period & t
within the period \\
Initial inventory and static capacity & Grain held at a warehouse & t \\
Reception and shipping rates & Grain that can be handled each operating
day & t/day \\
Freight coefficient & Cost multiplied by grain quantity and directed
road distance & BRL/(t km) \\
Inventory holding & Cost per unit held in a planning period & BRL/(t
period) \\
Capacity investment & Fixed cost and cost per investment-tonne,
according to facility eligibility & BRL and BRL/t \\
\end{longtable}

Static capacity limits inventory, whereas reception and shipping limit
flows within a period. The experiments use 30 operating days per period.
Candidate and expansion investments add 0.20 t/day of reception and
shipping capacity per investment-tonne, corresponding to a coupling
coefficient of 0.20 day\(^{-1}\). Bulk-handling conversion uses 1.0
day\(^{-1}\) on its declared investment-tonne basis and adds handling
capacity without adding static storage.

Two actual gold-workbook records illustrate these distinctions. Existing
facility CDA 35.0287.0002-6 has static capacity 890 t, reception 5 t/day
and shipping 6 t/day. Its period reception and shipping limits are
therefore 150 t and 180 t before eligible investment effects. Candidate
CDA 79.G590.0003-9 has zero base capacity in all three categories and an
investment upper bound of 307079 t. That bound is neither installed
capacity nor an investment chosen by the optimizer. These values
illustrate input interpretation, not observed grain movements.

The
\href{https://github.com/unb-lamfo-or-ai-hpc/model-agrologistic/blob/1b54023f4e0c0c27aba64883b1e9cf25817ad0f8/src/logic/excel_loader.py}{Excel
reader} and
\href{https://github.com/unb-lamfo-or-ai-hpc/model-agrologistic/blob/1b54023f4e0c0c27aba64883b1e9cf25817ad0f8/src/logic/mathematical_contract.py}{mathematical
input contract} implement the unit conventions and transformations.
Facility-specific eligibility and maxima take precedence over general
fallback values.

\subsection{D3 Investment, operating costs and penalty
coefficients}\label{d3-investment-operating-costs-and-penalty-coefficients}

Table D.2 reports the investment cost schedule used when a facility has
no explicit effective coefficient. The average of each low/high pair
provides the corresponding fallback. An explicit facility coefficient is
retained instead of replacing every facility by a common average.

\begin{longtable}[]{@{}
  >{\raggedright\arraybackslash}p{(\linewidth - 4\tabcolsep) * \real{0.2727}}
  >{\raggedleft\arraybackslash}p{(\linewidth - 4\tabcolsep) * \real{0.3636}}
  >{\raggedleft\arraybackslash}p{(\linewidth - 4\tabcolsep) * \real{0.3636}}@{}}
\caption{Table D.2. Investment cost schedule and arithmetic-mean
fallback coefficients.}\tabularnewline
\toprule\noalign{}
\begin{minipage}[b]{\linewidth}\raggedright
Investment type
\end{minipage} & \begin{minipage}[b]{\linewidth}\raggedleft
Low/high source values, BRL/t
\end{minipage} & \begin{minipage}[b]{\linewidth}\raggedleft
Fallback coefficient, BRL/t
\end{minipage} \\
\midrule\noalign{}
\endfirsthead
\toprule\noalign{}
\begin{minipage}[b]{\linewidth}\raggedright
Investment type
\end{minipage} & \begin{minipage}[b]{\linewidth}\raggedleft
Low/high source values, BRL/t
\end{minipage} & \begin{minipage}[b]{\linewidth}\raggedleft
Fallback coefficient, BRL/t
\end{minipage} \\
\midrule\noalign{}
\endhead
\bottomrule\noalign{}
\endlastfoot
New conventional warehouse & 1000 / 1500 & 1250 \\
New metal silo & 785 / 1310 & 1047.5 \\
Existing-site expansion & 600 / 1200 & 900 \\
Bulk-handling conversion & 380 / 660 & 520 \\
\end{longtable}

The candidate-cost fallback uses the conventional-warehouse value of
1250 BRL/t. Where a candidate has a total investment estimate and a
positive maximum capacity, the capacity-proportional coefficient is
obtained by dividing that estimate by the maximum capacity, with the
corresponding fixed component set to zero. The same project estimate is
consequently not charged again as an additional fixed cost. The general
expansion limit is 0.25 of static capacity when no facility-specific
maximum is supplied.

Inventory-holding tariffs are 45 BRL/(t period) for soybean and 48
BRL/(t period) for corn, for both public and private operators. The
additional Outros tariff in the source does not create a third product
in the evaluated instances.

\begin{longtable}[]{@{}lr@{}}
\caption{Table D.3. Actual Frete!A2:B28 values in the registered gold
workbook.}\tabularnewline
\toprule\noalign{}
State code(s) & Source freight coefficient, BRL/(t km) \\
\midrule\noalign{}
\endfirsthead
\toprule\noalign{}
State code(s) & Source freight coefficient, BRL/(t km) \\
\midrule\noalign{}
\endhead
\bottomrule\noalign{}
\endlastfoot
MT & 0.1842 \\
AM & 0.1969 \\
PE & 0.2004 \\
MS & 0.2174 \\
BA & 0.2584 \\
AC, AL, AP, CE, ES, PA, PB, RJ, RN, RR, SE & 0.2648 \\
DF & 0.294 \\
MA & 0.2943 \\
PI & 0.2979 \\
PR & 0.312 \\
SP & 0.3279 \\
RO & 0.3346 \\
RS & 0.3408 \\
GO & 0.3682 \\
TO & 0.3826 \\
MG & 0.3971 \\
SC & 0.4014 \\
\end{longtable}

Transport cost is calculated from grain quantity, the applicable state
freight coefficient and directed road distance. These coefficients are a
declared cost approximation, rather than observed prices for each road
movement.

The scalar objective uses three artificial coefficients: 90000 BRL per
tonne of unmet domestic demand, 45000 BRL per tonne of static exceedance
in each period, and 45000 BRL per tonne of reception overflow. They are
identical across the reported scalar experiments. Static exceedance
summed across periods is expressed in tonne-periods; reception overflow
is already a period flow and is not multiplied by operating days again.

These weights allow the optimizer to trade investment and operation
against specified shortages or exceedances. They are not calibrated
economic damages. Financial expenditure and artificial penalties are
therefore reported separately. A small gap for their sum does not
establish the same relative accuracy for the financial subtotal or a
particular domestic service level.

All monetary results are in Brazilian reais over the 60-period horizon,
without a newly imposed currency conversion, price-year adjustment,
annualization or discount rate. The
\href{https://github.com/unb-lamfo-or-ai-hpc/model-agrologistic/blob/1b54023f4e0c0c27aba64883b1e9cf25817ad0f8/docs/evidence/appendix_input_reproduction_20261028.json}{parameter
guide} retains the complete numerical mapping to the effective model.

\subsection{D4 Warehouse populations}\label{d4-warehouse-populations}

The h215 reference population contains 146 existing facilities and 69
candidate anchors. Larger instances retain these facilities and add
candidates from a deterministic, proportionally stratified ranking of
the eligible registry. Strata use state, facility type and
static-capacity quartile. The selection seed and ordered memberships are
retained with the
\href{https://github.com/unb-lamfo-or-ai-hpc/model-agrologistic/blob/1b54023f4e0c0c27aba64883b1e9cf25817ad0f8/docs/evidence/scalar_h300_population_lineage_20261020.json}{population
lineage}.

\begin{longtable}[]{@{}
  >{\raggedright\arraybackslash}p{(\linewidth - 6\tabcolsep) * \real{0.2000}}
  >{\raggedleft\arraybackslash}p{(\linewidth - 6\tabcolsep) * \real{0.2667}}
  >{\raggedleft\arraybackslash}p{(\linewidth - 6\tabcolsep) * \real{0.2667}}
  >{\raggedleft\arraybackslash}p{(\linewidth - 6\tabcolsep) * \real{0.2667}}@{}}
\caption{Table D.4. Facility populations available to the
optimization.}\tabularnewline
\toprule\noalign{}
\begin{minipage}[b]{\linewidth}\raggedright
Instance
\end{minipage} & \begin{minipage}[b]{\linewidth}\raggedleft
Existing facilities
\end{minipage} & \begin{minipage}[b]{\linewidth}\raggedleft
Candidate options
\end{minipage} & \begin{minipage}[b]{\linewidth}\raggedleft
Candidates added beyond h215
\end{minipage} \\
\midrule\noalign{}
\endfirsthead
\toprule\noalign{}
\begin{minipage}[b]{\linewidth}\raggedright
Instance
\end{minipage} & \begin{minipage}[b]{\linewidth}\raggedleft
Existing facilities
\end{minipage} & \begin{minipage}[b]{\linewidth}\raggedleft
Candidate options
\end{minipage} & \begin{minipage}[b]{\linewidth}\raggedleft
Candidates added beyond h215
\end{minipage} \\
\midrule\noalign{}
\endhead
\bottomrule\noalign{}
\endlastfoot
h215 & 146 & 69 & 0 \\
h300 & 146 & 154 & 85 \\
h400 & 146 & 254 & 185 \\
h500 & 146 & 354 & 285 \\
\end{longtable}

The
\href{https://github.com/unb-lamfo-or-ai-hpc/model-agrologistic/blob/1b54023f4e0c0c27aba64883b1e9cf25817ad0f8/src/logic/warehouse_population.py}{population
selector} orders the eligible candidates. The
\href{https://github.com/unb-lamfo-or-ai-hpc/model-agrologistic/blob/1b54023f4e0c0c27aba64883b1e9cf25817ad0f8/src/logic/warehouse_workbook.py}{workbook
constructor} retains the canonical rows and materializes the selected
additions. For each new candidate, base static, reception and shipping
capacities are zero. Its positive registry static capacity becomes the
maximum investable capacity. The new candidate has no expansion or
bulk-conversion eligibility. Thus, adding a registry location creates an
investment option, not an already operating warehouse.

The facility sets are nested, but the optimizer's chosen capacities and
flows need not be. In addition, route screening is recalculated for each
population; nested facility membership does not guarantee nested
selected arc sets or feasible regions. Objective and runtime changes
therefore cannot be attributed to warehouse count alone.

\subsection{D5 Road distances, network selection and
scenarios}\label{d5-road-distances-network-selection-and-scenarios}

The retained audits bind the Brazil OpenStreetMap extract
\texttt{brazil-260901} to its PBF SHA-256
\texttt{95c432ac3219338bbc7d4f5f7d20a5de709348ea928f656efd34609a51319fde}.
The
\href{https://github.com/unb-lamfo-or-ai-hpc/model-agrologistic/blob/965ca52777a5899163af72c9630e591b82c61de8/scripts/materialize_policy_osrm_population.slurm}{NPAD
materialization workflow} records Singularity 3.7.1, OSRM 6.0.0, the OCI
and SIF image digests and the driving-profile digest. It extracts,
partitions and customizes the fixed road network for multilevel Dijkstra
routing. OSRM Table responses give directed distances in metres, which
are converted to kilometres for the workbook. The ordered coordinate
pair, dataset and profile identify a cache entry. The historical
workbooks can reuse those entries directly. The
\href{https://github.com/unb-lamfo-or-ai-hpc/model-agrologistic/blob/965ca52777a5899163af72c9630e591b82c61de8/docs/evidence/osrm_release_preparation_20261101.json}{materialization
audit index} binds the original dataset and each frozen workbook.

OSRM materialization supplies directed road distances for the fixed
Brazil road dataset and driving profile. The retained h215, h300, h400
and h500 workbooks contain respectively 63289, 113269, 190569 and 287869
directed distance pairs. No Haversine fallback is recorded for these
materializations. The processed workbooks add a Distancias sheet to the
gold structure, giving 15 worksheets.

The distance matrix precedes product-specific network selection. The
evaluated sparse policy uses 20\% connectivity-preserving Pareto
selection and one export sink per product. Materialized pairs are
consequently not counts of retained arcs or flow variables. Topology W
excludes direct origin-to-destination movement; D includes it in
addition to warehouse routes. Availability of a distance does not itself
authorize a flow variable.

STO combines supply multipliers 0.85, 1 and 1.15 with domestic-demand
multipliers 0.95, 1 and 1.05. Their Cartesian product gives nine
complete trajectories, each with probability 1/9. These are controlled
scenarios, rather than empirically estimated event probabilities. DET
uses the original deterministic supply and demand arrays. Investments
are shared across the STO trajectories; operations vary by scenario.

Road lengths determine transport costs but do not supply route
polylines. Geographic figures therefore show the available verified node
and investment locations. Straight lines between nodes are not
substituted for actual road paths. Transport quantities can still be
analyzed from exported flows without claiming a verified geographic
path.

\subsection{D6 Reproducing an
instance}\label{d6-reproducing-an-instance}

First obtain the source files identified in the
\href{https://github.com/unb-lamfo-or-ai-hpc/model-agrologistic/blob/1b54023f4e0c0c27aba64883b1e9cf25817ad0f8/data/manifests/mvp_data_contract.json}{data
contract} and input guide, subject to their access conditions. Compare
full file hashes before processing. The runtime workbook, facility
order, road materialization and experiment definition must belong to the
same instance.

The deposited input archive contains several workbook generations. For
the road-distance instances discussed here, select the four members in
Table D.5. The common filename model\_input.xlsx does not distinguish a
pre-routing workbook from a frozen OSRM input. In the table, P denotes
data/processed/policy\_population\_v020\_osrm and N denotes
data/processed/nine\_population\_osrm\_v1; both paths are relative to
the archive root.

\begin{longtable}[]{@{}lll@{}}
\caption{Table D.5. Exact OSRM workbook members in the version 0.2.0
input archive.}\tabularnewline
\toprule\noalign{}
Instance & Archive member, using the prefixes above & SHA-256 prefix \\
\midrule\noalign{}
\endfirsthead
\toprule\noalign{}
Instance & Archive member, using the prefixes above & SHA-256 prefix \\
\midrule\noalign{}
\endhead
\bottomrule\noalign{}
\endlastfoot
h215 & P/warehouses\_215/model\_input.xlsx & \texttt{40edb66585da} \\
h300 & N/warehouses\_300/model\_input.xlsx & \texttt{7761cce77cf9} \\
h400 & N/warehouses\_400/model\_input.xlsx & \texttt{c30854569977} \\
h500 & N/warehouses\_500/model\_input.xlsx & \texttt{888fbff42e4f} \\
\end{longtable}

The 12-character prefixes help locate a record; verification requires
the complete 64-character SHA-256 in the
\href{https://github.com/unb-lamfo-or-ai-hpc/model-agrologistic/blob/1b54023f4e0c0c27aba64883b1e9cf25817ad0f8/docs/evidence/appendix_input_reproduction_20261028.json}{input
reproduction guide}, which also records worksheet counts and ordered
warehouse identities.

Each workbook directory contains osrm\_materialization\_audit.json.
Reuse it with the common warehouse\_population\_manifest.json and
warehouse\_population\_order.csv under
data/processed/warehouse\_population\_v020. The candidate registry is
the separate archive member
external-inputs/Warehouses\_Existing\_Candidate\_All.xlsx; it supplies
the selection pool, rather than replacing a processed workbook at
execution time.

The SHA-256 for the complete agrologistic-v020-inputs.zip archive is

\texttt{583b04afb49c4cba13c9d0aaa18131e12a18bac1da410cd801616bd0b639811f}

Its archive hash and each selected member hash serve different checks:
the former identifies the package, while the latter binds a calculation
to its effective input. These processed workbooks are archive members,
rather than files distributed as GitHub source. A newly generated
workbook should be recorded separately; changes in spreadsheet metadata
or recalculated caches can alter its byte identity.

Next verify existing/candidate roles, ordered CDA memberships and the 69
canonical anchors. Reuse the recorded candidate order and OSRM
materialization. A new random draw or routing query creates a different
input. Apply the recorded capacity coupling, terminal-stock rule, route
policy, topology and scenario design. Check the effective coefficients
listed above before constructing the model.

Finally, bind the input to the source version and numerical settings of
its experiment. The original
\href{https://github.com/unb-lamfo-or-ai-hpc/model-agrologistic/blob/1b54023f4e0c0c27aba64883b1e9cf25817ad0f8/docs/evidence/scalar_h215_2217990_review.json}{h215},
\href{https://github.com/unb-lamfo-or-ai-hpc/model-agrologistic/blob/1b54023f4e0c0c27aba64883b1e9cf25817ad0f8/docs/evidence/scalar_h300_2218193_review.json}{h300},
\href{https://github.com/unb-lamfo-or-ai-hpc/model-agrologistic/blob/1b54023f4e0c0c27aba64883b1e9cf25817ad0f8/docs/evidence/scalar_h400_2219884_review.json}{h400}
and
\href{https://github.com/unb-lamfo-or-ai-hpc/model-agrologistic/blob/1b54023f4e0c0c27aba64883b1e9cf25817ad0f8/docs/evidence/scalar_h500_2220382_review.json}{h500}
records identify those versions, inputs and validation results. The
\href{https://github.com/unb-lamfo-or-ai-hpc/model-agrologistic/blob/1b54023f4e0c0c27aba64883b1e9cf25817ad0f8/docs/evidence/scalar_wd_h215_input_inspection_20261022.json}{W/D
comparison} documents the common graph and costs for the topology
comparison. Appendix C explains small verification cases and
original-unit solution validation before large results are interpreted.

\subsection{D7 Reproduction and evidence
limits}\label{d7-reproduction-and-evidence-limits}

Public source code, the gold workbook, data inventories and experiment
records allow inspection of the formulation, effective coefficients and
reported claims. Repeating a large calculation additionally requires
access to its exact processed inputs, a compatible solver and runtime,
sufficient resources and the declared stopping rules.

Rebuilding the article reproduces its document and derived tables; it
does not rerun the optimization. Reviewing an existing certificate
checks retained evidence. A fresh numerical reproduction instead
requires newly exported decisions and independent residual validation.
Alternative optima, solver versions and shared-node conditions can
produce different decision vectors or logs even when objective quality
and physical feasibility agree.

Access to a deposited archive and permission to redistribute its
components are separate questions. Project-authored source code,
external facility records and OpenStreetMap-derived distances retain
their respective terms. Readers should retain the accompanying source
notices when obtaining or reusing the workbooks; the code's MIT license
does not relicense every worksheet or upstream dataset.

Traceability makes these distinctions inspectable. It does not by itself
establish economic calibration, unrestricted data access or
reproducibility under arbitrary computing conditions.

\subsection{D8 Dimensional interpretation of the calculation and its
outputs}\label{d8-dimensional-interpretation-of-the-calculation-and-its-outputs}

A planning period is one step of the input time grid. The experiments
use 60 such steps and 30 operating days per step. Supply, demand and
flow entries are quantities of grain within a step, not daily rates. End
stock and static capacity are quantities held at an instant. Reception
and shipping inputs, by contrast, are rates in t/day. Multiplying those
rates by operating days produces the period-tonne limits used in the
flow inequalities.

For example, an initial reception rate of 5 t/day permits 150 t in a
30-day period. Installing 10 t of capacity with a handling gain of 0.20
day\(^{-1}\) adds 2 t/day, hence 60 t of period reception. The same
conversion is applied to shipping. Bulk-handling investment changes
these handling rates but does not add static storage. A change from 30
to another number of operating days would change the experiment and must
be recorded.

Holding tariffs are quoted per tonne per modeled period. The objective
uses the charge for one period of stock at each time index; the
one-period duration is implicit in its coefficient. Thus, holding 2 t
for each of two periods at 45 BRL/(t period) costs 180 BRL. Multiplying
again by 30 operating days would apply an extra, unsupported time
factor. The stock-exceedance penalty follows the same per-period
assessment, whereas shortage and reception overflow are quantities of
grain within their periods.

\begin{longtable}[]{@{}
  >{\raggedright\arraybackslash}p{(\linewidth - 4\tabcolsep) * \real{0.3333}}
  >{\raggedright\arraybackslash}p{(\linewidth - 4\tabcolsep) * \real{0.3333}}
  >{\raggedright\arraybackslash}p{(\linewidth - 4\tabcolsep) * \real{0.3333}}@{}}
\caption{Table D.6. Dimensions and aggregation scopes of the main
outputs.}\tabularnewline
\toprule\noalign{}
\begin{minipage}[b]{\linewidth}\raggedright
Reported quantity
\end{minipage} & \begin{minipage}[b]{\linewidth}\raggedright
Aggregation and unit
\end{minipage} & \begin{minipage}[b]{\linewidth}\raggedright
Interpretation
\end{minipage} \\
\midrule\noalign{}
\endfirsthead
\toprule\noalign{}
\begin{minipage}[b]{\linewidth}\raggedright
Reported quantity
\end{minipage} & \begin{minipage}[b]{\linewidth}\raggedright
Aggregation and unit
\end{minipage} & \begin{minipage}[b]{\linewidth}\raggedright
Interpretation
\end{minipage} \\
\midrule\noalign{}
\endhead
\bottomrule\noalign{}
\endlastfoot
Added candidate or expansion capacity & First-stage investment tonnes &
Installed storage; charged once across scenarios \\
Unmet demand & Probability-weighted sum over products, destinations and
periods, t & Expected unmet horizon quantity \\
Stock exceedance & Probability-weighted stock excess summed across
periods, t period & Accumulated exposure; not installed capacity \\
Reception overflow & Probability-weighted excess arrivals summed across
periods, t & Period flow; a daily equivalent divides each record by its
operating days \\
Financial subtotal and artificial penalties & Each reconstructed in BRL,
then reported separately & Their sum is the scalar objective; physical
tonnes are not added to monetary costs \\
Relative objective gap & Dimensionless fraction; multiplied by 100 when
shown as \% & A statement about the penalized objective, not a service
percentage \\
\end{longtable}

Four tonnes of stock excess in each of two periods represent 8 t period;
four tonnes of reception overflow in each period represent 8 t. These
quantities must not be added and interpreted as 16 t of required
installed capacity. The historical generic emergency-capacity output
field therefore does not define a combined physical investment
requirement.

Scenario weights average alternative operating trajectories; they do not
multiply warehouse capacity or create nine successive physical
shipments. Domestic service is the ratio of probability-weighted served
tonnes to probability-weighted demand tonnes over the horizon. It is not
generally the unweighted mean of scenario service percentages. Likewise,
adding all transport legs counts transshipped grain more than once;
endpoint deliveries and individual movement families have separate
meanings (Appendix E).

OSRM's service returns metres and seconds. Materialization divides road
length by 1000 once to populate Distancia\_km; Duracao\_s remains in
seconds. The MILP uses distance for freight and does not impose
transit-time constraints from the duration column. A freight rate in
BRL/(t km), multiplied by kilometres and transported tonnes, gives BRL.
The road length follows the fastest route under the frozen driving
profile, rather than necessarily the shortest route by distance, as
specified in the
\href{https://project-osrm.org/docs/v5.24.0/api/\#table-service}{OSRM
Table API}.

Timing tables distinguish seconds from minutes and native from
application elapsed time. Memory in GiB is bytes divided by \(2^{30}\);
a decimal GB is \(10^9\) bytes. Solver-managed memory, sampled
process-tree RSS, Slurm MaxRSS and requested memory retain their
separate scopes.

\subsection{D9 Locating source code, distance matrices and solution
data}\label{d9-locating-source-code-distance-matrices-and-solution-data}

The
\href{https://github.com/unb-lamfo-or-ai-hpc/model-agrologistic/blob/43acc42ffce0350cbce6fdc4b0632529a503547d/data/manifests/research_data_distribution_v1.json}{research-data
catalog} connects repository paths, deposited archive members and
prepared distance tables to full checksums. Its distribution and access
fields distinguish an existing file from a proposed release location.

GitHub contains the canonical gold schema, source manifests,
mathematical and solver code, small reviewed investment CSVs and the
metadata needed to identify an experiment. The full processed workbooks
and directed distance matrices belong to the versioned data archive.
Table D.5 identifies the original workbook members; the catalog also
identifies the corresponding seven-column distance tables, with endpoint
IDs, movement family, kilometres, seconds, provider and fallback reason.
The complete original worksheets remain the reference even when a CSV
copy is used for inspection.

Three objects must be distinguished when reproducing a plan: the full
road distance table, the product-specific selected admissible arcs, and
the positive flows returned by an optimizer. The distance table precedes
route screening. A zero or absent exported flow does not determine
admissibility. A reproduction therefore requires the recorded route
policy and selected graph identity in addition to the workbook and
warehouse order.

Full solution vectors, flow-distance companions, inventory records and
original validation outputs are organized separately from small article
tables. They retain configuration, product, period and scenario
identifiers and the source used for aggregation. Access and release
status are recorded per asset in the catalog; a prepared local file is
not presented as an already published Zenodo download.

\subsection{D10 Reading the exported numerical
tables}\label{d10-reading-the-exported-numerical-tables}

The tables separate shared infrastructure decisions from
scenario-specific operation. A warehouse decision has no scenario or
period index: it applies to the complete planning horizon. Flow, stock
and shortage rows instead retain the original scenario, period and
product identifiers. Their \texttt{value} fields are original
quantities, not probability-weighted contributions.

\begin{longtable}[]{@{}
  >{\raggedright\arraybackslash}p{(\linewidth - 4\tabcolsep) * \real{0.3333}}
  >{\raggedright\arraybackslash}p{(\linewidth - 4\tabcolsep) * \real{0.3333}}
  >{\raggedright\arraybackslash}p{(\linewidth - 4\tabcolsep) * \real{0.3333}}@{}}
\caption{Table D.7. Original output files and their physical
meaning.}\tabularnewline
\toprule\noalign{}
\begin{minipage}[b]{\linewidth}\raggedright
Original file
\end{minipage} & \begin{minipage}[b]{\linewidth}\raggedright
Principal coordinates
\end{minipage} & \begin{minipage}[b]{\linewidth}\raggedright
Meaning of the quantities
\end{minipage} \\
\midrule\noalign{}
\endfirsthead
\toprule\noalign{}
\begin{minipage}[b]{\linewidth}\raggedright
Original file
\end{minipage} & \begin{minipage}[b]{\linewidth}\raggedright
Principal coordinates
\end{minipage} & \begin{minipage}[b]{\linewidth}\raggedright
Meaning of the quantities
\end{minipage} \\
\midrule\noalign{}
\endhead
\bottomrule\noalign{}
\endlastfoot
warehouse\_decisions.csv & Warehouse & Activation flags and installed
candidate/expansion tonnes; bulk-handling basis is separate \\
flows.csv & Scenario, movement family, endpoints, product, period &
Tonnes moved on one directed leg within that period \\
inventories.csv & Scenario, warehouse, product, period & Tonnes held at
period end \\
unmet\_demand.csv & Scenario, domestic destination, product, period &
Tonnes of unmet period demand \\
emergency\_capacity.csv & Scenario, warehouse, period, capacity\_type &
Static stock excess or reception-flow overflow, distinguished by type \\
emergency\_capacity\_daily.csv & Same coordinates plus days\_in\_period
& Reception overflow divided by operating days; the static
daily-equivalent field is blank \\
\end{longtable}

The
\href{https://github.com/unb-lamfo-or-ai-hpc/model-agrologistic/blob/8f33a360d8a1ca3f18c5c7a3aa0eac2cd2ad5a90/docs/evidence/paired_h300_release_dictionary_20261101.json}{source-bound
file and column dictionary} records member paths, byte sizes, SHA-256
values, column names and row counts for the Gurobi h300 pilot. The
separate
\href{https://github.com/unb-lamfo-or-ai-hpc/model-agrologistic/blob/083ef96cc99a346204df22551ce0739c10e8a2f7/docs/evidence/paired_h300_highs_release_dictionary_20261101.json}{HiGHS
dictionary} binds the companion output. The
\href{https://github.com/unb-lamfo-or-ai-hpc/model-agrologistic/blob/c493be8cd127c5175fc7eb111400229e64b8e3c0/src/logic/experiment_runner.py}{export
implementation} defines those fields. Its expected-horizon summaries
multiply each operational row by its scenario probability exactly once
and sum over the intended coordinates. First-stage investments are not
counted once per scenario.

For example, static records contribute to stock exposure in
tonne-periods when summed over time. Reception records contribute tonnes
of overflow, or a daily equivalent when each record is divided by its
own operating days. Neither is automatically an engineering capacity
requirement. Blank static daily equivalents do not mean that stock
excess was zero. Similarly, the sparse exported tables do not list every
feasible zero-flow arc; Appendix E explains the distinction between
admissibility and a returned positive flow.

For the first matched h300 pair, a data-only reconstruction from the
exact workbook and executed policy provides 45,240 product-specific
admissible arcs. This is the total number of directed endpoint-product
records across both products, before expansion over operating periods
and scenarios. The counts are 4,440 origin-to-warehouse, 4,800
warehouse-to-destination and 36,000 warehouse-to-warehouse arcs, with no
direct or repair arcs. Its graph and mathematical-contract fingerprints
match both original input receipts. The
\href{https://github.com/unb-lamfo-or-ai-hpc/model-agrologistic/blob/0f63c10eb1abf88318a4d9c6afc75fcac97f6266/docs/evidence/h300_selected_graph_release_20261101.json}{selected-graph
receipt} and
\href{https://github.com/unb-lamfo-or-ai-hpc/model-agrologistic/blob/0f63c10eb1abf88318a4d9c6afc75fcac97f6266/data/manifests/zenodo_next_release_inventory_20261101.json}{file-level
release inventory} distinguish these locally prepared files from the
existing published deposit. They identify the graph separately from the
two positive-flow companions; they do not claim that a new Zenodo
download is already available.

\subsection{D11 Data access and versioned
evidence}\label{d11-data-access-and-versioned-evidence}

Reproducing the population comparison requires the ordered warehouse
list and the effective workbook for each case. The shared 60-month input
horizon defines the period totals reported in the article. The nine STO
scenario probabilities weight those totals to give expected values. In
the DET case the single scenario has probability one. National annual
production and installed storage statistics provide external context and
have their own reporting dates and units.

Zenodo 0.2.0 (Celestino 2026) is the public release cited in the
article. It contains the historical inputs, population records, road
audits and inventories listed in the
\href{https://doi.org/10.5281/zenodo.22751909}{deposit record}. Complete
derivative distance tables, admissible graphs and solution companions
have separate source identities. Their access status and redistribution
conditions are recorded in the
\href{https://github.com/unb-lamfo-or-ai-hpc/model-agrologistic/blob/2c1ad568bbc8eb0ee45b42ad59b5a3a442e0ab3c/data/manifests/zenodo_release_plan_20261102.json}{repository
inventory}. The derivative files remain in preparation, with source
notices and private-field checks before public deposit.

Table D.8 distinguishes repository metadata from the versioned numerical
records they describe.

\begin{longtable}[]{@{}
  >{\raggedright\arraybackslash}p{(\linewidth - 4\tabcolsep) * \real{0.3333}}
  >{\raggedright\arraybackslash}p{(\linewidth - 4\tabcolsep) * \real{0.3333}}
  >{\raggedright\arraybackslash}p{(\linewidth - 4\tabcolsep) * \real{0.3333}}@{}}
\caption{Table D.8. Complementary repository and versioned-data delivery
roles.}\tabularnewline
\toprule\noalign{}
\begin{minipage}[b]{\linewidth}\raggedright
Data object
\end{minipage} & \begin{minipage}[b]{\linewidth}\raggedright
GitHub role
\end{minipage} & \begin{minipage}[b]{\linewidth}\raggedright
Versioned numerical archive role
\end{minipage} \\
\midrule\noalign{}
\endfirsthead
\toprule\noalign{}
\begin{minipage}[b]{\linewidth}\raggedright
Data object
\end{minipage} & \begin{minipage}[b]{\linewidth}\raggedright
GitHub role
\end{minipage} & \begin{minipage}[b]{\linewidth}\raggedright
Versioned numerical archive role
\end{minipage} \\
\midrule\noalign{}
\endhead
\bottomrule\noalign{}
\endlastfoot
Model and experimental protocol & Source, tests, settings and exact
software versions & Link each dataset version to the executed source and
environment \\
Inputs and complete directed distances & Gold schema, source members,
materialization audits and checksums & Preserve original workbooks and
provide full endpoint/family distance tables \\
Selected admissible network & Screening policy, counts and
graph/contract identities & Export the product-specific admissibility
mask separately from complete distances \\
Accepted investment and operating decisions & Small lists and reviewed
figures/tables & Full coordinate tables and SQLite companions with
period/scenario/product identities \\
Computational outcome & Configuration/status tables and evidence
references & Reviewed native/validation/timing/resource records,
including censored and unsuccessful outcomes \\
\end{longtable}

Source-specific terms govern sharing of external inputs and their
derivatives. The software license applies to the project code. Published
data, review summaries and experimental originals retain their own
identifiers and access conditions.

\newpage{}

\section{Appendix E. Geographic decisions and origin-destination
flows}\label{appendix-e.-geographic-decisions-and-origin-destination-flows}

\subsection{E1 Reading the geographic
results}\label{e1-reading-the-geographic-results}

The maps in this appendix locate the investments in each accepted
planning solution. An orange circle denotes a candidate warehouse that
is opened, a red square denotes an expansion, and a magenta diamond
denotes conversion to bulk handling. Purple stars identify
\textbf{export ports}. A warehouse can receive more than one investment,
and different facilities can share coordinates. Their CDA identifiers
remain separate in the investment tables even when symbols overlap.

An opened candidate, installed capacity and operational use answer
different questions. Opening is a binary investment decision; installed
capacity is a quantity in tonnes; use is established by a positive
exported flow. For an existing warehouse, the original availability flag
does not signify a new opening. Likewise, the bulk-handling investment
basis is not additional static storage. The tables distinguish candidate
storage, static expansion and the basis used to calculate handling
gains.

The configuration summaries count positive investment quantities above
\(10^{-7}\) t. Binary activation flags and the original quantities are
retained separately in the complete lists, allowing activation without
material capacity to be distinguished from a positive investment. This
display threshold does not replace the original feasibility tolerances.

The individual maps share the national extent, IBGE boundaries and
Albers equal-area projection. Map lines delineate study areas and do not
necessarily depict accepted national boundaries. Capacity views use
circle area proportional to added candidate and expansion storage, in
tonnes. OSRM road distances were calculated and stored in the processed
workbooks listed in Appendix D. The flow-distance tables associate those
values with the exported flows. The maps show facility positions and
investments and omit route lines for legibility.

The configurations comprise the deterministic and stochastic
warehouse-only plans, the larger stochastic populations, the two h215
plans permitting direct origin-destination movements, and the accepted
Pyomo solutions. The historical native-Gurobi h300 runs with 1, 2, 4, 8
and 16 solver threads have byte-identical original warehouse, flow,
inventory and shortage exports. Their common investment map and flow
matrices are therefore represented by the single-thread vector; their
runtime observations remain distinct. The accepted Pyomo Gurobi and
HiGHS solutions retain their own exported vectors, attained gaps and
investment summaries. The two common-Pyomo h300 pilots are separate
observations. Their own maps and investment lists follow the historical
configurations; their full flow-distance companions retain separate
source and solution identities.

To compare returned investment plans, join the two lists on the original
facility ID and retain a separate column for each investment type and
its quantity. The first matched h300 pair opens 154 candidates in each
solution, and the opening sets coincide. The positive expansion sets
contain 70 sites for Gurobi and 69 for HiGHS, with 69 shared sites. The
positive bulk-handling sets contain 76 and 77 sites, respectively, with
76 shared sites. These counts use the stated \(10^{-7}\) t threshold and
identify common sites, with capacities retained separately in the lists.
The repeated timing grid in Appendix C describes solver observations.
The maps here describe the particular accepted vectors linked to each
figure. The h400 E1--E2 follow-up contributes computational receipts in
Appendix C. The h400 map and flow summaries here retain their original
native Gurobi vector. Equal reported objective values in the follow-up
grid do not establish identical facility-level decisions.

\subsection{E2 From exported flows to sparse
matrices}\label{e2-from-exported-flows-to-sparse-matrices}

Let \(x^{k}_{ijpts}\) denote the exported tonnes on directed movement
family \(k\), from endpoint \(i\) to endpoint \(j\), for product \(p\),
period \(t\) and scenario \(s\). The full-resolution companion retains
these indices and the original probability \(\pi_s\). Its
expected-horizon matrix is

\[
F^{k,p}_{ij}=\sum_{s\in S}\pi_s\sum_{t\in T}x^{k}_{ijpts}.
\]

For a deterministic solution there is one scenario with probability one.
For the reported stochastic configurations, nine Cartesian supply-demand
scenarios each have probability \(1/9\). The aggregation sums
transported tonnes over the original 60 monthly periods. It neither
annualizes the quantities nor treats different scenarios as successive
physical shipments.

The paired distance matrix is \(D^k_{ij}=d^k_{ij}\), where \(d^k_{ij}\)
is the OSRM road distance in kilometres in the exact processed workbook
used by that configuration. Distances are joined by \textbf{movement
family and directed endpoint identities}, never by the order of rows or
by geographic proximity. A destination's domestic or export role is
retained separately. Reversing a movement requires its own recorded
distance. Distances are not summed when periods or scenarios are
aggregated.

The implementation's legacy family codes are OD for origin-warehouse, DD
for warehouse-warehouse, DC for warehouse-destination, and OC for a
direct origin-destination movement. Thus, the code OD does not mean a
direct origin-customer movement. OC occurs only in the D configurations.
The warehouse-only W configurations do not admit that movement family.

The companions use coordinate-form sparse matrices: each row names the
product, directed endpoints, movement family, expected tonnes and OSRM
kilometres. The unaggregated companion also gives period, scenario,
probability and original tonnes. This representation preserves the
positive flows without printing large mostly empty rectangular arrays.
Summing all movement families would double-count grain transported
through successive network legs; domestic and export arrivals must
instead be evaluated at their appropriate destinations.

The original exporter omits flow magnitudes at or below \(10^{-7}\) t.
An absent row therefore must not be treated as proof of an exact zero or
as evidence that an arc was inadmissible. Zero flow, an arc excluded by
screening, an unexported small value and missing road geometry are
distinct states. Reconstructing the complete admissibility mask requires
the frozen network definition identified with the input data. All
positive records in the prepared companions have an exact matching
workbook distance; no distance or road polyline is interpolated.

These are \textbf{flows of accepted solutions}, rather than claims of
exact zero-gap optima. Each configuration below retains its attained
relative objective gap. The original independent primal certificates
were accepted; post-processing their exports does not constitute a new
optimization or a fresh residual check.

The
\href{https://github.com/unb-lamfo-or-ai-hpc/model-agrologistic/blob/0788de130253dddc790b179c8831d8ee29dabcbd/manuscript/supplement/appendix_e_sources.json}{configuration
and original-export index},
\href{https://github.com/unb-lamfo-or-ai-hpc/model-agrologistic/blob/0788de130253dddc790b179c8831d8ee29dabcbd/manuscript/supplement/appendix_e_summary.json}{numerical
companion inventory} and
\href{https://github.com/unb-lamfo-or-ai-hpc/model-agrologistic/blob/0788de130253dddc790b179c8831d8ee29dabcbd/scripts/build_appendix_e.py}{reconstruction
program} identify the precise versions and aggregation procedure. The
complete warehouse lists are linked under each configuration.
Full-resolution flow-distance tables are prepared as compressed CSVs and
a separate SQLite database; they are not currently public GitHub
downloads. Access to their processed-source data and redistribution
conditions follows the availability statement in Appendix D. The routing
pair cache supports workbook preparation. The geographic SQLite binds
coordinates and decisions for the maps. The separate numerical SQLite
organizes sparse flows and directed distances from the accepted exports
and frozen workbook. These stores serve different stages and carry
separate source identities.

The versioned delivery sequence in Appendix D11 applies to these
companions: the present GitHub indices identify their contents and
original sources, while the planned Zenodo version will provide reviewed
full-resolution files. They are currently prepared locally, not newly
deposited. Original coordinate rows and scenario weights remain
available in the companions rather than being replaced by a rendered
map. Later campaign records without a usable incumbent have no
investment map or flow matrix to contribute.

\subsection{E3 h215 DET/W, Gurobi, native gurobipy, 4
threads}\label{e3-h215-detw-gurobi-native-gurobipy-4-threads}

\begin{figure}[H]

{\centering \pandocbounded{\includegraphics[keepaspectratio]{scalar_h215_det.pdf}}

}

\caption{Figure E.1. Investments for h215 DET/W, Gurobi, native
gurobipy, 4 threads.}

\end{figure}%

The accepted solution has a relative objective gap of
\textbf{0.000812\%}. It opens 69 candidates, expands 80 existing sites
and converts 66 sites for bulk handling.

The
\href{https://github.com/unb-lamfo-or-ai-hpc/model-agrologistic/blob/0788de130253dddc790b179c8831d8ee29dabcbd/manuscript/supplement/investments/scalar_h215_det.csv}{complete
warehouse investment list} provides CDA, municipality, UF, role,
activation flags, candidate capacity, expansion and bulk-handling basis
in tonnes. Existing-site availability is excluded from candidate
openings.

{\def\LTcaptype{none} % do not increment counter
\begin{longtable}[]{@{}
  >{\raggedright\arraybackslash}p{(\linewidth - 4\tabcolsep) * \real{0.3333}}
  >{\raggedleft\arraybackslash}p{(\linewidth - 4\tabcolsep) * \real{0.3333}}
  >{\raggedleft\arraybackslash}p{(\linewidth - 4\tabcolsep) * \real{0.3333}}@{}}
\toprule\noalign{}
\begin{minipage}[b]{\linewidth}\raggedright
Movement
\end{minipage} & \begin{minipage}[b]{\linewidth}\raggedleft
Positive exported records
\end{minipage} & \begin{minipage}[b]{\linewidth}\raggedleft
Expected arc-tonnes over the horizon
\end{minipage} \\
\midrule\noalign{}
\endhead
\bottomrule\noalign{}
\endlastfoot
origin-warehouse (OD) & 7,773 & 1,374,420,353.085 \\
warehouse-warehouse (DD) & 3,339 & 167,041,802.967 \\
warehouse-destination (DC) & 10,731 & 1,371,457,378.297 \\
origin-destination (OC) & 0 & 0.000 \\
\end{longtable}
}

These are network-arc volumes. Adding different movement families would
count a consignment at successive legs; their sum is not delivered
demand. The full-resolution companion retains 21,843 flow records, while
the expected-horizon sparse matrix contains 3,015 positive product/arc
coordinates.

The following five largest product/arc entries illustrate the paired
flow-distance matrices. They are a selected view, not the complete
matrix.

{\def\LTcaptype{none} % do not increment counter
\begin{longtable}[]{@{}
  >{\raggedright\arraybackslash}p{(\linewidth - 10\tabcolsep) * \real{0.1667}}
  >{\raggedright\arraybackslash}p{(\linewidth - 10\tabcolsep) * \real{0.1667}}
  >{\raggedright\arraybackslash}p{(\linewidth - 10\tabcolsep) * \real{0.1667}}
  >{\raggedright\arraybackslash}p{(\linewidth - 10\tabcolsep) * \real{0.1667}}
  >{\raggedleft\arraybackslash}p{(\linewidth - 10\tabcolsep) * \real{0.1667}}
  >{\raggedleft\arraybackslash}p{(\linewidth - 10\tabcolsep) * \real{0.1667}}@{}}
\toprule\noalign{}
\begin{minipage}[b]{\linewidth}\raggedright
Family
\end{minipage} & \begin{minipage}[b]{\linewidth}\raggedright
From
\end{minipage} & \begin{minipage}[b]{\linewidth}\raggedright
To / role
\end{minipage} & \begin{minipage}[b]{\linewidth}\raggedright
Product
\end{minipage} & \begin{minipage}[b]{\linewidth}\raggedleft
Expected horizon t
\end{minipage} & \begin{minipage}[b]{\linewidth}\raggedleft
OSRM km
\end{minipage} \\
\midrule\noalign{}
\endhead
\bottomrule\noalign{}
\endlastfoot
OD & Toledo - PR & 66.0333.0108-8 & Milho & 45,119,038.044 & 279.7462 \\
OD & Sorriso - MT & 56.F194.0007-8 & Soja & 37,989,237.118 & 30.9958 \\
OD & Sorriso - MT & 56.F194.0007-8 & Milho & 36,955,955.290 & 30.9958 \\
DC & 56.F194.0007-8 & Itaituba - PA / export & Milho & 36,649,167.943 &
1,084.3545 \\
DC & 56.F194.0007-8 & Itaituba - PA / export & Soja & 36,489,237.118 &
1,084.3545 \\
\end{longtable}
}

\begin{figure}[H]

{\centering \pandocbounded{\includegraphics[keepaspectratio]{scalar_h215_det_capacity.pdf}}

}

\caption{Figure E.2. Added static capacity for h215 DET/W, Gurobi,
native gurobipy, 4 threads.}

\end{figure}%

Circle area represents added candidate plus expansion storage in tonnes.
Bulk-handling basis is excluded from the circles; it determines handling
gains rather than additional static storage.

\subsection{E4 h215 STO/W, Gurobi, native gurobipy, 4
threads}\label{e4-h215-stow-gurobi-native-gurobipy-4-threads}

\begin{figure}[H]

{\centering \pandocbounded{\includegraphics[keepaspectratio]{scalar_h215_sto.pdf}}

}

\caption{Figure E.3. Investments for h215 STO/W, Gurobi, native
gurobipy, 4 threads.}

\end{figure}%

The accepted solution has a relative objective gap of
\textbf{0.740159\%}. It opens 69 candidates, expands 75 existing sites
and converts 71 sites for bulk handling.

The
\href{https://github.com/unb-lamfo-or-ai-hpc/model-agrologistic/blob/0788de130253dddc790b179c8831d8ee29dabcbd/manuscript/supplement/investments/scalar_h215_sto.csv}{complete
warehouse investment list} provides CDA, municipality, UF, role,
activation flags, candidate capacity, expansion and bulk-handling basis
in tonnes. Existing-site availability is excluded from candidate
openings.

{\def\LTcaptype{none} % do not increment counter
\begin{longtable}[]{@{}
  >{\raggedright\arraybackslash}p{(\linewidth - 4\tabcolsep) * \real{0.3333}}
  >{\raggedleft\arraybackslash}p{(\linewidth - 4\tabcolsep) * \real{0.3333}}
  >{\raggedleft\arraybackslash}p{(\linewidth - 4\tabcolsep) * \real{0.3333}}@{}}
\toprule\noalign{}
\begin{minipage}[b]{\linewidth}\raggedright
Movement
\end{minipage} & \begin{minipage}[b]{\linewidth}\raggedleft
Positive exported records
\end{minipage} & \begin{minipage}[b]{\linewidth}\raggedleft
Expected arc-tonnes over the horizon
\end{minipage} \\
\midrule\noalign{}
\endhead
\bottomrule\noalign{}
\endlastfoot
origin-warehouse (OD) & 68,943 & 1,374,420,353.085 \\
warehouse-warehouse (DD) & 30,950 & 165,214,630.554 \\
warehouse-destination (DC) & 94,595 & 1,371,554,638.383 \\
origin-destination (OC) & 0 & 0.000 \\
\end{longtable}
}

These are network-arc volumes. Adding different movement families would
count a consignment at successive legs; their sum is not delivered
demand. The full-resolution companion retains 194,488 flow records,
while the expected-horizon sparse matrix contains 4,090 positive
product/arc coordinates.

The following five largest product/arc entries illustrate the paired
flow-distance matrices. They are a selected view, not the complete
matrix.

{\def\LTcaptype{none} % do not increment counter
\begin{longtable}[]{@{}
  >{\raggedright\arraybackslash}p{(\linewidth - 10\tabcolsep) * \real{0.1667}}
  >{\raggedright\arraybackslash}p{(\linewidth - 10\tabcolsep) * \real{0.1667}}
  >{\raggedright\arraybackslash}p{(\linewidth - 10\tabcolsep) * \real{0.1667}}
  >{\raggedright\arraybackslash}p{(\linewidth - 10\tabcolsep) * \real{0.1667}}
  >{\raggedleft\arraybackslash}p{(\linewidth - 10\tabcolsep) * \real{0.1667}}
  >{\raggedleft\arraybackslash}p{(\linewidth - 10\tabcolsep) * \real{0.1667}}@{}}
\toprule\noalign{}
\begin{minipage}[b]{\linewidth}\raggedright
Family
\end{minipage} & \begin{minipage}[b]{\linewidth}\raggedright
From
\end{minipage} & \begin{minipage}[b]{\linewidth}\raggedright
To / role
\end{minipage} & \begin{minipage}[b]{\linewidth}\raggedright
Product
\end{minipage} & \begin{minipage}[b]{\linewidth}\raggedleft
Expected horizon t
\end{minipage} & \begin{minipage}[b]{\linewidth}\raggedleft
OSRM km
\end{minipage} \\
\midrule\noalign{}
\endhead
\bottomrule\noalign{}
\endlastfoot
OD & Toledo - PR & 66.0333.0108-8 & Milho & 44,337,337.208 & 279.7462 \\
OD & Sorriso - MT & 56.F194.0007-8 & Soja & 37,682,338.094 & 30.9958 \\
OD & Sorriso - MT & 56.F194.0007-8 & Milho & 36,820,331.500 & 30.9958 \\
DC & 56.F194.0007-8 & Itaituba - PA / export & Milho & 36,479,315.850 &
1,084.3545 \\
DC & 56.F194.0007-8 & Itaituba - PA / export & Soja & 36,249,004.760 &
1,084.3545 \\
\end{longtable}
}

\begin{figure}[H]

{\centering \pandocbounded{\includegraphics[keepaspectratio]{scalar_h215_sto_capacity.pdf}}

}

\caption{Figure E.4. Added static capacity for h215 STO/W, Gurobi,
native gurobipy, 4 threads.}

\end{figure}%

Circle area represents added candidate plus expansion storage in tonnes.
Bulk-handling basis is excluded from the circles; it determines handling
gains rather than additional static storage.

\subsection{E5 h300 STO/W, Gurobi, native gurobipy, 1
thread}\label{e5-h300-stow-gurobi-native-gurobipy-1-thread}

\begin{figure}[H]

{\centering \pandocbounded{\includegraphics[keepaspectratio]{scalar_h300_t1.pdf}}

}

\caption{Figure E.5. Investments for h300 STO/W, Gurobi, native
gurobipy, 1 thread.}

\end{figure}%

The accepted solution has a relative objective gap of
\textbf{0.004119\%}. It opens 154 candidates, expands 70 existing sites
and converts 76 sites for bulk handling.

The
\href{https://github.com/unb-lamfo-or-ai-hpc/model-agrologistic/blob/0788de130253dddc790b179c8831d8ee29dabcbd/manuscript/supplement/investments/scalar_h300_t1.csv}{complete
warehouse investment list} provides CDA, municipality, UF, role,
activation flags, candidate capacity, expansion and bulk-handling basis
in tonnes. Existing-site availability is excluded from candidate
openings.

{\def\LTcaptype{none} % do not increment counter
\begin{longtable}[]{@{}
  >{\raggedright\arraybackslash}p{(\linewidth - 4\tabcolsep) * \real{0.3333}}
  >{\raggedleft\arraybackslash}p{(\linewidth - 4\tabcolsep) * \real{0.3333}}
  >{\raggedleft\arraybackslash}p{(\linewidth - 4\tabcolsep) * \real{0.3333}}@{}}
\toprule\noalign{}
\begin{minipage}[b]{\linewidth}\raggedright
Movement
\end{minipage} & \begin{minipage}[b]{\linewidth}\raggedleft
Positive exported records
\end{minipage} & \begin{minipage}[b]{\linewidth}\raggedleft
Expected arc-tonnes over the horizon
\end{minipage} \\
\midrule\noalign{}
\endhead
\bottomrule\noalign{}
\endlastfoot
origin-warehouse (OD) & 84,572 & 1,374,420,353.085 \\
warehouse-warehouse (DD) & 33,468 & 189,900,398.314 \\
warehouse-destination (DC) & 109,474 & 1,371,084,300.466 \\
origin-destination (OC) & 0 & 0.000 \\
\end{longtable}
}

These are network-arc volumes. Adding different movement families would
count a consignment at successive legs; their sum is not delivered
demand. The full-resolution companion retains 227,514 flow records,
while the expected-horizon sparse matrix contains 5,058 positive
product/arc coordinates.

The following five largest product/arc entries illustrate the paired
flow-distance matrices. They are a selected view, not the complete
matrix.

{\def\LTcaptype{none} % do not increment counter
\begin{longtable}[]{@{}
  >{\raggedright\arraybackslash}p{(\linewidth - 10\tabcolsep) * \real{0.1667}}
  >{\raggedright\arraybackslash}p{(\linewidth - 10\tabcolsep) * \real{0.1667}}
  >{\raggedright\arraybackslash}p{(\linewidth - 10\tabcolsep) * \real{0.1667}}
  >{\raggedright\arraybackslash}p{(\linewidth - 10\tabcolsep) * \real{0.1667}}
  >{\raggedleft\arraybackslash}p{(\linewidth - 10\tabcolsep) * \real{0.1667}}
  >{\raggedleft\arraybackslash}p{(\linewidth - 10\tabcolsep) * \real{0.1667}}@{}}
\toprule\noalign{}
\begin{minipage}[b]{\linewidth}\raggedright
Family
\end{minipage} & \begin{minipage}[b]{\linewidth}\raggedright
From
\end{minipage} & \begin{minipage}[b]{\linewidth}\raggedright
To / role
\end{minipage} & \begin{minipage}[b]{\linewidth}\raggedright
Product
\end{minipage} & \begin{minipage}[b]{\linewidth}\raggedleft
Expected horizon t
\end{minipage} & \begin{minipage}[b]{\linewidth}\raggedleft
OSRM km
\end{minipage} \\
\midrule\noalign{}
\endhead
\bottomrule\noalign{}
\endlastfoot
OD & Toledo - PR & 66.0333.0108-8 & Milho & 39,287,786.033 & 279.7462 \\
OD & Sorriso - MT & 56.F194.0007-8 & Soja & 37,193,870.178 & 30.9958 \\
OD & Sorriso - MT & 56.F194.0007-8 & Milho & 36,544,449.063 & 30.9958 \\
DC & 56.F194.0007-8 & Itaituba - PA / export & Milho & 36,141,115.356 &
1,084.3545 \\
DC & 56.F194.0007-8 & Itaituba - PA / export & Soja & 35,693,870.178 &
1,084.3545 \\
\end{longtable}
}

\begin{figure}[H]

{\centering \pandocbounded{\includegraphics[keepaspectratio]{scalar_h300_t1_capacity.pdf}}

}

\caption{Figure E.6. Added static capacity for h300 STO/W, Gurobi,
native gurobipy, 1 thread.}

\end{figure}%

Circle area represents added candidate plus expansion storage in tonnes.
Bulk-handling basis is excluded from the circles; it determines handling
gains rather than additional static storage.

\subsection{E6 h400 STO/W, Gurobi, native gurobipy, 1
thread}\label{e6-h400-stow-gurobi-native-gurobipy-1-thread}

\begin{figure}[H]

{\centering \pandocbounded{\includegraphics[keepaspectratio]{scalar_h400_t1.pdf}}

}

\caption{Figure E.7. Investments for h400 STO/W, Gurobi, native
gurobipy, 1 thread.}

\end{figure}%

The accepted solution has a relative objective gap of
\textbf{0.836784\%}. It opens 254 candidates, expands 69 existing sites
and converts 77 sites for bulk handling.

The
\href{https://github.com/unb-lamfo-or-ai-hpc/model-agrologistic/blob/0788de130253dddc790b179c8831d8ee29dabcbd/manuscript/supplement/investments/scalar_h400_t1.csv}{complete
warehouse investment list} provides CDA, municipality, UF, role,
activation flags, candidate capacity, expansion and bulk-handling basis
in tonnes. Existing-site availability is excluded from candidate
openings.

{\def\LTcaptype{none} % do not increment counter
\begin{longtable}[]{@{}
  >{\raggedright\arraybackslash}p{(\linewidth - 4\tabcolsep) * \real{0.3333}}
  >{\raggedleft\arraybackslash}p{(\linewidth - 4\tabcolsep) * \real{0.3333}}
  >{\raggedleft\arraybackslash}p{(\linewidth - 4\tabcolsep) * \real{0.3333}}@{}}
\toprule\noalign{}
\begin{minipage}[b]{\linewidth}\raggedright
Movement
\end{minipage} & \begin{minipage}[b]{\linewidth}\raggedleft
Positive exported records
\end{minipage} & \begin{minipage}[b]{\linewidth}\raggedleft
Expected arc-tonnes over the horizon
\end{minipage} \\
\midrule\noalign{}
\endhead
\bottomrule\noalign{}
\endlastfoot
origin-warehouse (OD) & 100,664 & 1,374,420,353.085 \\
warehouse-warehouse (DD) & 40,205 & 200,859,544.376 \\
warehouse-destination (DC) & 125,724 & 1,370,916,024.951 \\
origin-destination (OC) & 0 & 0.000 \\
\end{longtable}
}

These are network-arc volumes. Adding different movement families would
count a consignment at successive legs; their sum is not delivered
demand. The full-resolution companion retains 266,593 flow records,
while the expected-horizon sparse matrix contains 6,143 positive
product/arc coordinates.

The following five largest product/arc entries illustrate the paired
flow-distance matrices. They are a selected view, not the complete
matrix.

{\def\LTcaptype{none} % do not increment counter
\begin{longtable}[]{@{}
  >{\raggedright\arraybackslash}p{(\linewidth - 10\tabcolsep) * \real{0.1667}}
  >{\raggedright\arraybackslash}p{(\linewidth - 10\tabcolsep) * \real{0.1667}}
  >{\raggedright\arraybackslash}p{(\linewidth - 10\tabcolsep) * \real{0.1667}}
  >{\raggedright\arraybackslash}p{(\linewidth - 10\tabcolsep) * \real{0.1667}}
  >{\raggedleft\arraybackslash}p{(\linewidth - 10\tabcolsep) * \real{0.1667}}
  >{\raggedleft\arraybackslash}p{(\linewidth - 10\tabcolsep) * \real{0.1667}}@{}}
\toprule\noalign{}
\begin{minipage}[b]{\linewidth}\raggedright
Family
\end{minipage} & \begin{minipage}[b]{\linewidth}\raggedright
From
\end{minipage} & \begin{minipage}[b]{\linewidth}\raggedright
To / role
\end{minipage} & \begin{minipage}[b]{\linewidth}\raggedright
Product
\end{minipage} & \begin{minipage}[b]{\linewidth}\raggedleft
Expected horizon t
\end{minipage} & \begin{minipage}[b]{\linewidth}\raggedleft
OSRM km
\end{minipage} \\
\midrule\noalign{}
\endhead
\bottomrule\noalign{}
\endlastfoot
OD & Toledo - PR & 66.0333.0108-8 & Milho & 37,524,219.236 & 279.7462 \\
OD & Sorriso - MT & 56.F194.0007-8 & Soja & 36,677,336.940 & 30.9958 \\
OD & Sorriso - MT & 56.F194.0007-8 & Milho & 36,383,126.177 & 30.9958 \\
DC & 56.F194.0007-8 & Itaituba - PA / export & Milho & 36,020,561.649 &
1,084.3545 \\
DC & 56.F194.0007-8 & Itaituba - PA / export & Soja & 35,213,796.112 &
1,084.3545 \\
\end{longtable}
}

\begin{figure}[H]

{\centering \pandocbounded{\includegraphics[keepaspectratio]{scalar_h400_t1_capacity.pdf}}

}

\caption{Figure E.8. Added static capacity for h400 STO/W, Gurobi,
native gurobipy, 1 thread.}

\end{figure}%

Circle area represents added candidate plus expansion storage in tonnes.
Bulk-handling basis is excluded from the circles; it determines handling
gains rather than additional static storage.

\subsection{E7 h500 STO/W, Gurobi, native gurobipy, 1
thread}\label{e7-h500-stow-gurobi-native-gurobipy-1-thread}

\begin{figure}[H]

{\centering \pandocbounded{\includegraphics[keepaspectratio]{scalar_h500_t1.pdf}}

}

\caption{Figure E.9. Investments for h500 STO/W, Gurobi, native
gurobipy, 1 thread.}

\end{figure}%

The accepted solution has a relative objective gap of
\textbf{0.207432\%}. It opens 354 candidates, expands 72 existing sites
and converts 74 sites for bulk handling.

The
\href{https://github.com/unb-lamfo-or-ai-hpc/model-agrologistic/blob/0788de130253dddc790b179c8831d8ee29dabcbd/manuscript/supplement/investments/scalar_h500_t1.csv}{complete
warehouse investment list} provides CDA, municipality, UF, role,
activation flags, candidate capacity, expansion and bulk-handling basis
in tonnes. Existing-site availability is excluded from candidate
openings.

{\def\LTcaptype{none} % do not increment counter
\begin{longtable}[]{@{}
  >{\raggedright\arraybackslash}p{(\linewidth - 4\tabcolsep) * \real{0.3333}}
  >{\raggedleft\arraybackslash}p{(\linewidth - 4\tabcolsep) * \real{0.3333}}
  >{\raggedleft\arraybackslash}p{(\linewidth - 4\tabcolsep) * \real{0.3333}}@{}}
\toprule\noalign{}
\begin{minipage}[b]{\linewidth}\raggedright
Movement
\end{minipage} & \begin{minipage}[b]{\linewidth}\raggedleft
Positive exported records
\end{minipage} & \begin{minipage}[b]{\linewidth}\raggedleft
Expected arc-tonnes over the horizon
\end{minipage} \\
\midrule\noalign{}
\endhead
\bottomrule\noalign{}
\endlastfoot
origin-warehouse (OD) & 110,385 & 1,374,420,353.085 \\
warehouse-warehouse (DD) & 44,294 & 209,001,024.443 \\
warehouse-destination (DC) & 135,106 & 1,370,709,301.045 \\
origin-destination (OC) & 0 & 0.000 \\
\end{longtable}
}

These are network-arc volumes. Adding different movement families would
count a consignment at successive legs; their sum is not delivered
demand. The full-resolution companion retains 289,785 flow records,
while the expected-horizon sparse matrix contains 7,023 positive
product/arc coordinates.

The following five largest product/arc entries illustrate the paired
flow-distance matrices. They are a selected view, not the complete
matrix.

{\def\LTcaptype{none} % do not increment counter
\begin{longtable}[]{@{}
  >{\raggedright\arraybackslash}p{(\linewidth - 10\tabcolsep) * \real{0.1667}}
  >{\raggedright\arraybackslash}p{(\linewidth - 10\tabcolsep) * \real{0.1667}}
  >{\raggedright\arraybackslash}p{(\linewidth - 10\tabcolsep) * \real{0.1667}}
  >{\raggedright\arraybackslash}p{(\linewidth - 10\tabcolsep) * \real{0.1667}}
  >{\raggedleft\arraybackslash}p{(\linewidth - 10\tabcolsep) * \real{0.1667}}
  >{\raggedleft\arraybackslash}p{(\linewidth - 10\tabcolsep) * \real{0.1667}}@{}}
\toprule\noalign{}
\begin{minipage}[b]{\linewidth}\raggedright
Family
\end{minipage} & \begin{minipage}[b]{\linewidth}\raggedright
From
\end{minipage} & \begin{minipage}[b]{\linewidth}\raggedright
To / role
\end{minipage} & \begin{minipage}[b]{\linewidth}\raggedright
Product
\end{minipage} & \begin{minipage}[b]{\linewidth}\raggedleft
Expected horizon t
\end{minipage} & \begin{minipage}[b]{\linewidth}\raggedleft
OSRM km
\end{minipage} \\
\midrule\noalign{}
\endhead
\bottomrule\noalign{}
\endlastfoot
OD & Toledo - PR & 66.0333.0108-8 & Milho & 38,378,996.796 & 279.7462 \\
OD & Sorriso - MT & 56.F194.0007-8 & Soja & 36,347,505.229 & 30.9958 \\
DC & 56.F194.0007-8 & Itaituba - PA / export & Soja & 34,905,524.715 &
1,084.3545 \\
OD & Sorriso - MT & 56.F194.0007-8 & Milho & 34,141,365.299 & 30.9958 \\
DC & 56.F194.0007-8 & Itaituba - PA / export & Milho & 33,715,764.420 &
1,084.3545 \\
\end{longtable}
}

\begin{figure}[H]

{\centering \pandocbounded{\includegraphics[keepaspectratio]{scalar_h500_t1_capacity.pdf}}

}

\caption{Figure E.10. Added static capacity for h500 STO/W, Gurobi,
native gurobipy, 1 thread.}

\end{figure}%

Circle area represents added candidate plus expansion storage in tonnes.
Bulk-handling basis is excluded from the circles; it determines handling
gains rather than additional static storage.

\subsection{E8 h215 DET/D, Gurobi, native gurobipy, 4
threads}\label{e8-h215-detd-gurobi-native-gurobipy-4-threads}

\begin{figure}[H]

{\centering \pandocbounded{\includegraphics[keepaspectratio]{scalar_h215_d_det.pdf}}

}

\caption{Figure E.11. Investments for h215 DET/D, Gurobi, native
gurobipy, 4 threads.}

\end{figure}%

The accepted solution has a relative objective gap of
\textbf{0.552427\%}. It opens 69 candidates, expands 142 existing sites
and converts 4 sites for bulk handling.

The
\href{https://github.com/unb-lamfo-or-ai-hpc/model-agrologistic/blob/0788de130253dddc790b179c8831d8ee29dabcbd/manuscript/supplement/investments/scalar_h215_d_det.csv}{complete
warehouse investment list} provides CDA, municipality, UF, role,
activation flags, candidate capacity, expansion and bulk-handling basis
in tonnes. Existing-site availability is excluded from candidate
openings.

{\def\LTcaptype{none} % do not increment counter
\begin{longtable}[]{@{}
  >{\raggedright\arraybackslash}p{(\linewidth - 4\tabcolsep) * \real{0.3333}}
  >{\raggedleft\arraybackslash}p{(\linewidth - 4\tabcolsep) * \real{0.3333}}
  >{\raggedleft\arraybackslash}p{(\linewidth - 4\tabcolsep) * \real{0.3333}}@{}}
\toprule\noalign{}
\begin{minipage}[b]{\linewidth}\raggedright
Movement
\end{minipage} & \begin{minipage}[b]{\linewidth}\raggedleft
Positive exported records
\end{minipage} & \begin{minipage}[b]{\linewidth}\raggedleft
Expected arc-tonnes over the horizon
\end{minipage} \\
\midrule\noalign{}
\endhead
\bottomrule\noalign{}
\endlastfoot
origin-warehouse (OD) & 4,080 & 246,625,348.240 \\
warehouse-warehouse (DD) & 4,654 & 105,637,499.209 \\
warehouse-destination (DC) & 5,393 & 244,772,691.525 \\
origin-destination (OC) & 2,586 & 1,127,795,004.845 \\
\end{longtable}
}

These are network-arc volumes. Adding different movement families would
count a consignment at successive legs; their sum is not delivered
demand. The full-resolution companion retains 16,713 flow records, while
the expected-horizon sparse matrix contains 2,914 positive product/arc
coordinates.

The following five largest product/arc entries illustrate the paired
flow-distance matrices. They are a selected view, not the complete
matrix.

{\def\LTcaptype{none} % do not increment counter
\begin{longtable}[]{@{}
  >{\raggedright\arraybackslash}p{(\linewidth - 10\tabcolsep) * \real{0.1667}}
  >{\raggedright\arraybackslash}p{(\linewidth - 10\tabcolsep) * \real{0.1667}}
  >{\raggedright\arraybackslash}p{(\linewidth - 10\tabcolsep) * \real{0.1667}}
  >{\raggedright\arraybackslash}p{(\linewidth - 10\tabcolsep) * \real{0.1667}}
  >{\raggedleft\arraybackslash}p{(\linewidth - 10\tabcolsep) * \real{0.1667}}
  >{\raggedleft\arraybackslash}p{(\linewidth - 10\tabcolsep) * \real{0.1667}}@{}}
\toprule\noalign{}
\begin{minipage}[b]{\linewidth}\raggedright
Family
\end{minipage} & \begin{minipage}[b]{\linewidth}\raggedright
From
\end{minipage} & \begin{minipage}[b]{\linewidth}\raggedright
To / role
\end{minipage} & \begin{minipage}[b]{\linewidth}\raggedright
Product
\end{minipage} & \begin{minipage}[b]{\linewidth}\raggedleft
Expected horizon t
\end{minipage} & \begin{minipage}[b]{\linewidth}\raggedleft
OSRM km
\end{minipage} \\
\midrule\noalign{}
\endhead
\bottomrule\noalign{}
\endlastfoot
OC & Sorriso - MT & Itaituba - PA / export & Soja & 203,616,292.657 &
1,092.5900 \\
OC & Sorriso - MT & Itaituba - PA / export & Milho & 131,290,493.804 &
1,092.5900 \\
OC & Tibagi - PR & Paranaguá - PR / export & Soja & 76,138,750.461 &
307.0004 \\
OC & Tupanciretã - RS & Rio Grande - RS / export & Soja & 59,395,881.979
& 437.8240 \\
OC & Rio Verde - GO & Santos - SP / export & Soja & 58,128,305.549 &
1,049.2323 \\
\end{longtable}
}

\begin{figure}[H]

{\centering \pandocbounded{\includegraphics[keepaspectratio]{scalar_h215_d_det_capacity.pdf}}

}

\caption{Figure E.12. Added static capacity for h215 DET/D, Gurobi,
native gurobipy, 4 threads.}

\end{figure}%

Circle area represents added candidate plus expansion storage in tonnes.
Bulk-handling basis is excluded from the circles; it determines handling
gains rather than additional static storage.

\subsection{E9 h215 STO/D, Gurobi, native gurobipy, 4
threads}\label{e9-h215-stod-gurobi-native-gurobipy-4-threads}

\begin{figure}[H]

{\centering \pandocbounded{\includegraphics[keepaspectratio]{scalar_h215_d_sto.pdf}}

}

\caption{Figure E.13. Investments for h215 STO/D, Gurobi, native
gurobipy, 4 threads.}

\end{figure}%

The accepted solution has a relative objective gap of
\textbf{0.623718\%}. It opens 69 candidates, expands 142 existing sites
and converts 4 sites for bulk handling.

The
\href{https://github.com/unb-lamfo-or-ai-hpc/model-agrologistic/blob/0788de130253dddc790b179c8831d8ee29dabcbd/manuscript/supplement/investments/scalar_h215_d_sto.csv}{complete
warehouse investment list} provides CDA, municipality, UF, role,
activation flags, candidate capacity, expansion and bulk-handling basis
in tonnes. Existing-site availability is excluded from candidate
openings.

{\def\LTcaptype{none} % do not increment counter
\begin{longtable}[]{@{}
  >{\raggedright\arraybackslash}p{(\linewidth - 4\tabcolsep) * \real{0.3333}}
  >{\raggedleft\arraybackslash}p{(\linewidth - 4\tabcolsep) * \real{0.3333}}
  >{\raggedleft\arraybackslash}p{(\linewidth - 4\tabcolsep) * \real{0.3333}}@{}}
\toprule\noalign{}
\begin{minipage}[b]{\linewidth}\raggedright
Movement
\end{minipage} & \begin{minipage}[b]{\linewidth}\raggedleft
Positive exported records
\end{minipage} & \begin{minipage}[b]{\linewidth}\raggedleft
Expected arc-tonnes over the horizon
\end{minipage} \\
\midrule\noalign{}
\endhead
\bottomrule\noalign{}
\endlastfoot
origin-warehouse (OD) & 36,698 & 251,548,979.450 \\
warehouse-warehouse (DD) & 42,249 & 92,859,992.548 \\
warehouse-destination (DC) & 49,237 & 249,712,101.193 \\
origin-destination (OC) & 23,143 & 1,122,871,373.635 \\
\end{longtable}
}

These are network-arc volumes. Adding different movement families would
count a consignment at successive legs; their sum is not delivered
demand. The full-resolution companion retains 151,327 flow records,
while the expected-horizon sparse matrix contains 3,933 positive
product/arc coordinates.

The following five largest product/arc entries illustrate the paired
flow-distance matrices. They are a selected view, not the complete
matrix.

{\def\LTcaptype{none} % do not increment counter
\begin{longtable}[]{@{}
  >{\raggedright\arraybackslash}p{(\linewidth - 10\tabcolsep) * \real{0.1667}}
  >{\raggedright\arraybackslash}p{(\linewidth - 10\tabcolsep) * \real{0.1667}}
  >{\raggedright\arraybackslash}p{(\linewidth - 10\tabcolsep) * \real{0.1667}}
  >{\raggedright\arraybackslash}p{(\linewidth - 10\tabcolsep) * \real{0.1667}}
  >{\raggedleft\arraybackslash}p{(\linewidth - 10\tabcolsep) * \real{0.1667}}
  >{\raggedleft\arraybackslash}p{(\linewidth - 10\tabcolsep) * \real{0.1667}}@{}}
\toprule\noalign{}
\begin{minipage}[b]{\linewidth}\raggedright
Family
\end{minipage} & \begin{minipage}[b]{\linewidth}\raggedright
From
\end{minipage} & \begin{minipage}[b]{\linewidth}\raggedright
To / role
\end{minipage} & \begin{minipage}[b]{\linewidth}\raggedright
Product
\end{minipage} & \begin{minipage}[b]{\linewidth}\raggedleft
Expected horizon t
\end{minipage} & \begin{minipage}[b]{\linewidth}\raggedleft
OSRM km
\end{minipage} \\
\midrule\noalign{}
\endhead
\bottomrule\noalign{}
\endlastfoot
OC & Sorriso - MT & Itaituba - PA / export & Soja & 202,957,786.378 &
1,092.5900 \\
OC & Sorriso - MT & Itaituba - PA / export & Milho & 128,091,894.532 &
1,092.5900 \\
OC & Tibagi - PR & Paranaguá - PR / export & Soja & 76,021,985.003 &
307.0004 \\
OC & Tupanciretã - RS & Rio Grande - RS / export & Soja & 59,459,489.217
& 437.8240 \\
OC & Rio Verde - GO & Santos - SP / export & Soja & 59,092,087.619 &
1,049.2323 \\
\end{longtable}
}

\begin{figure}[H]

{\centering \pandocbounded{\includegraphics[keepaspectratio]{scalar_h215_d_sto_capacity.pdf}}

}

\caption{Figure E.14. Added static capacity for h215 STO/D, Gurobi,
native gurobipy, 4 threads.}

\end{figure}%

Circle area represents added candidate plus expansion storage in tonnes.
Bulk-handling basis is excluded from the circles; it determines handling
gains rather than additional static storage.

\subsection{E10 h215 STO/W, Gurobi, Pyomo/APPSI, 4
threads}\label{e10-h215-stow-gurobi-pyomoappsi-4-threads}

\begin{figure}[H]

{\centering \pandocbounded{\includegraphics[keepaspectratio]{pyomo_gurobi_h215_w_native8h_license_v2.pdf}}

}

\caption{Figure E.15. Investments for h215 STO/W, Gurobi, Pyomo/APPSI, 4
threads.}

\end{figure}%

The accepted solution has a relative objective gap of
\textbf{0.740159\%}. It opens 69 candidates, expands 75 existing sites
and converts 71 sites for bulk handling.

The
\href{https://github.com/unb-lamfo-or-ai-hpc/model-agrologistic/blob/0788de130253dddc790b179c8831d8ee29dabcbd/manuscript/supplement/investments/pyomo_gurobi_h215_w_native8h_license_v2.csv}{complete
warehouse investment list} provides CDA, municipality, UF, role,
activation flags, candidate capacity, expansion and bulk-handling basis
in tonnes. Existing-site availability is excluded from candidate
openings.

{\def\LTcaptype{none} % do not increment counter
\begin{longtable}[]{@{}
  >{\raggedright\arraybackslash}p{(\linewidth - 4\tabcolsep) * \real{0.3333}}
  >{\raggedleft\arraybackslash}p{(\linewidth - 4\tabcolsep) * \real{0.3333}}
  >{\raggedleft\arraybackslash}p{(\linewidth - 4\tabcolsep) * \real{0.3333}}@{}}
\toprule\noalign{}
\begin{minipage}[b]{\linewidth}\raggedright
Movement
\end{minipage} & \begin{minipage}[b]{\linewidth}\raggedleft
Positive exported records
\end{minipage} & \begin{minipage}[b]{\linewidth}\raggedleft
Expected arc-tonnes over the horizon
\end{minipage} \\
\midrule\noalign{}
\endhead
\bottomrule\noalign{}
\endlastfoot
origin-warehouse (OD) & 68,935 & 1,374,420,353.085 \\
warehouse-warehouse (DD) & 30,146 & 166,261,943.882 \\
warehouse-destination (DC) & 94,444 & 1,371,554,638.383 \\
origin-destination (OC) & 0 & 0.000 \\
\end{longtable}
}

These are network-arc volumes. Adding different movement families would
count a consignment at successive legs; their sum is not delivered
demand. The full-resolution companion retains 193,525 flow records,
while the expected-horizon sparse matrix contains 4,115 positive
product/arc coordinates.

The following five largest product/arc entries illustrate the paired
flow-distance matrices. They are a selected view, not the complete
matrix.

{\def\LTcaptype{none} % do not increment counter
\begin{longtable}[]{@{}
  >{\raggedright\arraybackslash}p{(\linewidth - 10\tabcolsep) * \real{0.1667}}
  >{\raggedright\arraybackslash}p{(\linewidth - 10\tabcolsep) * \real{0.1667}}
  >{\raggedright\arraybackslash}p{(\linewidth - 10\tabcolsep) * \real{0.1667}}
  >{\raggedright\arraybackslash}p{(\linewidth - 10\tabcolsep) * \real{0.1667}}
  >{\raggedleft\arraybackslash}p{(\linewidth - 10\tabcolsep) * \real{0.1667}}
  >{\raggedleft\arraybackslash}p{(\linewidth - 10\tabcolsep) * \real{0.1667}}@{}}
\toprule\noalign{}
\begin{minipage}[b]{\linewidth}\raggedright
Family
\end{minipage} & \begin{minipage}[b]{\linewidth}\raggedright
From
\end{minipage} & \begin{minipage}[b]{\linewidth}\raggedright
To / role
\end{minipage} & \begin{minipage}[b]{\linewidth}\raggedright
Product
\end{minipage} & \begin{minipage}[b]{\linewidth}\raggedleft
Expected horizon t
\end{minipage} & \begin{minipage}[b]{\linewidth}\raggedleft
OSRM km
\end{minipage} \\
\midrule\noalign{}
\endhead
\bottomrule\noalign{}
\endlastfoot
OD & Toledo - PR & 66.0333.0108-8 & Milho & 44,337,337.208 & 279.7462 \\
OD & Sorriso - MT & 56.F194.0007-8 & Soja & 37,809,356.051 & 30.9958 \\
OD & Sorriso - MT & 56.F194.0007-8 & Milho & 36,693,313.543 & 30.9958 \\
DC & 56.F194.0007-8 & Itaituba - PA / export & Soja & 36,376,022.718 &
1,084.3545 \\
DC & 56.F194.0007-8 & Itaituba - PA / export & Milho & 36,352,297.893 &
1,084.3545 \\
\end{longtable}
}

\begin{figure}[H]

{\centering \pandocbounded{\includegraphics[keepaspectratio]{pyomo_gurobi_h215_w_native8h_license_v2_capacity.pdf}}

}

\caption{Figure E.16. Added static capacity for h215 STO/W, Gurobi,
Pyomo/APPSI, 4 threads.}

\end{figure}%

Circle area represents added candidate plus expansion storage in tonnes.
Bulk-handling basis is excluded from the circles; it determines handling
gains rather than additional static storage.

\subsection{E11 h215 STO/W, HiGHS, Pyomo/APPSI, 4
threads}\label{e11-h215-stow-highs-pyomoappsi-4-threads}

\begin{figure}[H]

{\centering \pandocbounded{\includegraphics[keepaspectratio]{scalar_highs_h215_w_native8h_startup_v2.pdf}}

}

\caption{Figure E.17. Investments for h215 STO/W, HiGHS, Pyomo/APPSI, 4
threads.}

\end{figure}%

The accepted solution has a relative objective gap of
\textbf{0.000196\%}. It opens 69 candidates, expands 73 existing sites
and converts 73 sites for bulk handling.

The
\href{https://github.com/unb-lamfo-or-ai-hpc/model-agrologistic/blob/0788de130253dddc790b179c8831d8ee29dabcbd/manuscript/supplement/investments/scalar_highs_h215_w_native8h_startup_v2.csv}{complete
warehouse investment list} provides CDA, municipality, UF, role,
activation flags, candidate capacity, expansion and bulk-handling basis
in tonnes. Existing-site availability is excluded from candidate
openings.

{\def\LTcaptype{none} % do not increment counter
\begin{longtable}[]{@{}
  >{\raggedright\arraybackslash}p{(\linewidth - 4\tabcolsep) * \real{0.3333}}
  >{\raggedleft\arraybackslash}p{(\linewidth - 4\tabcolsep) * \real{0.3333}}
  >{\raggedleft\arraybackslash}p{(\linewidth - 4\tabcolsep) * \real{0.3333}}@{}}
\toprule\noalign{}
\begin{minipage}[b]{\linewidth}\raggedright
Movement
\end{minipage} & \begin{minipage}[b]{\linewidth}\raggedleft
Positive exported records
\end{minipage} & \begin{minipage}[b]{\linewidth}\raggedleft
Expected arc-tonnes over the horizon
\end{minipage} \\
\midrule\noalign{}
\endhead
\bottomrule\noalign{}
\endlastfoot
origin-warehouse (OD) & 68,279 & 1,374,420,353.085 \\
warehouse-warehouse (DD) & 17,597 & 77,160,138.174 \\
warehouse-destination (DC) & 96,397 & 1,371,586,837.177 \\
origin-destination (OC) & 0 & 0.000 \\
\end{longtable}
}

These are network-arc volumes. Adding different movement families would
count a consignment at successive legs; their sum is not delivered
demand. The full-resolution companion retains 182,273 flow records,
while the expected-horizon sparse matrix contains 3,880 positive
product/arc coordinates.

The following five largest product/arc entries illustrate the paired
flow-distance matrices. They are a selected view, not the complete
matrix.

{\def\LTcaptype{none} % do not increment counter
\begin{longtable}[]{@{}
  >{\raggedright\arraybackslash}p{(\linewidth - 10\tabcolsep) * \real{0.1667}}
  >{\raggedright\arraybackslash}p{(\linewidth - 10\tabcolsep) * \real{0.1667}}
  >{\raggedright\arraybackslash}p{(\linewidth - 10\tabcolsep) * \real{0.1667}}
  >{\raggedright\arraybackslash}p{(\linewidth - 10\tabcolsep) * \real{0.1667}}
  >{\raggedleft\arraybackslash}p{(\linewidth - 10\tabcolsep) * \real{0.1667}}
  >{\raggedleft\arraybackslash}p{(\linewidth - 10\tabcolsep) * \real{0.1667}}@{}}
\toprule\noalign{}
\begin{minipage}[b]{\linewidth}\raggedright
Family
\end{minipage} & \begin{minipage}[b]{\linewidth}\raggedright
From
\end{minipage} & \begin{minipage}[b]{\linewidth}\raggedright
To / role
\end{minipage} & \begin{minipage}[b]{\linewidth}\raggedright
Product
\end{minipage} & \begin{minipage}[b]{\linewidth}\raggedleft
Expected horizon t
\end{minipage} & \begin{minipage}[b]{\linewidth}\raggedleft
OSRM km
\end{minipage} \\
\midrule\noalign{}
\endhead
\bottomrule\noalign{}
\endlastfoot
OD & Toledo - PR & 66.0333.0108-8 & Milho & 44,248,444.895 & 279.7462 \\
OD & Sorriso - MT & 56.F194.0007-8 & Soja & 37,644,913.612 & 30.9958 \\
OD & Sorriso - MT & 56.F194.0007-8 & Milho & 36,819,569.223 & 30.9958 \\
DC & 56.F194.0007-8 & Itaituba - PA / export & Milho & 36,425,453.697 &
1,084.3545 \\
DC & 56.F194.0007-8 & Itaituba - PA / export & Soja & 36,210,245.083 &
1,084.3545 \\
\end{longtable}
}

\begin{figure}[H]

{\centering \pandocbounded{\includegraphics[keepaspectratio]{scalar_highs_h215_w_native8h_startup_v2_capacity.pdf}}

}

\caption{Figure E.18. Added static capacity for h215 STO/W, HiGHS,
Pyomo/APPSI, 4 threads.}

\end{figure}%

Circle area represents added candidate plus expansion storage in tonnes.
Bulk-handling basis is excluded from the circles; it determines handling
gains rather than additional static storage.

\subsection{E12 h300 STO/W, Gurobi, Pyomo/APPSI, one
thread}\label{e12-h300-stow-gurobi-pyomoappsi-one-thread}

\begin{figure}[H]

{\centering \pandocbounded{\includegraphics[keepaspectratio]{h300_w_sto_gurobi_t1_r1.pdf}}

}

\caption{Figure E.19. Investments in the common-Pyomo h300 Gurobi
solution.}

\end{figure}%

The original accepted plan has 154 candidates with positive capacity, 70
positive expansions and 76 positive bulk-handling investments. The
\href{https://github.com/unb-lamfo-or-ai-hpc/model-agrologistic/blob/083ef96cc99a346204df22551ce0739c10e8a2f7/manuscript/supplement/investments/h300_w_sto_gurobi_t1_r1.csv}{complete
investment list} retains facility IDs, flags and quantities in tonnes.

\begin{figure}[H]

{\centering \pandocbounded{\includegraphics[keepaspectratio]{h300_w_sto_gurobi_t1_r1_capacity.pdf}}

}

\caption{Figure E.20. Added static capacity in the common-Pyomo h300
Gurobi solution.}

\end{figure}%

Circle area represents added candidate plus expansion static tonnes;
bulk handling is a separate handling gain.

The source-bound companion contains 225,891 positive flow records and
5,079 expected-horizon OD coordinates. Each flow is joined to its exact
directed distance in the frozen workbook. The
\href{https://github.com/unb-lamfo-or-ai-hpc/model-agrologistic/blob/083ef96cc99a346204df22551ce0739c10e8a2f7/docs/evidence/paired_h300_flow_companions_20261101.json}{companion
inventory} records hashes and aggregation scopes. Full CSV.gz and SQLite
files are prepared locally for a subsequent reviewed data release; no
new public download or exact-zero optimum is implied.

\subsection{E13 h300 STO/W, HiGHS, Pyomo/APPSI, one
thread}\label{e13-h300-stow-highs-pyomoappsi-one-thread}

\begin{figure}[H]

{\centering \pandocbounded{\includegraphics[keepaspectratio]{h300_w_sto_highs_t1_r1.pdf}}

}

\caption{Figure E.21. Investments in the common-Pyomo h300 HiGHS
solution.}

\end{figure}%

The original accepted plan has 154 candidates with positive capacity, 69
positive expansions and 77 positive bulk-handling investments. The
\href{https://github.com/unb-lamfo-or-ai-hpc/model-agrologistic/blob/083ef96cc99a346204df22551ce0739c10e8a2f7/manuscript/supplement/investments/h300_w_sto_highs_t1_r1.csv}{complete
investment list} retains facility IDs, flags and quantities in tonnes.

\begin{figure}[H]

{\centering \pandocbounded{\includegraphics[keepaspectratio]{h300_w_sto_highs_t1_r1_capacity.pdf}}

}

\caption{Figure E.22. Added static capacity in the common-Pyomo h300
HiGHS solution.}

\end{figure}%

Circle area represents added candidate plus expansion static tonnes;
bulk handling is a separate handling gain.

The source-bound companion contains 215,100 positive flow records and
4,902 expected-horizon OD coordinates. Each flow is joined to its exact
directed distance in the frozen workbook. The
\href{https://github.com/unb-lamfo-or-ai-hpc/model-agrologistic/blob/083ef96cc99a346204df22551ce0739c10e8a2f7/docs/evidence/paired_h300_flow_companions_20261101.json}{companion
inventory} records hashes and aggregation scopes. Full CSV.gz and SQLite
files are prepared locally for a subsequent reviewed data release; no
new public download or exact-zero optimum is implied.

Facilities at the same recorded coordinates can overlap on the national
maps. Consequently, visually identical marks do not establish identical
investment vectors. For example, facility 47.0287.0007-0 receives
expansion in the Gurobi plan and bulk handling in the HiGHS plan, while
another facility at the same coordinates also has bulk handling. The
complete ID-based lists, rather than symbol counts, determine the
reported investment differences.



\section{Appendix F. Research data, provenance and
access}\label{appendix-f.-research-data-provenance-and-access}

\subsection{F1. Frozen input data}\label{f1.-frozen-input-data}

The \href{https://doi.org/10.5281/zenodo.22751909}{versioned
Agrologistic v0.2.0 record} contains the historical input archive.
Appendix D identifies the workbooks, their SHA-256 digests, the
warehouse populations and the directed OSRM distances used for each
model size. Cite the version DOI and match the workbook digest to the
experimental record before comparing results. The terms applicable to
the original registry and OpenStreetMap-derived road data are documented
separately from the software licence.

\subsection{F2. National investment
summaries}\label{f2.-national-investment-summaries}

The proposed companion dataset contains national summaries of candidate
opening, existing-site expansion and bulk-handling conversion. It
contains 10 case--decision rows across seven Gurobi configurations. The
data dictionary distinguishes the binary activation count from modeled
capacity in metric tonnes. Each value can be recomputed from the
corresponding GitHub investment list. The package documents coverage,
source identities, field definitions, licence scope and checksums.

The companion dataset has been prepared and verified locally and is
awaiting deposit. The v0.2.0 DOI identifies the historical inputs. The
GitHub investment lists provide the source values for the national
summaries.

\subsection{F3. Full-resolution
outputs}\label{f3.-full-resolution-outputs}

The full positive-flow tables, endpoint-level distance exports, selected
graphs and SQLite companions are outside the proposed deposit. The
national investment summaries support verification of those totals.
Exact numerical reproduction uses the input, network, model and
execution records identified in Appendices C--E. The release
documentation records access and rights by asset. Full-resolution
transfers depend on the applicable third-party terms and disclosure
review.

\end{document}
